\documentclass[12pt,letterpaper,oneside]{article}

%------------------------- Idioma y codificacion ---------------------
\usepackage[utf8]{inputenc}
\usepackage[T1]{fontenc}
\usepackage[spanish,es-noquoting,es-nodecimaldot]{babel}
\addto\captionsspanish{\renewcommand{\tablename}{Tabla}}
\addto\captionsspanish{\renewcommand{\listtablename}{ÍNDICE DE TABLAS}}
\AtBeginDocument{\shorthandoff{"}}

%------------------------- Tipografia --------------------------------
\IfFileExists{newtxtext.sty}
  {\usepackage{newtxtext}\usepackage{newtxmath}}
  {\usepackage{mathptmx}}

%------------------------- Margenes y pagina -------------------------
\usepackage[
  letterpaper,
  top=2.6cm, bottom=2.5cm, left=3.0cm, right=3.0cm,
  headheight=52pt, headsep=14pt, footskip=32pt
]{geometry}

\usepackage{graphicx}
\usepackage{longtable}
\usepackage{array}
\usepackage{tabularx}
\usepackage{caption}
\usepackage{xcolor}
\usepackage{setspace}
\usepackage{fancyhdr}
\usepackage{titlesec}
\usepackage{ragged2e}
\usepackage{float}
\usepackage[hidelinks,plainpages=false]{hyperref}
\hypersetup{pdftitle={TD: Gamma Lab 2.0}, pdfauthor={Equipo Gamma Lab}}

\onehalfspacing
\setlength{\parindent}{1.0cm}
\setcounter{secnumdepth}{3}
\setcounter{tocdepth}{2}
\renewcommand{\contentsname}{TABLA DE CONTENIDO}
\setlength{\emergencystretch}{3em}
\DeclareRobustCommand{\_}{\textunderscore\allowbreak}

%=====================================================================
%  DATOS DEL DOCUMENTO
%=====================================================================
\newcommand{\tituloProyecto}{GAMMA LAB 2.0}
\newcommand{\tituloCorto}{GAMMA LAB 2.0}
\newcommand{\autorUno}{Juan Felipe Romero Ortiz}
\newcommand{\autorDos}{Laura Juliana Garzón Arias}
\newcommand{\autorTres}{Sebastián Almanza Galvis}
\newcommand{\autorCuatro}{Xiomara Herrera Suárez}
\newcommand{\nombreDocumento}{\textit{TEST DOCUMENT} (TD)}
\newcommand{\director}{Ing. Cécile Gauthier Umaña}
\newcommand{\juradoUno}{Ing. Nombre del Primer Jurado}
\newcommand{\juradoDos}{Ing. Nombre del Segundo Jurado}
\newcommand{\universidad}{PONTIFICIA UNIVERSIDAD JAVERIANA}
\newcommand{\facultad}{FACULTAD DE INGENIERÍA}
\newcommand{\programa}{PROGRAMA DE INGENIERÍA DE SISTEMAS}
\newcommand{\ciudad}{BOGOTÁ, D.C.}
\newcommand{\fechaLarga}{Noviembre, 2026}
\newcommand{\logoUniversidad}{imagen_logo_universidad}
\newcommand{\logoProyecto}{imagen_logo_gammalab}
%=====================================================================

\newcommand{\imgSegura}[3]{%
  \IfFileExists{#1.png}{\includegraphics[height=#2]{#1.png}}{%
  \IfFileExists{#1.pdf}{\includegraphics[height=#2]{#1.pdf}}{%
  \IfFileExists{#1.jpg}{\includegraphics[height=#2]{#1.jpg}}{%
  \textcolor{gray!60}{\framebox[#3]{\rule{0pt}{#2}\scriptsize logo}}}}}}

% Texto de relleno para lo que se completara con el usuario.
\newcommand{\pendiente}[1][Escriba aquí el contenido de esta sección.]{%
  {\color{gray!75}\itshape [#1]}\par}
% Casilla de seleccion para los cuadros que se diligencian a mano.
\newcommand{\casilla}{\framebox[0.75em]{\rule{0pt}{0.45em}}}

\newcommand{\capitulo}[1]{\clearpage\section{#1}}

%------------------- Encabezado y pie de pagina ----------------------
\pagestyle{fancy}
\fancyhf{}
\fancyhead[L]{%
  \raisebox{-0.35\height}{\imgSegura{\logoUniversidad}{1.05cm}{2.2cm}}%
  \hspace{0.5cm}%
  \raisebox{-0.35\height}{\imgSegura{\logoProyecto}{0.85cm}{1.6cm}}%
}
\fancyhead[R]{\textbf{TD: \tituloCorto}}
\fancyfoot[R]{\textbf{Página \thepage}}
\renewcommand{\headrulewidth}{1.1pt}
\renewcommand{\footrulewidth}{1.1pt}

%------------------- Estilo de titulos -------------------------------
\titleformat{\section}{\normalfont\Large\bfseries}{\thesection}{0.8em}{}
\titlespacing*{\section}{0pt}{0pt}{16pt}
\titleformat{\subsection}{\normalfont\large\bfseries}{\thesubsection}{0.7em}{}
\titlespacing*{\subsection}{0pt}{14pt}{8pt}
\titleformat{\subsubsection}{\normalfont\normalsize\bfseries}{\thesubsubsection}{0.6em}{}
\titlespacing*{\subsubsection}{0pt}{12pt}{6pt}

%=====================================================================
\begin{document}
\pagenumbering{roman}

%------------------------ PORTADA -------------------------------------
\begin{center}
  {\Large\bfseries \tituloProyecto}\par
  \vspace{1.4cm}

  {\bfseries Autores:}\par
  \vspace{0.5cm}
  \autorUno\par
  \autorDos\par
  \autorTres\par
  \autorCuatro\par
  \vspace{1.4cm}

  \nombreDocumento\par
  \vspace{1.4cm}

  {\bfseries Director}\par
  \vspace{0.4cm}
  \director\par
  \vspace{1.2cm}

  {\bfseries Jurados del Trabajo Final de Grado}\par
  \vspace{0.4cm}
  \juradoUno\par
  \juradoDos\par
  \vspace{1.8cm}

  \universidad\par\vspace{0.6cm}
  \facultad\par\vspace{0.6cm}
  \programa\par\vspace{0.6cm}
  \ciudad\par\vspace{0.6cm}
  \fechaLarga\par
\end{center}
\clearpage

%------------------------ TABLA DE CONTENIDO -------------------------
\tableofcontents
\clearpage
\listoftables

\clearpage
\pagenumbering{arabic}

\section{INTRODUCCIÓN}

El presente \textit{Test Document} (TD) describe el plan de pruebas de GammaLab 2.0 y reúne los resultados obtenidos al ejecutarlo. GammaLab 2.0 conserva la arquitectura de \textit{microkernel} de la versión anterior: un núcleo de servicios esenciales al que se conectan extensiones, o \textit{plugins}, que aportan cada técnica de análisis. Sobre esa base, la segunda versión corrige las limitaciones identificadas en la primera, incorpora el análisis de acoplamiento fase-amplitud (PAC), el índice de modulación y la consistencia de fase entre canales (PPC), introduce la gestión de proyectos y traslada el cómputo pesado a un servicio de tareas del núcleo.

El plan de pruebas se organiza en seis frentes: las pruebas unitarias y de integración de los servicios del núcleo; las pruebas de los módulos de análisis; la validación de los resultados frente al sistema de referencia BOARD-FTD-PACC; la verificación de los requisitos no funcionales mediante pruebas y mediciones de rendimiento; las pruebas de usabilidad, y las pruebas de aceptación del usuario. Los cuatro primeros frentes son automáticos o medidos por el equipo de desarrollo, y sus resultados se presentan en este documento. Los dos últimos requieren la participación del usuario del Laboratorio de Electrofisiología y, a la fecha de esta versión del documento, se encuentran pendientes de ejecución; por ello, sus casos y cuadros de registro se presentan preparados para diligenciarse durante las sesiones con el usuario.

Este documento toma como punto de partida el TD de GammaLab 1.0 \cite{td2025}, del que conserva la estructura y el formato de los cuadros de aceptación, y se apoya en la especificación de requisitos (SRS), en el documento de arquitectura (SAD) y en el documento de diseño (SDD) de la presente versión.

\capitulo{ALCANCE}

\subsection{Enfoque general}
Las pruebas se centran en la funcionalidad crítica del sistema y en la experiencia del usuario final. Se busca garantizar que los servicios del núcleo y las extensiones funcionen de forma correcta dentro del entorno definido, que los resultados numéricos sean coherentes con los del sistema de referencia BOARD-FTD-PACC y que se cumplan los límites de tiempo, memoria y estabilidad fijados en el SRS. Las pruebas se desarrollan sobre señales suministradas por el laboratorio y sobre señales sintéticas cuya respuesta correcta se conoce de antemano.

\subsection{Elementos en alcance}
\begin{itemize}
    \item \textbf{Carga de señales y gestión de proyectos:} lectura de archivos ABF, EDF y MAT, y guardado, apertura y cierre de proyectos \textit{.glab}.
    \item \textbf{Preprocesamiento:} segmentación en \textit{trials}, selección y descarte, remoción de artefactos y filtrado, sin modificar los datos originales.
    \item \textbf{Análisis:} dominios del tiempo, de la frecuencia y de tiempo-frecuencia, PAC de uno y de dos canales, índice de modulación y PPC.
    \item \textbf{Medidas y exportación:} pendiente y amplitud sobre las gráficas, y exportación de imágenes y tablas.
    \item \textbf{Servicios del núcleo:} almacén de datos, servicio de tareas, servicios compartidos y liberación de recursos.
    \item \textbf{Requisitos no funcionales:} usabilidad, rendimiento, confiabilidad y soporte.
    \item \textbf{Validación con el usuario:} usabilidad de la interfaz y aceptación de cada requisito funcional.
\end{itemize}

\subsection{Fuera de alcance}
\begin{itemize}
    \item Integraciones en la nube o con servicios externos no definidos en el SRS.
    \item Compatibilidad con sistemas operativos distintos de Windows 10 y 11 de 64 bits.
    \item Formatos de señal distintos de ABF, EDF y MAT.
    \item Pruebas estadísticas de significancia del acoplamiento, que el sistema no implementa y quedan a cargo del investigador.
    \item Extensiones desarrolladas fuera del alcance del proyecto.
\end{itemize}

\subsection{Supuestos y dependencias}
\begin{itemize}
    \item Disponibilidad de la señal \textit{17308005.abf}, suministrada por el laboratorio, y de las referencias numéricas exportadas desde BOARD-FTD-PACC, que se conservan en la carpeta \textit{test/data} del repositorio.
    \item Disponibilidad del usuario del Laboratorio de Electrofisiología para las sesiones de usabilidad y de aceptación.
    \item Instalación de GammaLab 2.0 mediante su instalador en el equipo en que se realicen las pruebas con el usuario.
\end{itemize}

\subsection{Herramientas}
\begin{itemize}
    \item \textbf{\textit{pytest} 8.4.2:} desarrollo y ejecución de las pruebas automáticas.
    \item \textbf{BOARD-FTD-PACC y MATLAB:} generación de las referencias numéricas con las que se comparan los análisis.
    \item \textbf{\textit{time.perf\_counter} de Python:} medición de los tiempos de respuesta de los requisitos de rendimiento.
    \item \textbf{\textit{Visual Studio Code}:} entorno de desarrollo.
\end{itemize}

\subsection{Entorno de ejecución de las pruebas automáticas}
Las pruebas automáticas y las mediciones se ejecutaron en el equipo descrito en la Tabla~\ref{tab:td_entorno}, con el comando \textit{python -m pytest test} desde la carpeta raíz del repositorio.

\begin{center}
\small
\begin{tabularx}{\textwidth}{|p{0.30\textwidth}|X|}
\hline
\textbf{Elemento} & \textbf{Valor} \\
\hline
Procesador & Intel Core i5-11400H a 2,70\,GHz \\
\hline
Memoria RAM & 32\,GB \\
\hline
Sistema operativo & Windows 11, 64 bits \\
\hline
Intérprete & Python 3.12.1 \\
\hline
Fecha de la ejecución & 8 de noviembre de 2026 \\
\hline
\end{tabularx}
\captionof{table}{Entorno de ejecución de las pruebas automáticas}
\label{tab:td_entorno}
\end{center}
\capitulo{ORGANIZACIÓN DE LAS PRUEBAS Y RESUMEN DE RESULTADOS}
\label{sec:td_resumen}

La suite automática reside en la carpeta \textit{test} del repositorio y se divide en tres subcarpetas: \textit{services\_test}, con las pruebas de los servicios del núcleo; \textit{plugins\_test}, con las pruebas de los módulos de análisis y las comparaciones con el sistema de referencia, y \textit{performance\_test}, con las pruebas de rendimiento. La Tabla~\ref{tab:td_resumen} resume el resultado de la ejecución del 8 de noviembre de 2026.

\begin{center}
\small
\begin{tabularx}{\textwidth}{|X|c|c|c|c|c|}
\hline
\textbf{Frente de prueba} & \textbf{Suites} & \textbf{Pruebas} & \textbf{Aprobadas} & \textbf{Fallo esperado} & \textbf{Fallidas} \\
\hline
Unitarias y de integración del núcleo (Sección~\ref{sec:td_unitarias}) & 34 & 292 & 291 & 0 & 1 \\
\hline
Módulos de análisis (Sección~\ref{sec:td_modulos}) & 16 & 417 & 417 & 0 & 0 \\
\hline
Validación frente a BOARD-FTD-PACC (Sección~\ref{sec:td_validacion}) & 9 & 27 & 22 & 5 & 0 \\
\hline
Rendimiento (Sección~\ref{sec:td_rendimiento}) & 6 & 25 & 25 & 0 & 0 \\
\hline
\textbf{Total} & \textbf{65} & \textbf{761} & \textbf{755} & \textbf{5} & \textbf{1} \\
\hline
\end{tabularx}
\captionof{table}{Resumen de la ejecución de la suite automática}
\label{tab:td_resumen}
\end{center}

Una prueba con \textbf{fallo esperado} es una prueba que se mantiene en la suite para documentar una diferencia conocida y explicada: se ejecuta en cada corrida y, si algún día dejara de fallar, la suite lo advertiría. Las 5 pruebas en esta condición corresponden a comparaciones con las referencias de MATLAB cuya causa se analiza en la Sección~\ref{sec:td_diferencias}.

La única prueba fallida, \textit{test\_el\_orden\_de\_botones\_viene\_de\_los\_metadatos}, de la suite \textit{test\_plugin\_sin\_tocar\_nucleo}, no revela un defecto del sistema sino una prueba desactualizada: comprueba el orden de los botones de la subcategoría tiempo-frecuencia esperando cuatro extensiones (PAC, índice de modulación, \textit{Wavelet} y \textit{Wavelet Average}), y desde la incorporación de PPC la subcategoría tiene cinco. La prueba debe actualizarse para incluir a PPC en la lista esperada.

Las pruebas de usabilidad (Sección~\ref{sec:td_usabilidad}) y de aceptación del usuario (Sección~\ref{sec:td_uat}) son manuales y se encuentran pendientes de ejecución.
\capitulo{PRUEBAS UNITARIAS Y DE INTEGRACIÓN}
\label{sec:td_unitarias}

\subsection{Objetivo}
Estas pruebas evalúan el funcionamiento interno de los servicios que el núcleo ofrece a las extensiones: la lectura de señales, el almacén de datos compartido, el servicio de tareas, la medición, la exportación, la gestión de proyectos y la liberación de recursos. Su propósito es asegurar que cada servicio cumpla su contrato de forma aislada y que los servicios colaboren correctamente entre sí.

\subsection{Implementación}
Las pruebas se implementaron con \textit{pytest}. Las unitarias ejercitan cada servicio con datos controlados y verifican estructuras, dimensiones, valores y excepciones; las de integración combinan varios servicios con la ventana principal o con el servicio de tareas en funcionamiento, para comprobar, por ejemplo, que un cálculo pesado no bloquea la interfaz o que un proyecto guardado se restaura completo. Cuando una prueba requiere la interfaz gráfica, se ejecuta sin mostrar ventanas en pantalla.

\subsection{Resultados}
La Tabla~\ref{tab:td_servicios} presenta cada suite, lo que verifica y su resultado. Entre paréntesis se indica el número de la matriz de requisitos que la suite verifica, cuando corresponde.

\begin{center}
\footnotesize
\begin{longtable}{|>{\raggedright\arraybackslash}p{0.22\textwidth}|>{\raggedright\arraybackslash}p{0.33\textwidth}|>{\centering\arraybackslash}p{0.09\textwidth}|>{\centering\arraybackslash}p{0.11\textwidth}|>{\raggedright\arraybackslash}p{0.10\textwidth}|}
\hline
\textbf{Suite} & \textbf{Qué verifica} & \textbf{Pruebas} & \textbf{Aprobadas} & \textbf{Resultado} \\
\hline
\endfirsthead
\hline
\textbf{Suite} & \textbf{Qué verifica} & \textbf{Pruebas} & \textbf{Aprobadas} & \textbf{Resultado} \\
\hline
\endhead
\hline
\endfoot
\hline
\caption{Resultados de las pruebas unitarias y de integración del núcleo} \label{tab:td_servicios} \\
\endlastfoot
\textit{test\_apertura\_proyecto} & La apertura de un proyecto lee el archivo fuera del hilo de interfaz y el guardado no vuelve a copiar la señal (R04). & 5 & 5 & Aprobada \\
\hline
\textit{test\_artifact\_remove\_orquestador} & \textit{Remove Artifact} calcula a través del servicio de tareas sin modificar los datos compartidos ni crear hilos propios. & 6 & 6 & Aprobada \\
\hline
\textit{test\_cache\_tabla\_vtk} & La tabla VTK de una señal se reutiliza y se descarta cuando la señal cambia (R72). & 3 & 3 & Aprobada \\
\hline
\textit{test\_cancelacion\_al\_cambiar\_seccion} & Al cambiar de sección se cancelan las tareas del módulo que se abandona, sin afectar a los demás. & 2 & 2 & Aprobada \\
\hline
\textit{test\_cu018\_sampling\_rate} & Corrección manual de la frecuencia de muestreo, validación del valor y restauración del original (CU-018). & 4 & 4 & Aprobada \\
\hline
\textit{test\_data\_store\_concurrencia} & El \textit{DataStore} admite un único escritor y lecturas seguras entre hilos (R54). & 10 & 10 & Aprobada \\
\hline
\textit{test\_dataset} & Operaciones básicas del almacén: claves, señales, señal activa y manejo de errores. & 7 & 7 & Aprobada \\
\hline
\textit{test\_emf\_export} & Exportación vectorial EMF y exportación de tablas de mapas de calor (R35). & 9 & 9 & Aprobada \\
\hline
\textit{test\_export} & Exportación de imágenes y tablas (PNG, CSV, JSON), cancelaciones y formatos no soportados. & 11 & 11 & Aprobada \\
\hline
\textit{test\_fileio\_abf} & Lectura de un archivo ABF válido y manejo de un archivo corrupto. & 2 & 2 & Aprobada \\
\hline
\textit{test\_fileio\_edf} & Lectura de un archivo EDF válido y manejo de un archivo corrupto. & 2 & 2 & Aprobada \\
\hline
\textit{test\_fileio\_sin\_duplicar} & Los lectores de archivos no duplican los datos de la señal en memoria. & 10 & 10 & Aprobada \\
\hline
\textit{test\_gabor\_en\_core} & La transformada de Gabor reside en el núcleo de algoritmos compartidos y no depende de las extensiones. & 6 & 6 & Aprobada \\
\hline
\textit{test\_guardado\_proyecto} & Guardado del proyecto en pasos, con la copia de la señal fuera del hilo de interfaz (R04, R07). & 8 & 8 & Aprobada \\
\hline
\textit{test\_indicadores\_ui} & Severidad de los mensajes, indicador del \textit{trial} actual, \textit{trials} usados y lectura X/Y bajo el cursor (R02, R26, R37, R38). & 24 & 24 & Aprobada \\
\hline
\textit{test\_instalacion\_rutas} & El programa instalado guarda sus ajustes en la carpeta del usuario y abre el archivo \textit{.glab} que recibe (R39). & 8 & 8 & Aprobada \\
\hline
\textit{test\_liberar\_senal} & Al cerrar el proyecto o cargar otra señal se libera la anterior junto con sus derivados (R11, R52). & 5 & 5 & Aprobada \\
\hline
\textit{test\_liberar\_ventanas} & Al salir de un módulo pesado se libera su ventana OpenGL sin afectar a los demás (R45). & 6 & 6 & Aprobada \\
\hline
\textit{test\_measurement} & Servicio de medición: estados, identificadores, contexto, cálculos y limpieza. & 28 & 28 & Aprobada \\
\hline
\textit{test\_medidas\_artefacto} & Las medidas del segmento se calculan y se piden confirmar antes de modificarlo (R40). & 6 & 6 & Aprobada \\
\hline
\textit{test\_niveles\_de\_mensaje} & Cada situación usa el nivel de severidad que le corresponde (R02). & 5 & 5 & Aprobada \\
\hline
\textit{test\_persistencia\_artifact\_remove} & Las modificaciones de \textit{Remove Artifact} se conservan al guardar y abrir el proyecto. & 15 & 15 & Aprobada \\
\hline
\textit{test\_plugin\_sin\_tocar\_nucleo} & Una extensión nueva se guarda con el proyecto y se ordena desde su \textit{properties.yml} sin modificar el núcleo (R77). & 4 & 3 & 1 fallida \\
\hline
\textit{test\_resultados\_xlsx} & Los paneles de resultados de pendiente y amplitud exportan a Excel. & 4 & 4 & Aprobada \\
\hline
\textit{test\_run\_task} & Patrón \textit{IPlugin.run\_task} para despachar cómputos al servicio de tareas (R70, R71). & 6 & 6 & Aprobada \\
\hline
\textit{test\_servicios\_compartidos} & Los servicios de ajustes, exportación y medición son únicos y los registra el núcleo (R69). & 7 & 7 & Aprobada \\
\hline
\textit{test\_submuestreo} & Se dibujan como máximo 2000 puntos por serie, sin alterar los datos (R73). & 9 & 9 & Aprobada \\
\hline
\textit{test\_task\_service} & Servicio de tareas: cola, una tarea a la vez, progreso, cancelación cooperativa y manejo de fallos (R46). & 14 & 14 & Aprobada \\
\hline
\textit{test\_task\_service\_integracion} & Una extensión calcula la transformada \textit{wavelet} por el servicio de tareas mientras la interfaz sigue respondiendo. & 2 & 2 & Aprobada \\
\hline
\textit{test\_trials} & Detección de estímulos, ventanas de corte en sus dos modos y ensamblaje de los \textit{trials}. & 23 & 23 & Aprobada \\
\hline
\textit{test\_trials\_derivados} & El preprocesamiento crea un conjunto derivado sin modificar los \textit{trials} originales (R79). & 16 & 16 & Aprobada \\
\hline
\textit{test\_trials\_por\_canal} & Consulta de los \textit{trials} de un canal concreto, y no solo del último generado (R68). & 12 & 12 & Aprobada \\
\hline
\textit{test\_vtk\_observers} & Los observadores de VTK se retiran al salir del módulo y se reponen al volver (R58). & 6 & 6 & Aprobada \\
\hline
\textit{test\_xlsx\_writer} & El escritor propio de archivos \textit{.xlsx} produce libros válidos y rápidos (R51). & 7 & 7 & Aprobada \\
\hline
\end{longtable}
\end{center}

\subsection{Análisis de resultados}
Las 292 pruebas de este frente se aprobaron, con la única excepción de la prueba desactualizada descrita en la Sección~\ref{sec:td_resumen}. Destacan tres conjuntos. Las pruebas del servicio de tareas y de su integración con la interfaz confirman que el cómputo pesado se ejecuta fuera del hilo de interfaz, que solo corre una tarea a la vez y que la cancelación es cooperativa. Las pruebas de liberación de la señal y de las ventanas gráficas verifican que la memoria se recupera al cerrar un proyecto, al cargar otra señal y al salir de un módulo. Por último, las pruebas de proyectos comprueban que el guardado y la apertura conservan la señal, los \textit{trials}, los descartes, el preprocesamiento, las medidas y los parámetros de los análisis.
\capitulo{PRUEBAS DE LOS MÓDULOS DE ANÁLISIS}
\label{sec:td_modulos}

\subsection{Objetivo}
Estas pruebas verifican el cálculo de cada módulo de análisis y su integración en la extensión correspondiente, con independencia de la comparación frente al sistema de referencia, que se trata en la Sección~\ref{sec:td_validacion}.

\subsection{Implementación}
Cada módulo de análisis separa el cálculo, que recibe y devuelve arreglos numéricos, de la lógica de la interfaz. Esa separación permite probar el cálculo sin interfaz gráfica: las pruebas emplean señales sintéticas con respuesta conocida, por ejemplo una oscilación lenta que modula la amplitud de una rápida con una fase preferente fijada, y verifican que el resultado reproduce esa propiedad. Las pruebas de integración comprueban que cada extensión envía su cálculo al servicio de tareas, exporta sus resultados y restaura sus parámetros desde el proyecto.

\subsection{Resultados}

\begin{center}
\footnotesize
\begin{longtable}{|>{\raggedright\arraybackslash}p{0.22\textwidth}|>{\raggedright\arraybackslash}p{0.33\textwidth}|>{\centering\arraybackslash}p{0.09\textwidth}|>{\centering\arraybackslash}p{0.11\textwidth}|>{\raggedright\arraybackslash}p{0.10\textwidth}|}
\hline
\textbf{Suite} & \textbf{Qué verifica} & \textbf{Pruebas} & \textbf{Aprobadas} & \textbf{Resultado} \\
\hline
\endfirsthead
\hline
\textbf{Suite} & \textbf{Qué verifica} & \textbf{Pruebas} & \textbf{Aprobadas} & \textbf{Resultado} \\
\hline
\endhead
\hline
\endfoot
\hline
\caption{Resultados de las pruebas de los módulos de análisis} \label{tab:td_modulos} \\
\endlastfoot
\textit{test\_average\_seleccion} & \textit{Average} promedia solo los \textit{trials} seleccionados por el usuario (R41). & 8 & 8 & Aprobada \\
\hline
\textit{test\_ayudas\_controles} & Todos los controles de cada módulo tienen descripción emergente (R42). & 21 & 21 & Aprobada \\
\hline
\textit{test\_correcciones\_interaccion} & Correcciones de navegación e interacción: proyecto, bienvenida, cierre, restablecimiento de parámetros, lectura del cursor, selección de \textit{trials}, \textit{zoom}, \textit{Butterfly}, submuestreo, filtro y apertura (R07, R10, R12, R25, R26, R27, R30, R32, R73, R76, R78). & 25 & 25 & Aprobada \\
\hline
\textit{test\_erp\_compute} & Cálculo del mapa de calor de ERP, separado del dibujo. & 13 & 13 & Aprobada \\
\hline
\textit{test\_erp\_vtk\_llenado} & El llenado vectorizado del mapa de calor de ERP produce el mismo resultado que el llenado elemento por elemento (R59). & 7 & 7 & Aprobada \\
\hline
\textit{test\_escalograma\_rango\_color} & El rango de valores de color del escalograma puede fijarse manualmente (R33). & 7 & 7 & Aprobada \\
\hline
\textit{test\_modulation\_index\_core} & Comodulograma del índice de modulación, portado de la función \textit{f\_PAC\_sing} de BOARD-FTD-PACC. & 27 & 27 & Aprobada \\
\hline
\textit{test\_pac\_core} & Núcleo de PAC: filtros IIR, transformada de Morse y pasos del cálculo frente al comportamiento de BOARD-FTD-PACC. & 124 & 124 & Aprobada \\
\hline
\textit{test\_pac\_grados} & Eje de fase de PAC expresado en grados entre 0 y 360. & 19 & 19 & Aprobada \\
\hline
\textit{test\_pac\_intercanal} & PAC Psel entre canales: ciclos y fase de un canal, amplitud de otro (R60, R61). & 16 & 16 & Aprobada \\
\hline
\textit{test\_pac\_intercanal\_exportacion} & Exportación numérica del PAC entre canales, con el canal de amplitud entre los parámetros (R65). & 4 & 4 & Aprobada \\
\hline
\textit{test\_pac\_intercanal\_plugin} & Integración del modo PAC Psel inter-canal en la extensión PAC (R62). & 15 & 15 & Aprobada \\
\hline
\textit{test\_ppc\_core} & Núcleo de PPC frente a señales sintéticas de verdad conocida (R63, R64). & 84 & 84 & Aprobada \\
\hline
\textit{test\_ppc\_plugin\_orquestado} & Ejecución de PPC a través del servicio de tareas, con progreso y cancelación. & 23 & 23 & Aprobada \\
\hline
\textit{test\_relative\_psd\_exportacion} & Exportación numérica de la PSD relativa: parámetros, potencias y fracción por banda (R05). & 4 & 4 & Aprobada \\
\hline
\textit{test\_submuestreo\_wavelet} & Submuestreo de la señal antes de la transformada \textit{wavelet} y uso de la tasa efectiva. & 20 & 20 & Aprobada \\
\hline
\end{longtable}
\end{center}

\subsection{Análisis de resultados}
Las 417 pruebas de este frente se aprobaron. El volumen se concentra en los módulos nuevos de la versión 2.0: 124 pruebas para el núcleo de PAC, 84 para el de PPC y 27 para el índice de modulación, que verifican cada paso del cálculo y su respuesta ante entradas que comprometen el resultado, como bandas fuera del rango de Nyquist, canales no alineados o \textit{trials} demasiado cortos.
\capitulo{PRUEBAS DE VALIDACIÓN FRENTE AL SISTEMA DE REFERENCIA}
\label{sec:td_validacion}

\subsection{Objetivo}
La intención de estas pruebas es comparar los resultados de GammaLab 2.0 con los del sistema de referencia BOARD-FTD-PACC, para garantizar la exactitud de los análisis y documentar con precisión las diferencias que existen entre ambas implementaciones.

\subsection{Implementación}
Para estas pruebas se utilizó la señal \textit{17308005.abf}, suministrada por el laboratorio, que cuenta con dos canales. En ambos sistemas se emplearon los mismos parámetros de segmentación:

\begin{itemize}
    \item Canal: CA1.
    \item Canal de estímulo: IN7.
    \item Umbral: 0,7.
    \item Número de estímulos por \textit{trial}: 1.
    \item Ventana: de $-0{,}05$ a $4{,}0$\,s respecto al estímulo.
    \item Modo de cierre: \textit{Cut to the next stim}.
\end{itemize}

Cada análisis trabaja con la totalidad de los \textit{trials}. Para cada prueba se ejecuta el cálculo de la extensión sin su parte gráfica y se compara el resultado con la matriz exportada desde BOARD-FTD-PACC. Los análisis espectrales usan una frecuencia de remuestreo de 1000\,Hz y un rango de 0 a 500\,Hz; PSD y PSD promedio usan ventana de Hamming, segmentos de 1024 muestras, solapamiento de 128 y 1024 puntos de transformada; la transformada \textit{wavelet} usa un rango de 1 a 500\,Hz, dos ciclos y eje de frecuencias lineal, el mismo de la referencia.

Cada análisis se somete a tres criterios: la \textbf{estructura}, que exige las mismas dimensiones que la matriz de referencia; la \textbf{correlación} de Pearson por columna o por fila, que debe superar un umbral, y la comparación \textbf{punto a punto}, que exige que al menos el 95\,\% de los puntos difiera de la referencia en menos de una tolerancia absoluta. La Tabla~\ref{tab:td_validacion} recoge los umbrales y los resultados.

\subsection{Resultados}

\begin{center}
\small
\begin{longtable}{|>{\raggedright\arraybackslash}p{0.15\textwidth}|>{\centering\arraybackslash}p{0.13\textwidth}|>{\centering\arraybackslash}p{0.13\textwidth}|>{\centering\arraybackslash}p{0.14\textwidth}|>{\centering\arraybackslash}p{0.14\textwidth}|>{\centering\arraybackslash}p{0.14\textwidth}|}
\hline
\textbf{Análisis} & \textbf{Correlación mínima} & \textbf{Tolerancia punto a punto} & \textbf{Estructura} & \textbf{Correlación} & \textbf{Punto a punto} \\
\hline
\endfirsthead
\hline
\textbf{Análisis} & \textbf{Correlación mínima} & \textbf{Tolerancia punto a punto} & \textbf{Estructura} & \textbf{Correlación} & \textbf{Punto a punto} \\
\hline
\endhead
\hline
\endfoot
\hline
\caption{Resultados de la validación frente a BOARD-FTD-PACC} \label{tab:td_validacion} \\
\endlastfoot
\textit{Average} & 0,98 & 0,001 & Aprobada & Aprobada & Aprobada \\
\hline
ERP & 0,98 & 0,001 & Aprobada & Aprobada & Aprobada \\
\hline
FFT & 0,98 & 0,75 & Aprobada & Aprobada & Aprobada \\
\hline
\textit{FFT Average} & 0,98 & 0,15 & Aprobada & Aprobada & Aprobada \\
\hline
PSD & 0,88 & 0,0001 & Aprobada & Aprobada & Aprobada \\
\hline
\textit{PSD Average} & 0,95 & 0,0001 & Aprobada & Fallo esperado & Aprobada \\
\hline
\textit{Wavelet} & 0,80 & 0,0001 & Aprobada & Fallo esperado & Fallo esperado \\
\hline
\textit{Wavelet Average} & 0,80 & 0,0001 & Aprobada & Fallo esperado & Fallo esperado \\
\hline
\end{longtable}
\end{center}

\textit{Average}, ERP, FFT, \textit{FFT Average} y PSD superan los tres criterios. \textit{PSD Average} supera la estructura y la comparación punto a punto, pero no el umbral de correlación. \textit{Wavelet} y \textit{Wavelet Average} superan la estructura, pero no la correlación ni la comparación punto a punto. Estas cinco comparaciones se mantienen en la suite como fallos esperados, porque su causa está identificada y medida, como se explica a continuación.

\subsection{Análisis de las diferencias}
\label{sec:td_diferencias}
La suite \textit{test\_referencias\_matlab\_wavelet\_psd}, cuyas 3 pruebas se aprobaron, fija numéricamente las diferencias con la referencia, de modo que cualquier cambio en el cálculo que las agrande sea detectado:

\begin{itemize}
    \item \textbf{Transformada \textit{wavelet}.} BOARD-FTD-PACC entrega la amplitud sin la normalización por la raíz de la escala que aplica la librería PyWavelets, y con un factor de escala de unas 2,05 veces. Al dividir el escalograma de GammaLab entre la raíz de la escala y multiplicarlo por dos, la correlación global con la referencia sube de 0,83 a 0,98 para un \textit{trial} y a 0,997 para el promedio. La diferencia es, por tanto, de convención de amplitud y no de forma del escalograma.
    \item \textbf{PSD promedio.} GammaLab retira la tendencia constante de cada segmento antes de transformarlo, que es el comportamiento por defecto de SciPy, y BOARD-FTD-PACC no lo hace. Sin retirarla, la potencia total coincide con la referencia dentro del 1\,\% y el espectro en escala logarítmica correlaciona 0,999; la diferencia restante se concentra en los primeros intervalos de frecuencia, entre 0 y 3\,Hz.
\end{itemize}

\subsection{Validación de PAC, índice de modulación y PPC}
Los módulos nuevos de la versión 2.0 se validaron de forma distinta. PAC y el índice de modulación se portaron de las funciones de BOARD-FTD-PACC, y sus pruebas (124 para PAC y 27 para el índice de modulación) verifican cada paso del cálculo frente al comportamiento del código de referencia, incluidos los filtros IIR, la transformada de Morse y la agrupación de la amplitud por fase. PPC, en cambio, se aparta deliberadamente de la referencia: BOARD-FTD-PACC calcula un valor de bloqueo de fase cuyo sesgo hace que las frecuencias bajas aparezcan perfectamente acopladas con independencia de los datos, mientras que GammaLab calcula la consistencia de fase por pares sobre muestras independientes. Por ello, sus 84 pruebas de núcleo no se comparan con la referencia sino con señales sintéticas de verdad conocida: dos canales con desfase constante deben dar consistencia alta, y canales independientes, consistencia cercana a cero.

\subsection{Comparación visual con el usuario}
Como complemento de la comparación numérica, el usuario revisará las gráficas de cada análisis en ambos sistemas. La Tabla~\ref{tab:td_visual} queda preparada para registrar esa revisión.

\begin{center}
\small
\begin{longtable}{|>{\raggedright\arraybackslash}p{0.20\textwidth}|>{\raggedright\arraybackslash}p{0.48\textwidth}|>{\raggedright\arraybackslash}p{0.22\textwidth}|}
\hline
\textbf{Análisis} & \textbf{Observación del usuario} & \textbf{Resultado} \\
\hline
\endfirsthead
\hline
\textbf{Análisis} & \textbf{Observación del usuario} & \textbf{Resultado} \\
\hline
\endhead
\hline
\endfoot
\hline
\caption{Registro de la comparación visual con el usuario} \label{tab:td_visual} \\
\endlastfoot
\textit{Average} & \rule{0pt}{2.2em} & \casilla{} Equivalente \newline \casilla{} Diferencia aceptada \newline \casilla{} Diferencia no aceptada \\
\hline
ERP & \rule{0pt}{2.2em} & \casilla{} Equivalente \newline \casilla{} Diferencia aceptada \newline \casilla{} Diferencia no aceptada \\
\hline
FFT & \rule{0pt}{2.2em} & \casilla{} Equivalente \newline \casilla{} Diferencia aceptada \newline \casilla{} Diferencia no aceptada \\
\hline
\textit{FFT Average} & \rule{0pt}{2.2em} & \casilla{} Equivalente \newline \casilla{} Diferencia aceptada \newline \casilla{} Diferencia no aceptada \\
\hline
PSD & \rule{0pt}{2.2em} & \casilla{} Equivalente \newline \casilla{} Diferencia aceptada \newline \casilla{} Diferencia no aceptada \\
\hline
\textit{PSD Average} & \rule{0pt}{2.2em} & \casilla{} Equivalente \newline \casilla{} Diferencia aceptada \newline \casilla{} Diferencia no aceptada \\
\hline
\textit{Relative PSD} & \rule{0pt}{2.2em} & \casilla{} Equivalente \newline \casilla{} Diferencia aceptada \newline \casilla{} Diferencia no aceptada \\
\hline
\textit{Wavelet} & \rule{0pt}{2.2em} & \casilla{} Equivalente \newline \casilla{} Diferencia aceptada \newline \casilla{} Diferencia no aceptada \\
\hline
\textit{Wavelet Average} & \rule{0pt}{2.2em} & \casilla{} Equivalente \newline \casilla{} Diferencia aceptada \newline \casilla{} Diferencia no aceptada \\
\hline
PAC & \rule{0pt}{2.2em} & \casilla{} Equivalente \newline \casilla{} Diferencia aceptada \newline \casilla{} Diferencia no aceptada \\
\hline
PAC entre canales & \rule{0pt}{2.2em} & \casilla{} Equivalente \newline \casilla{} Diferencia aceptada \newline \casilla{} Diferencia no aceptada \\
\hline
\textit{Modulation Index} & \rule{0pt}{2.2em} & \casilla{} Equivalente \newline \casilla{} Diferencia aceptada \newline \casilla{} Diferencia no aceptada \\
\hline
PPC & \rule{0pt}{2.2em} & \casilla{} Equivalente \newline \casilla{} Diferencia aceptada \newline \casilla{} Diferencia no aceptada \\
\hline
\end{longtable}
\end{center}
\capitulo{PRUEBAS DE RENDIMIENTO Y DE REQUISITOS NO FUNCIONALES}
\label{sec:td_rendimiento}

\subsection{Objetivo}
Estas pruebas verifican los requisitos no funcionales del SRS que no dependen de la percepción del usuario: los límites de tiempo y de memoria (rendimiento), la estabilidad ante la concurrencia y la cancelación (confiabilidad) y las propiedades de la arquitectura que facilitan el mantenimiento (soporte). Los requisitos de usabilidad se evalúan con el usuario en la Sección~\ref{sec:td_usabilidad}.

\subsection{Implementación}
La verificación combina tres medios. Las pruebas automáticas comprueban los comportamientos que pueden expresarse como una aserción, por ejemplo que tras cerrar un proyecto no quede ninguna copia de la señal en memoria. Las 25 pruebas de rendimiento de la carpeta \textit{performance\_test}, todas aprobadas, miden el tiempo de las operaciones aisladas. Por último, los límites que dependen de la aplicación completa, como el tiempo de arranque o el de cambio de módulo, se midieron con \textit{time.perf\_counter} sobre la aplicación en funcionamiento en el equipo de la Tabla~\ref{tab:td_entorno}.

\begin{center}
\footnotesize
\begin{longtable}{|>{\raggedright\arraybackslash}p{0.22\textwidth}|>{\raggedright\arraybackslash}p{0.33\textwidth}|>{\centering\arraybackslash}p{0.09\textwidth}|>{\centering\arraybackslash}p{0.11\textwidth}|>{\raggedright\arraybackslash}p{0.10\textwidth}|}
\hline
\textbf{Suite} & \textbf{Qué verifica} & \textbf{Pruebas} & \textbf{Aprobadas} & \textbf{Resultado} \\
\hline
\endfirsthead
\hline
\textbf{Suite} & \textbf{Qué verifica} & \textbf{Pruebas} & \textbf{Aprobadas} & \textbf{Resultado} \\
\hline
\endhead
\hline
\endfoot
\hline
\caption{Resultados de las pruebas de rendimiento} \label{tab:td_perf} \\
\endlastfoot
\textit{test\_compute\_bench} & Tiempo de cómputo de los análisis (FFT, PSD, promedio, \textit{wavelet} y sus variantes) sobre una señal real. & 7 & 7 & Aprobada \\
\hline
\textit{test\_io\_bench} & Tiempo de carga de archivos, de corte de \textit{trials} y de las operaciones del almacén en memoria. & 4 & 4 & Aprobada \\
\hline
\textit{test\_kernel\_bench} & Costo del arranque: descubrimiento e instanciación de extensiones y registro en el núcleo. & 5 & 5 & Aprobada \\
\hline
\textit{test\_preprocessing\_bench} & Tiempo del filtro pasa-banda Butterworth que se aplica antes de los análisis. & 2 & 2 & Aprobada \\
\hline
\textit{test\_render\_bench} & Tiempo de construcción de los mapas de calor de VTK de \textit{wavelet}, \textit{wavelet} promedio y ERP. & 3 & 3 & Aprobada \\
\hline
\textit{test\_vtk\_adapters\_bench} & Tiempo de conversión vectorizada de señales y \textit{trials} a las tablas de VTK. & 4 & 4 & Aprobada \\
\hline
\end{longtable}
\end{center}

\subsection{Resultados de los requisitos de rendimiento}
La Tabla~\ref{tab:td_nfr_p} contrasta cada requisito de rendimiento con su verificación y su medición.

\begin{center}
\small
\begin{longtable}{|>{\raggedright\arraybackslash}p{0.13\textwidth}|>{\raggedright\arraybackslash}p{0.19\textwidth}|>{\raggedright\arraybackslash}p{0.40\textwidth}|>{\raggedright\arraybackslash}p{0.14\textwidth}|}
\hline
\textbf{Requisito} & \textbf{Criterio} & \textbf{Verificación y medición} & \textbf{Resultado} \\
\hline
\endfirsthead
\hline
\textbf{Requisito} & \textbf{Criterio} & \textbf{Verificación y medición} & \textbf{Resultado} \\
\hline
\endhead
\hline
\endfoot
\hline
\caption{Verificación de los requisitos de rendimiento} \label{tab:td_nfr_p} \\
\endlastfoot
NFR-P-01 (R11) & Liberar la señal al cerrar el proyecto & Pruebas automáticas y medición: tras cerrar el proyecto no queda ninguna copia viva de la señal (\textit{test\_liberar\_senal}). & Cumple. \\
\hline
NFR-P-02 (R45) & Liberar recursos gráficos al salir del módulo & Pruebas automáticas de la liberación. Medición: al recorrer los módulos de análisis, el proceso pasa de 3.118 MB sin liberar a 2.585 MB con la configuración actual (\textit{test\_liberar\_ventanas}). & Cumple. \\
\hline
NFR-P-03 (R47) & Cargar el archivo en 2 s o menos & Pruebas automáticas de la lectura sin duplicación. Medición con un archivo de 2 canales de 1,8 millones de muestras: 0,04 s de lectura y 0,2 s hasta verlo en pantalla (límite: 2 s) (\textit{test\_io\_bench}, \textit{test\_fileio\_sin\_duplicar}). & Cumple. \\
\hline
NFR-P-04 (R48) & \textit{Wavelet} en 10 s o menos & Medición: de 1,1 a 1,4 s para un \textit{trial} de 30.500 muestras con los parámetros por defecto (límite: 10 s) (\textit{test\_compute\_bench}). & Cumple. \\
\hline
NFR-P-05 (R49) & Graficar 3 canales sin degradar la interacción & Pruebas automáticas del submuestreo. Medición en el equipo de desarrollo, con tres canales de 1,8 millones de muestras: mediana de 29 a 44 FPS al hacer \textit{zoom} y de 24 a 36 FPS al desplazar (límite: 24 FPS); el peor cuadro de cada corrida varió entre 16 y 31 FPS según la carga del equipo. Reducir los puntos dibujados hasta 250 no modificó la tasa, por lo que el tiempo restante es fijo del renderizado (\textit{test\_submuestreo}, \textit{test\_render\_bench}). & Cumple en la mediana; el peor cuadro de algunas corridas baja de 24 FPS. \\
\hline
NFR-P-06 (R50) & Cambiar de módulo en 500 ms o menos & Medición: el cambio de módulo tarda entre 10 y 310 ms con una señal de 1,7 millones de muestras (límite: 500 ms). & Cumple. \\
\hline
NFR-P-07 (R51) & Exportar 100\,000 filas en 2 s o menos & Pruebas automáticas del escritor. Medición con 100.000 filas: 0,39 s en Excel y 0,11 s en CSV (límite: 2 s) (\textit{test\_xlsx\_writer}). & Cumple. \\
\hline
NFR-P-08 (R52) & Liberar la señal anterior en menos de 1 s & Pruebas automáticas y medición: la señal anterior se libera en 243 a 353 ms; tras cuatro cargas seguidas queda una sola señal viva y la memoria sube 178 MB en vez de 631 MB (\textit{test\_liberar\_senal}). & Cumple. \\
\hline
NFR-P-09 (R53) & Cierre en menos de 3 s con recursos liberados & Medición: la ventana se cierra en 172 ms y el proceso termina en 0,63 s (límite: 3 s). & Cumple. \\
\hline
NFR-P-10 (R55) & Interfaz responsiva (0,3 s o menos) & Medición: la interfaz se congela entre 270 y 480 ms la primera vez que se abre PAC, según la corrida, y entre 110 y 200 ms en adelante (límite: 300 ms). & Cumple parcialmente: la primera apertura de PAC puede superar el límite. \\
\hline
NFR-P-11 (R56) & Promedio de todos los \textit{trials} en 1 s o menos & Medición: de 0,32 a 0,41 s para 50 \textit{trials} de 30.500 muestras (límite: 1 s) (\textit{test\_compute\_bench}). & Cumple. \\
\hline
NFR-P-12 (R57) & Exportar una gráfica en 1 s (3 s si es vectorial) & Pruebas automáticas de la exportación. Medición: PNG en 0,26 s, TIFF en 0,06 s y EMF en 0,12 s (límites: 1 s y 3 s) (\textit{test\_export}, \textit{test\_emf\_export}). & Cumple. \\
\hline
NFR-P-13 (R59) & Conversión en bloque a las estructuras de VTK & Medición: 1,8 millones de filas convertidas en 0,03 s. Prueba automática del llenado vectorizado del mapa de calor (\textit{test\_vtk\_adapters\_bench}, \textit{test\_erp\_vtk\_llenado}). & Cumple. \\
\hline
NFR-P-14 (R71) & Cómputo pesado por el servicio de tareas & Pruebas automáticas del patrón y de las extensiones que lo usan (\textit{test\_run\_task}, \textit{test\_task\_service}, \textit{test\_artifact\_remove\_orquestador}, \textit{test\_ppc\_plugin\_orquestado}). & Cumple. \\
\hline
NFR-P-15 (R73) & Hasta 2000 puntos por serie y 200 curvas & Pruebas automáticas del submuestreo: como máximo 2000 puntos por serie, sin alterar los datos (\textit{test\_submuestreo}, \textit{test\_correcciones\_interaccion}). & Cumple. \\
\hline
NFR-P-16 (R74) & PAC por el servicio de tareas & Medición: un comodulograma de 20×20 pares tarda entre 0,2 y 0,5 s (límite: 30 s); la cancelación se comprueba en la aplicación. & Cumple. \\
\hline
NFR-P-17 (R81) & Listo en 8 s o menos & Medición: 3,1 s hasta la ventana lista (límite: 8 s) y prueba del programa instalado (\textit{test\_kernel\_bench}, \textit{prueba\_congelado}). & Cumple. \\
\hline
\end{longtable}
\end{center}

Dos requisitos presentan matices. En NFR-P-05 la mediana de cuadros por segundo supera el límite de 24 en todas las corridas, pero el peor cuadro de algunas corridas cae por debajo de ese valor. En NFR-P-10 la interfaz responde dentro de los 300\,ms en el uso normal, pero la primera apertura del módulo PAC en una sesión puede congelarla entre 270 y 480\,ms, por la carga inicial de sus componentes.

\subsection{Resultados de los requisitos de confiabilidad y soporte}

\begin{center}
\small
\begin{longtable}{|>{\raggedright\arraybackslash}p{0.13\textwidth}|>{\raggedright\arraybackslash}p{0.19\textwidth}|>{\raggedright\arraybackslash}p{0.40\textwidth}|>{\raggedright\arraybackslash}p{0.14\textwidth}|}
\hline
\textbf{Requisito} & \textbf{Criterio} & \textbf{Verificación y medición} & \textbf{Resultado} \\
\hline
\endfirsthead
\hline
\textbf{Requisito} & \textbf{Criterio} & \textbf{Verificación y medición} & \textbf{Resultado} \\
\hline
\endhead
\hline
\endfoot
\hline
\caption{Verificación de los requisitos de confiabilidad y soporte} \label{tab:td_nfr_rs} \\
\endlastfoot
NFR-R-01 (R46) & Terminar los hilos al cancelar o cerrar & Pruebas automáticas de la cancelación cooperativa, del desligado a los dos segundos y del cambio de sección (\textit{test\_task\_service}, \textit{test\_cancelacion\_al\_cambiar\_seccion}). & Cumple. \\
\hline
NFR-R-02 (R54) & Sincronizar el estado compartido & Pruebas automáticas de la concurrencia sobre el almacén (\textit{test\_data\_store\_concurrencia}). & Cumple. \\
\hline
NFR-R-03 (R58) & Retirar los detectores de eventos de VTK & Pruebas automáticas del retiro, la reposición y el reemplazo del menú; verificación en la aplicación con cinco cálculos seguidos (\textit{test\_vtk\_observers}). & Cumple. \\
\hline
NFR-R-04 (R72) & Reutilizar la representación gráfica & Pruebas automáticas que verifican que la segunda petición devuelve la misma tabla, que corregir o restaurar la frecuencia la descarta y que la tabla nueva trae el eje corregido sin tocar la anterior (\textit{test\_cache\_tabla\_vtk}, \textit{test\_cu018\_sampling\_rate}). & Cumple. \\
\hline
NFR-R-05 (R79) & No modificar la señal ni los \textit{trials} originales & Pruebas automáticas de los conjuntos derivados (\textit{test\_trials\_derivados}). & Cumple. \\
\hline
NFR-S-01 (R39) & Ejecutable instalable que registra la extensión del proyecto & Pruebas automáticas de las rutas del programa instalado y prueba del ejecutable, además de la instalación real verificada en un equipo (\textit{test\_instalacion\_rutas}, \textit{prueba\_congelado}). & Cumple. \\
\hline
NFR-S-02 (R67) & Vía pública única sobre los \textit{trials} & Pruebas automáticas de los \textit{trials} originales, derivados y por canal (\textit{test\_trials\_derivados}, \textit{test\_trials\_por\_canal}). & Cumple. \\
\hline
NFR-S-03 (R69) & Servicios compartidos registrados en el núcleo & Prueba automática que recorre el código y falla si otro módulo crea uno de los servicios, y pruebas de cada servicio. Verificado en la aplicación con una exportación y una medición reales (\textit{test\_servicios\_compartidos}). & Cumple. \\
\hline
NFR-S-04 (R70) & Cómputo pesado por un servicio del núcleo & Pruebas automáticas del servicio de tareas y del patrón de uso (\textit{test\_task\_service}, \textit{test\_run\_task}). & Cumple. \\
\hline
NFR-S-05 (R77) & Nuevo \textit{plugin} sin modificar el núcleo & Pruebas automáticas con una extensión de prueba que se guarda, se restaura y se ordena sin tocar el núcleo (\textit{test\_plugin\_sin\_tocar\_nucleo}). & Cumple; una de sus cuatro pruebas falla por estar desactualizada (Sección~\ref{sec:td_resumen}). \\
\hline
\end{longtable}
\end{center}
\capitulo{PRUEBAS DE USABILIDAD}
\label{sec:td_usabilidad}

\subsection{Objetivo}
Estas pruebas evalúan la facilidad con que los investigadores del laboratorio completan las tareas habituales de su flujo de trabajo en GammaLab 2.0, y si la reorganización de la interfaz resuelve los problemas observados en las pruebas de usabilidad de la versión 1.0. Verifican, en particular, los requisitos de usabilidad NFR-U-01 a NFR-U-09.

\subsection{Metodología}
Cada sesión es individual y la conduce un moderador del equipo. El participante recibe la descripción de cada tarea, la realiza sin instrucciones adicionales y se le pide expresar en voz alta lo que piensa mientras trabaja. El moderador solo interviene si el participante solicita ayuda, y registra cada solicitud. Por cada caso se registran cuatro métricas: si la tarea se completó, el tiempo empleado, los errores cometidos y las solicitudes de ayuda. Al terminar los casos, el participante diligencia el cuestionario de satisfacción SUS (\textit{System Usability Scale}) \cite{brooke1996}.

\subsection{Entorno de la prueba}

\begin{center}
\small
\begin{tabularx}{\textwidth}{|X|X|X|X|X|}
\hline
\textbf{Equipo} & \textbf{Procesador} & \textbf{RAM} & \textbf{Sistema operativo} & \textbf{Versión de GammaLab} \\
\hline
\rule{0pt}{2em} & & & & 2.0, instalada con \textit{GammaLab-Setup-2.0.0.exe} \\
\hline
\end{tabularx}
\captionof{table}{Equipo empleado en las pruebas de usabilidad}
\label{tab:td_ut_entorno}
\end{center}

Los datos de prueba son la señal \textit{17308005.abf} y, a elección del participante, señales propias del laboratorio con los parámetros de su protocolo.

\subsection{Casos de prueba}

\begin{center}
\small
\begin{longtable}{|>{\raggedright\arraybackslash}p{0.27\textwidth}|>{\raggedright\arraybackslash}p{0.63\textwidth}|}
\hline
\endfirsthead
\hline
\endhead
\hline
\endfoot
\hline
\caption{Caso de prueba de usabilidad UT-01} \label{tab:ut-01} \\
\endlastfoot
\textbf{Id} & UT-01 \\
\hline
\textbf{Caso de prueba} & Primer arranque y orientación en la interfaz \\
\hline
\textbf{Caso de uso} & CU-011, CU-008, CU-010 \\
\hline
\textbf{Requisitos relacionados} & NFR-U-04, NFR-U-05, NFR-U-07, NFR-U-08, NFR-U-09 \\
\hline
\textbf{Descripción} & Evaluar si un investigador que abre GammaLab por primera vez comprende la organización de la interfaz y encuentra las funcionalidades sin ayuda externa. \\
\hline
\textbf{Datos y acciones de entrada} & Abrir GammaLab desde el acceso directo; leer el panel de bienvenida; recorrer las secciones \textit{Home}, \textit{Preprocessing} y \textit{Analysis}; situar el cursor sobre tres controles de cualquier módulo y leer su descripción emergente. \\
\hline
\textbf{Resultado esperado} & La pantalla de inicio informa el avance de la carga; el panel de bienvenida aparece solo en el primer arranque; cada funcionalidad tiene icono y cada control muestra su descripción emergente. \\
\hline
\textbf{Procedimientos especiales} & El moderador no ofrece ayuda salvo que el participante la solicite; cada solicitud se registra. \\
\hline
\textbf{Dependencias} & Ninguna \\
\hline
\multicolumn{2}{|c|}{\textbf{Información para el seguimiento}} \\
\hline
\textbf{Resultado obtenido} & \rule{0pt}{2.2em} \\
\hline
\textbf{Tiempo empleado (min)} &  \\
\hline
\textbf{Errores cometidos} &  \\
\hline
\textbf{Solicitudes de ayuda} &  \\
\hline
\textbf{Estado} & \casilla{} Superado \qquad \casilla{} Superado con dificultad \qquad \casilla{} No superado \\
\hline
\textbf{Fecha de ejecución} &  \\
\hline
\textbf{Observaciones} & \rule{0pt}{2.2em} \\
\hline
\end{longtable}
\end{center}
\begin{center}
\small
\begin{longtable}{|>{\raggedright\arraybackslash}p{0.27\textwidth}|>{\raggedright\arraybackslash}p{0.63\textwidth}|}
\hline
\endfirsthead
\hline
\endhead
\hline
\endfoot
\hline
\caption{Caso de prueba de usabilidad UT-02} \label{tab:ut-02} \\
\endlastfoot
\textbf{Id} & UT-02 \\
\hline
\textbf{Caso de prueba} & Cargar una señal e inspeccionarla \\
\hline
\textbf{Caso de uso} & CU-003, CU-012 \\
\hline
\textbf{Requisitos relacionados} & FR-IO-02, FR-UI-04, FR-UI-07, NFR-P-03 \\
\hline
\textbf{Descripción} & Evaluar la facilidad para cargar un registro e inspeccionar su forma, amplitud y tiempo. \\
\hline
\textbf{Datos y acciones de entrada} & Cargar el archivo \textit{17308005.abf} con \textit{File} > \textit{Open signal}; elegir los canales visibles; hacer \textit{zoom} en cada eje por separado; leer el valor X/Y de un punto de la señal. \\
\hline
\textbf{Resultado esperado} & La señal se muestra en 2 s o menos; el \textit{zoom} actúa sobre el eje elegido y la lectura del cursor muestra el valor del punto de la señal más cercano. \\
\hline
\textbf{Procedimientos especiales} & El moderador no ofrece ayuda salvo que el participante la solicite; cada solicitud se registra. \\
\hline
\textbf{Dependencias} & Ninguna \\
\hline
\multicolumn{2}{|c|}{\textbf{Información para el seguimiento}} \\
\hline
\textbf{Resultado obtenido} & \rule{0pt}{2.2em} \\
\hline
\textbf{Tiempo empleado (min)} &  \\
\hline
\textbf{Errores cometidos} &  \\
\hline
\textbf{Solicitudes de ayuda} &  \\
\hline
\textbf{Estado} & \casilla{} Superado \qquad \casilla{} Superado con dificultad \qquad \casilla{} No superado \\
\hline
\textbf{Fecha de ejecución} &  \\
\hline
\textbf{Observaciones} & \rule{0pt}{2.2em} \\
\hline
\end{longtable}
\end{center}
\begin{center}
\small
\begin{longtable}{|>{\raggedright\arraybackslash}p{0.27\textwidth}|>{\raggedright\arraybackslash}p{0.63\textwidth}|}
\hline
\endfirsthead
\hline
\endhead
\hline
\endfoot
\hline
\caption{Caso de prueba de usabilidad UT-03} \label{tab:ut-03} \\
\endlastfoot
\textbf{Id} & UT-03 \\
\hline
\textbf{Caso de prueba} & Generar, revisar y descartar \textit{trials} \\
\hline
\textbf{Caso de uso} & CU-004, CU-006 \\
\hline
\textbf{Requisitos relacionados} & FR-UI-01, FR-UI-02, FR-UI-05, FR-UI-08, FR-UI-11 \\
\hline
\textbf{Descripción} & Evaluar si el usuario segmenta la señal con los parámetros de su protocolo y depura el conjunto de \textit{trials}. \\
\hline
\textbf{Datos y acciones de entrada} & En \textit{Preprocessing} > \textit{Trials}, generar los \textit{trials} del canal CA1 con el canal de estímulo IN7 y el umbral en mV de su protocolo; recorrerlos con \textit{Previous} y \textit{Next}; descartar dos y recuperar uno. \\
\hline
\textbf{Resultado esperado} & Los \textit{trials} se generan según los parámetros; el umbral se expresa en mV; el indicador de la barra de título muestra en todo momento el \textit{trial} en curso y su estado. \\
\hline
\textbf{Procedimientos especiales} & El moderador no ofrece ayuda salvo que el participante la solicite; cada solicitud se registra. \\
\hline
\textbf{Dependencias} & UT-02 \\
\hline
\multicolumn{2}{|c|}{\textbf{Información para el seguimiento}} \\
\hline
\textbf{Resultado obtenido} & \rule{0pt}{2.2em} \\
\hline
\textbf{Tiempo empleado (min)} &  \\
\hline
\textbf{Errores cometidos} &  \\
\hline
\textbf{Solicitudes de ayuda} &  \\
\hline
\textbf{Estado} & \casilla{} Superado \qquad \casilla{} Superado con dificultad \qquad \casilla{} No superado \\
\hline
\textbf{Fecha de ejecución} &  \\
\hline
\textbf{Observaciones} & \rule{0pt}{2.2em} \\
\hline
\end{longtable}
\end{center}
\begin{center}
\small
\begin{longtable}{|>{\raggedright\arraybackslash}p{0.27\textwidth}|>{\raggedright\arraybackslash}p{0.63\textwidth}|}
\hline
\endfirsthead
\hline
\endhead
\hline
\endfoot
\hline
\caption{Caso de prueba de usabilidad UT-04} \label{tab:ut-04} \\
\endlastfoot
\textbf{Id} & UT-04 \\
\hline
\textbf{Caso de prueba} & Remover un artefacto con medidas \\
\hline
\textbf{Caso de uso} & CU-007 \\
\hline
\textbf{Requisitos relacionados} & FR-UI-13, NFR-R-05 \\
\hline
\textbf{Descripción} & Evaluar si el usuario corrige un artefacto decidiendo a partir de las medidas del segmento. \\
\hline
\textbf{Datos y acciones de entrada} & En \textit{Preprocessing} > \textit{Remove Artifact}, ubicar un \textit{trial} con artefacto, delimitar el tramo con \textit{Point A} y \textit{Point B}, revisar las medidas mostradas y confirmar la corrección; restaurar después los \textit{trials} originales. \\
\hline
\textbf{Resultado esperado} & Las medidas del segmento se muestran antes de modificarlo; la corrección se aplica solo al confirmar y la restauración devuelve los \textit{trials} originales. \\
\hline
\textbf{Procedimientos especiales} & El moderador no ofrece ayuda salvo que el participante la solicite; cada solicitud se registra. \\
\hline
\textbf{Dependencias} & UT-02 \\
\hline
\multicolumn{2}{|c|}{\textbf{Información para el seguimiento}} \\
\hline
\textbf{Resultado obtenido} & \rule{0pt}{2.2em} \\
\hline
\textbf{Tiempo empleado (min)} &  \\
\hline
\textbf{Errores cometidos} &  \\
\hline
\textbf{Solicitudes de ayuda} &  \\
\hline
\textbf{Estado} & \casilla{} Superado \qquad \casilla{} Superado con dificultad \qquad \casilla{} No superado \\
\hline
\textbf{Fecha de ejecución} &  \\
\hline
\textbf{Observaciones} & \rule{0pt}{2.2em} \\
\hline
\end{longtable}
\end{center}
\begin{center}
\small
\begin{longtable}{|>{\raggedright\arraybackslash}p{0.27\textwidth}|>{\raggedright\arraybackslash}p{0.63\textwidth}|}
\hline
\endfirsthead
\hline
\endhead
\hline
\endfoot
\hline
\caption{Caso de prueba de usabilidad UT-05} \label{tab:ut-05} \\
\endlastfoot
\textbf{Id} & UT-05 \\
\hline
\textbf{Caso de prueba} & Filtrar los \textit{trials} \\
\hline
\textbf{Caso de uso} & CU-016 \\
\hline
\textbf{Requisitos relacionados} & FR-PAC-20, NFR-R-05 \\
\hline
\textbf{Descripción} & Evaluar la claridad del flujo de vista previa y aplicación del filtro. \\
\hline
\textbf{Datos y acciones de entrada} & En \textit{Preprocessing} > \textit{Filter}, fijar una banda de 4 a 8 Hz y orden 8; pulsar \textit{Preview}; aplicar con \textit{Filter}. \\
\hline
\textbf{Resultado esperado} & La vista previa no modifica los \textit{trials}; la aplicación crea un conjunto derivado que usan los análisis posteriores. \\
\hline
\textbf{Procedimientos especiales} & El moderador no ofrece ayuda salvo que el participante la solicite; cada solicitud se registra. \\
\hline
\textbf{Dependencias} & UT-02 \\
\hline
\multicolumn{2}{|c|}{\textbf{Información para el seguimiento}} \\
\hline
\textbf{Resultado obtenido} & \rule{0pt}{2.2em} \\
\hline
\textbf{Tiempo empleado (min)} &  \\
\hline
\textbf{Errores cometidos} &  \\
\hline
\textbf{Solicitudes de ayuda} &  \\
\hline
\textbf{Estado} & \casilla{} Superado \qquad \casilla{} Superado con dificultad \qquad \casilla{} No superado \\
\hline
\textbf{Fecha de ejecución} &  \\
\hline
\textbf{Observaciones} & \rule{0pt}{2.2em} \\
\hline
\end{longtable}
\end{center}
\begin{center}
\small
\begin{longtable}{|>{\raggedright\arraybackslash}p{0.27\textwidth}|>{\raggedright\arraybackslash}p{0.63\textwidth}|}
\hline
\endfirsthead
\hline
\endhead
\hline
\endfoot
\hline
\caption{Caso de prueba de usabilidad UT-06} \label{tab:ut-06} \\
\endlastfoot
\textbf{Id} & UT-06 \\
\hline
\textbf{Caso de prueba} & Análisis PAC de un canal \\
\hline
\textbf{Caso de uso} & CU-001 \\
\hline
\textbf{Requisitos relacionados} & FR-PAC-01 a FR-PAC-11, NFR-U-02, NFR-U-03 \\
\hline
\textbf{Descripción} & Evaluar si el usuario configura y lee el análisis PAC en sus distintos modos. \\
\hline
\textbf{Datos y acciones de entrada} & En \textit{Analysis} > PAC, fijar la banda de fase de 3 a 8 Hz y la de amplitud de 25 a 100 Hz; calcular en los modos \textit{PAC}, \textit{PAC Average} y \textit{PAC Psel}; identificar la fase preferente en la brújula del vector medio. \\
\hline
\textbf{Resultado esperado} & Cada cálculo muestra la barra de progreso y su resultado en la misma pantalla; el mapa fase-frecuencia, el histograma y el vector medio son legibles para el usuario. \\
\hline
\textbf{Procedimientos especiales} & El moderador no ofrece ayuda salvo que el participante la solicite; cada solicitud se registra. \\
\hline
\textbf{Dependencias} & UT-02 \\
\hline
\multicolumn{2}{|c|}{\textbf{Información para el seguimiento}} \\
\hline
\textbf{Resultado obtenido} & \rule{0pt}{2.2em} \\
\hline
\textbf{Tiempo empleado (min)} &  \\
\hline
\textbf{Errores cometidos} &  \\
\hline
\textbf{Solicitudes de ayuda} &  \\
\hline
\textbf{Estado} & \casilla{} Superado \qquad \casilla{} Superado con dificultad \qquad \casilla{} No superado \\
\hline
\textbf{Fecha de ejecución} &  \\
\hline
\textbf{Observaciones} & \rule{0pt}{2.2em} \\
\hline
\end{longtable}
\end{center}
\begin{center}
\small
\begin{longtable}{|>{\raggedright\arraybackslash}p{0.27\textwidth}|>{\raggedright\arraybackslash}p{0.63\textwidth}|}
\hline
\endfirsthead
\hline
\endhead
\hline
\endfoot
\hline
\caption{Caso de prueba de usabilidad UT-07} \label{tab:ut-07} \\
\endlastfoot
\textbf{Id} & UT-07 \\
\hline
\textbf{Caso de prueba} & PAC entre canales y PPC \\
\hline
\textbf{Caso de uso} & CU-002 \\
\hline
\textbf{Requisitos relacionados} & FR-PAC-12 a FR-PAC-16 \\
\hline
\textbf{Descripción} & Evaluar si el usuario completa sin ayuda los análisis que combinan dos canales. \\
\hline
\textbf{Datos y acciones de entrada} & Generar los \textit{trials} de dos canales con los mismos parámetros; calcular \textit{PAC Psel inter-canal} eligiendo el canal de amplitud; calcular PPC sobre el mismo par de canales. \\
\hline
\textbf{Resultado esperado} & El usuario identifica qué canal aporta la fase y cuál la amplitud; los mapas se rotulan con el canal de origen y las celdas no medibles de PPC quedan en blanco. \\
\hline
\textbf{Procedimientos especiales} & El moderador no ofrece ayuda salvo que el participante la solicite; cada solicitud se registra. \\
\hline
\textbf{Dependencias} & UT-02 \\
\hline
\multicolumn{2}{|c|}{\textbf{Información para el seguimiento}} \\
\hline
\textbf{Resultado obtenido} & \rule{0pt}{2.2em} \\
\hline
\textbf{Tiempo empleado (min)} &  \\
\hline
\textbf{Errores cometidos} &  \\
\hline
\textbf{Solicitudes de ayuda} &  \\
\hline
\textbf{Estado} & \casilla{} Superado \qquad \casilla{} Superado con dificultad \qquad \casilla{} No superado \\
\hline
\textbf{Fecha de ejecución} &  \\
\hline
\textbf{Observaciones} & \rule{0pt}{2.2em} \\
\hline
\end{longtable}
\end{center}
\begin{center}
\small
\begin{longtable}{|>{\raggedright\arraybackslash}p{0.27\textwidth}|>{\raggedright\arraybackslash}p{0.63\textwidth}|}
\hline
\endfirsthead
\hline
\endhead
\hline
\endfoot
\hline
\caption{Caso de prueba de usabilidad UT-08} \label{tab:ut-08} \\
\endlastfoot
\textbf{Id} & UT-08 \\
\hline
\textbf{Caso de prueba} & Tomar medidas y exportar resultados \\
\hline
\textbf{Caso de uso} & CU-005, CU-017 \\
\hline
\textbf{Requisitos relacionados} & FR-IO-05, FR-IO-06, FR-IO-07, FR-WV-03 \\
\hline
\textbf{Descripción} & Evaluar la facilidad para medir sobre una gráfica y llevar los resultados a otras herramientas. \\
\hline
\textbf{Datos y acciones de entrada} & Medir una pendiente y una amplitud con el menú contextual; exportar la tabla de resultados a Excel; exportar un escalograma en EMF y su tabla en CSV. \\
\hline
\textbf{Resultado esperado} & Las medidas aparecen en el panel \textit{Results}; los archivos se generan con el nombre de la señal y abren correctamente en las herramientas de ofimática. \\
\hline
\textbf{Procedimientos especiales} & El moderador no ofrece ayuda salvo que el participante la solicite; cada solicitud se registra. \\
\hline
\textbf{Dependencias} & UT-02 \\
\hline
\multicolumn{2}{|c|}{\textbf{Información para el seguimiento}} \\
\hline
\textbf{Resultado obtenido} & \rule{0pt}{2.2em} \\
\hline
\textbf{Tiempo empleado (min)} &  \\
\hline
\textbf{Errores cometidos} &  \\
\hline
\textbf{Solicitudes de ayuda} &  \\
\hline
\textbf{Estado} & \casilla{} Superado \qquad \casilla{} Superado con dificultad \qquad \casilla{} No superado \\
\hline
\textbf{Fecha de ejecución} &  \\
\hline
\textbf{Observaciones} & \rule{0pt}{2.2em} \\
\hline
\end{longtable}
\end{center}
\begin{center}
\small
\begin{longtable}{|>{\raggedright\arraybackslash}p{0.27\textwidth}|>{\raggedright\arraybackslash}p{0.63\textwidth}|}
\hline
\endfirsthead
\hline
\endhead
\hline
\endfoot
\hline
\caption{Caso de prueba de usabilidad UT-09} \label{tab:ut-09} \\
\endlastfoot
\textbf{Id} & UT-09 \\
\hline
\textbf{Caso de prueba} & Cancelar un cálculo pesado \\
\hline
\textbf{Caso de uso} & CU-014, CU-020 \\
\hline
\textbf{Requisitos relacionados} & NFR-U-03, NFR-R-01 \\
\hline
\textbf{Descripción} & Evaluar si el usuario reconoce un cálculo en curso y lo detiene cuando lo necesita. \\
\hline
\textbf{Datos y acciones de entrada} & Lanzar \textit{Wavelet Average} sobre todos los \textit{trials}; cancelarlo desde la barra de progreso; cambiar de módulo mientras otro cálculo está en curso. \\
\hline
\textbf{Resultado esperado} & La barra de progreso muestra el avance y permite cancelar; tras la cancelación la aplicación vuelve a su estado anterior sin resultados parciales. \\
\hline
\textbf{Procedimientos especiales} & El moderador no ofrece ayuda salvo que el participante la solicite; cada solicitud se registra. \\
\hline
\textbf{Dependencias} & UT-02 \\
\hline
\multicolumn{2}{|c|}{\textbf{Información para el seguimiento}} \\
\hline
\textbf{Resultado obtenido} & \rule{0pt}{2.2em} \\
\hline
\textbf{Tiempo empleado (min)} &  \\
\hline
\textbf{Errores cometidos} &  \\
\hline
\textbf{Solicitudes de ayuda} &  \\
\hline
\textbf{Estado} & \casilla{} Superado \qquad \casilla{} Superado con dificultad \qquad \casilla{} No superado \\
\hline
\textbf{Fecha de ejecución} &  \\
\hline
\textbf{Observaciones} & \rule{0pt}{2.2em} \\
\hline
\end{longtable}
\end{center}
\begin{center}
\small
\begin{longtable}{|>{\raggedright\arraybackslash}p{0.27\textwidth}|>{\raggedright\arraybackslash}p{0.63\textwidth}|}
\hline
\endfirsthead
\hline
\endhead
\hline
\endfoot
\hline
\caption{Caso de prueba de usabilidad UT-10} \label{tab:ut-10} \\
\endlastfoot
\textbf{Id} & UT-10 \\
\hline
\textbf{Caso de prueba} & Guardar, cerrar y abrir un proyecto \\
\hline
\textbf{Caso de uso} & CU-005, CU-021, CU-023 \\
\hline
\textbf{Requisitos relacionados} & FR-IO-03, FR-IO-04, NFR-U-06 \\
\hline
\textbf{Descripción} & Evaluar si el usuario conserva y retoma su trabajo mediante proyectos. \\
\hline
\textbf{Datos y acciones de entrada} & Guardar el trabajo con \textit{File} > \textit{Save as}; cerrar el proyecto con \textit{File} > \textit{Close project}; abrirlo de nuevo con \textit{File} > \textit{Open project}. \\
\hline
\textbf{Resultado esperado} & El cierre se comunica y pide confirmación si hay cambios sin guardar; al abrir, se restauran la señal, los \textit{trials}, los descartes, las medidas y los parámetros de los análisis. \\
\hline
\textbf{Procedimientos especiales} & El moderador no ofrece ayuda salvo que el participante la solicite; cada solicitud se registra. \\
\hline
\textbf{Dependencias} & UT-02 \\
\hline
\multicolumn{2}{|c|}{\textbf{Información para el seguimiento}} \\
\hline
\textbf{Resultado obtenido} & \rule{0pt}{2.2em} \\
\hline
\textbf{Tiempo empleado (min)} &  \\
\hline
\textbf{Errores cometidos} &  \\
\hline
\textbf{Solicitudes de ayuda} &  \\
\hline
\textbf{Estado} & \casilla{} Superado \qquad \casilla{} Superado con dificultad \qquad \casilla{} No superado \\
\hline
\textbf{Fecha de ejecución} &  \\
\hline
\textbf{Observaciones} & \rule{0pt}{2.2em} \\
\hline
\end{longtable}
\end{center}

\subsection{Consolidado de resultados}

\begin{center}
\footnotesize
\begin{longtable}{|>{\raggedright\arraybackslash}p{0.09\textwidth}|>{\centering\arraybackslash}p{0.13\textwidth}|>{\centering\arraybackslash}p{0.12\textwidth}|>{\centering\arraybackslash}p{0.09\textwidth}|>{\centering\arraybackslash}p{0.09\textwidth}|>{\raggedright\arraybackslash}p{0.28\textwidth}|}
\hline
\textbf{Caso} & \textbf{Completado} & \textbf{Tiempo (min)} & \textbf{Errores} & \textbf{Ayudas} & \textbf{Observaciones} \\
\hline
\endfirsthead
\hline
\textbf{Caso} & \textbf{Completado} & \textbf{Tiempo (min)} & \textbf{Errores} & \textbf{Ayudas} & \textbf{Observaciones} \\
\hline
\endhead
\hline
\endfoot
\hline
\caption{Consolidado de las pruebas de usabilidad} \label{tab:td_ut_consolidado} \\
\endlastfoot
UT-01 & & & & & \rule{0pt}{1.6em} \\
\hline
UT-02 & & & & & \rule{0pt}{1.6em} \\
\hline
UT-03 & & & & & \rule{0pt}{1.6em} \\
\hline
UT-04 & & & & & \rule{0pt}{1.6em} \\
\hline
UT-05 & & & & & \rule{0pt}{1.6em} \\
\hline
UT-06 & & & & & \rule{0pt}{1.6em} \\
\hline
UT-07 & & & & & \rule{0pt}{1.6em} \\
\hline
UT-08 & & & & & \rule{0pt}{1.6em} \\
\hline
UT-09 & & & & & \rule{0pt}{1.6em} \\
\hline
UT-10 & & & & & \rule{0pt}{1.6em} \\
\hline
\end{longtable}
\end{center}

\subsection{Cuestionario de satisfacción}
El participante califica cada afirmación de la Tabla~\ref{tab:td_sus} de 1 (totalmente en desacuerdo) a 5 (totalmente de acuerdo). La puntuación SUS se obtiene restando 1 a la calificación de las afirmaciones impares y restando de 5 la calificación de las pares, sumando los diez resultados y multiplicando la suma por 2,5, lo que da un valor entre 0 y 100 \cite{brooke1996}.

\begin{center}
\small
\begin{longtable}{|>{\centering\arraybackslash}p{0.04\textwidth}|>{\raggedright\arraybackslash}p{0.56\textwidth}|c|c|c|c|c|}
\hline
\textbf{N.\textsuperscript{o}} & \textbf{Afirmación} & \textbf{1} & \textbf{2} & \textbf{3} & \textbf{4} & \textbf{5} \\
\hline
\endfirsthead
\hline
\textbf{N.\textsuperscript{o}} & \textbf{Afirmación} & \textbf{1} & \textbf{2} & \textbf{3} & \textbf{4} & \textbf{5} \\
\hline
\endhead
\hline
\endfoot
\hline
\caption{Cuestionario SUS} \label{tab:td_sus} \\
\endlastfoot
1 & Creo que me gustaría usar este sistema con frecuencia. & & & & & \\
\hline
2 & Encontré el sistema innecesariamente complejo. & & & & & \\
\hline
3 & Pensé que el sistema era fácil de usar. & & & & & \\
\hline
4 & Creo que necesitaría el apoyo de una persona con conocimientos técnicos para poder usar este sistema. & & & & & \\
\hline
5 & Encontré que las distintas funciones del sistema estaban bien integradas. & & & & & \\
\hline
6 & Pensé que había demasiada inconsistencia en el sistema. & & & & & \\
\hline
7 & Imagino que la mayoría de las personas aprenderían a usar este sistema muy rápidamente. & & & & & \\
\hline
8 & Encontré el sistema muy engorroso de usar. & & & & & \\
\hline
9 & Me sentí muy seguro al usar el sistema. & & & & & \\
\hline
10 & Necesité aprender muchas cosas antes de poder empezar a usar el sistema. & & & & & \\
\hline
\multicolumn{2}{|r|}{\textbf{Puntuación SUS (0 a 100)}} & \multicolumn{5}{c|}{} \\
\hline
\end{longtable}
\end{center}

\subsection{Conclusiones de las pruebas de usabilidad}
\pendiente[Se redactarán una vez ejecutadas las sesiones con el usuario, a partir del consolidado de resultados y del cuestionario de satisfacción.]
\capitulo{PRUEBAS DE ACEPTACIÓN DEL USUARIO (UAT)}
\label{sec:td_uat}

\subsection{Objetivo}
El objetivo de estas pruebas es verificar que GammaLab 2.0 cumple los requisitos funcionales y las expectativas del usuario, validando su desempeño en condiciones reales de uso y determinando si el sistema está listo para emplearse en el trabajo de análisis del laboratorio.

\subsection{Implementación}
El usuario ejecuta, con señales del laboratorio y con los parámetros de su flujo de trabajo, la funcionalidad asociada a cada uno de los 45 requisitos funcionales del SRS. Por cada requisito se diligencia un cuadro con la retroalimentación del usuario, el comportamiento observado en la versión 2.0 y el resultado de la prueba, que puede ser \textbf{aceptado}, \textbf{aceptado parcialmente} o \textbf{no aceptado}. El campo \textit{Antes} describe, cuando se conoce, el comportamiento de GammaLab 1.0 según su TD \cite{td2025} y la matriz de requisitos, para hacer visible el cambio entre versiones. Las pruebas se realizan en el equipo de la Tabla~\ref{tab:td_ut_entorno}.

\subsection{Cuadros de aceptación}

\begin{table}[H]
\centering
\small
\begin{tabularx}{\textwidth}{|>{\raggedright\arraybackslash}p{2.7cm}|X|>{\raggedright\arraybackslash}p{2.6cm}|}
\hline
\textbf{ID del requisito} & \textbf{Descripción breve} & \textbf{Rol del usuario} \\
\hline
FR-IO-01 (R08) & Persistir la señal cargada & Investigador \\
\hline
\multicolumn{3}{|>{\raggedright\arraybackslash}p{\dimexpr\textwidth-2\tabcolsep-2\arrayrulewidth}|}{\centering\textbf{Descripción}} \\
\hline
\multicolumn{3}{|>{\raggedright\arraybackslash}p{\dimexpr\textwidth-2\tabcolsep-2\arrayrulewidth}|}{El sistema debe persistir la señal cargada} \\
\hline
\multicolumn{3}{|>{\raggedright\arraybackslash}p{\dimexpr\textwidth-2\tabcolsep-2\arrayrulewidth}|}{\centering\textbf{Criterios de aceptación}} \\
\hline
\multicolumn{3}{|>{\raggedright\arraybackslash}p{\dimexpr\textwidth-2\tabcolsep-2\arrayrulewidth}|}{Al ejecutar la funcionalidad con señales reales del laboratorio, el usuario comprueba que el sistema persiste la señal cargada, sin errores y con un resultado coherente con lo que muestra la interfaz.} \\
\hline
\multicolumn{3}{|>{\raggedright\arraybackslash}p{\dimexpr\textwidth-2\tabcolsep-2\arrayrulewidth}|}{\centering\textbf{Antes (GammaLab 1.0)}} \\
\hline
\multicolumn{3}{|>{\raggedright\arraybackslash}p{\dimexpr\textwidth-2\tabcolsep-2\arrayrulewidth}|}{La versión 1.0 dejaba de mostrar la señal abierta al cambiar entre funcionalidades.} \\
\hline
\multicolumn{3}{|>{\raggedright\arraybackslash}p{\dimexpr\textwidth-2\tabcolsep-2\arrayrulewidth}|}{\centering\textbf{Después (GammaLab 2.0, observado durante la prueba)}} \\
\hline
\multicolumn{3}{|>{\raggedright\arraybackslash}p{\dimexpr\textwidth-2\tabcolsep-2\arrayrulewidth}|}{\rule{0pt}{2.2em}} \\
\hline
\multicolumn{3}{|>{\raggedright\arraybackslash}p{\dimexpr\textwidth-2\tabcolsep-2\arrayrulewidth}|}{\centering\textbf{Retroalimentación del usuario}} \\
\hline
\multicolumn{3}{|>{\raggedright\arraybackslash}p{\dimexpr\textwidth-2\tabcolsep-2\arrayrulewidth}|}{\rule{0pt}{2.8em}} \\
\hline
\multicolumn{3}{|>{\raggedright\arraybackslash}p{\dimexpr\textwidth-2\tabcolsep-2\arrayrulewidth}|}{\centering\textbf{Resultado de la prueba}} \\
\hline
\multicolumn{3}{|>{\raggedright\arraybackslash}p{\dimexpr\textwidth-2\tabcolsep-2\arrayrulewidth}|}{\casilla{} Aceptado \qquad \casilla{} Aceptado parcialmente \qquad \casilla{} No aceptado} \\
\hline
\multicolumn{3}{|>{\raggedright\arraybackslash}p{\dimexpr\textwidth-2\tabcolsep-2\arrayrulewidth}|}{\centering\textbf{Observaciones}} \\
\hline
\multicolumn{3}{|>{\raggedright\arraybackslash}p{\dimexpr\textwidth-2\tabcolsep-2\arrayrulewidth}|}{\rule{0pt}{2.2em}} \\
\hline
\end{tabularx}
\caption{Prueba de aceptación del requisito FR-IO-01 (R08)}
\label{tab:uat_fr-io-01}
\end{table}
\begin{table}[H]
\centering
\small
\begin{tabularx}{\textwidth}{|>{\raggedright\arraybackslash}p{2.7cm}|X|>{\raggedright\arraybackslash}p{2.6cm}|}
\hline
\textbf{ID del requisito} & \textbf{Descripción breve} & \textbf{Rol del usuario} \\
\hline
FR-IO-02 (R80) & Cargar señales ABF, EDF y MAT & Investigador \\
\hline
\multicolumn{3}{|>{\raggedright\arraybackslash}p{\dimexpr\textwidth-2\tabcolsep-2\arrayrulewidth}|}{\centering\textbf{Descripción}} \\
\hline
\multicolumn{3}{|>{\raggedright\arraybackslash}p{\dimexpr\textwidth-2\tabcolsep-2\arrayrulewidth}|}{El sistema debe cargar señales en los formatos ABF, EDF y MAT.} \\
\hline
\multicolumn{3}{|>{\raggedright\arraybackslash}p{\dimexpr\textwidth-2\tabcolsep-2\arrayrulewidth}|}{\centering\textbf{Criterios de aceptación}} \\
\hline
\multicolumn{3}{|>{\raggedright\arraybackslash}p{\dimexpr\textwidth-2\tabcolsep-2\arrayrulewidth}|}{Al ejecutar la funcionalidad con señales reales del laboratorio, el usuario comprueba que el sistema carga señales en los formatos ABF, EDF y MAT, sin errores y con un resultado coherente con lo que muestra la interfaz.} \\
\hline
\multicolumn{3}{|>{\raggedright\arraybackslash}p{\dimexpr\textwidth-2\tabcolsep-2\arrayrulewidth}|}{\centering\textbf{Antes (GammaLab 1.0)}} \\
\hline
\multicolumn{3}{|>{\raggedright\arraybackslash}p{\dimexpr\textwidth-2\tabcolsep-2\arrayrulewidth}|}{La versión 1.0 cargaba archivos ABF y EDF; la lectura de archivos MAT no estaba disponible.} \\
\hline
\multicolumn{3}{|>{\raggedright\arraybackslash}p{\dimexpr\textwidth-2\tabcolsep-2\arrayrulewidth}|}{\centering\textbf{Después (GammaLab 2.0, observado durante la prueba)}} \\
\hline
\multicolumn{3}{|>{\raggedright\arraybackslash}p{\dimexpr\textwidth-2\tabcolsep-2\arrayrulewidth}|}{\rule{0pt}{2.2em}} \\
\hline
\multicolumn{3}{|>{\raggedright\arraybackslash}p{\dimexpr\textwidth-2\tabcolsep-2\arrayrulewidth}|}{\centering\textbf{Retroalimentación del usuario}} \\
\hline
\multicolumn{3}{|>{\raggedright\arraybackslash}p{\dimexpr\textwidth-2\tabcolsep-2\arrayrulewidth}|}{\rule{0pt}{2.8em}} \\
\hline
\multicolumn{3}{|>{\raggedright\arraybackslash}p{\dimexpr\textwidth-2\tabcolsep-2\arrayrulewidth}|}{\centering\textbf{Resultado de la prueba}} \\
\hline
\multicolumn{3}{|>{\raggedright\arraybackslash}p{\dimexpr\textwidth-2\tabcolsep-2\arrayrulewidth}|}{\casilla{} Aceptado \qquad \casilla{} Aceptado parcialmente \qquad \casilla{} No aceptado} \\
\hline
\multicolumn{3}{|>{\raggedright\arraybackslash}p{\dimexpr\textwidth-2\tabcolsep-2\arrayrulewidth}|}{\centering\textbf{Observaciones}} \\
\hline
\multicolumn{3}{|>{\raggedright\arraybackslash}p{\dimexpr\textwidth-2\tabcolsep-2\arrayrulewidth}|}{\rule{0pt}{2.2em}} \\
\hline
\end{tabularx}
\caption{Prueba de aceptación del requisito FR-IO-02 (R80)}
\label{tab:uat_fr-io-02}
\end{table}
\begin{table}[H]
\centering
\small
\begin{tabularx}{\textwidth}{|>{\raggedright\arraybackslash}p{2.7cm}|X|>{\raggedright\arraybackslash}p{2.6cm}|}
\hline
\textbf{ID del requisito} & \textbf{Descripción breve} & \textbf{Rol del usuario} \\
\hline
FR-IO-03 (R07) & Guardar un proyecto & Investigador \\
\hline
\multicolumn{3}{|>{\raggedright\arraybackslash}p{\dimexpr\textwidth-2\tabcolsep-2\arrayrulewidth}|}{\centering\textbf{Descripción}} \\
\hline
\multicolumn{3}{|>{\raggedright\arraybackslash}p{\dimexpr\textwidth-2\tabcolsep-2\arrayrulewidth}|}{El sistema debe poder guardar un proyecto en un archivo con la extensión .glab que registre el nombre y la ruta del archivo de señal de origen junto con el trabajo realizado sobre ella, sin incluir los datos de la señal dentro de ese archivo.} \\
\hline
\multicolumn{3}{|>{\raggedright\arraybackslash}p{\dimexpr\textwidth-2\tabcolsep-2\arrayrulewidth}|}{\centering\textbf{Criterios de aceptación}} \\
\hline
\multicolumn{3}{|>{\raggedright\arraybackslash}p{\dimexpr\textwidth-2\tabcolsep-2\arrayrulewidth}|}{Al ejecutar la funcionalidad con señales reales del laboratorio, el usuario comprueba que el sistema puede guardar un proyecto en un archivo con la extensión .glab que registre el nombre y la ruta del archivo de señal de origen junto con el trabajo realizado sobre ella, sin incluir los datos de la señal dentro de ese archivo, sin errores y con un resultado coherente con lo que muestra la interfaz.} \\
\hline
\multicolumn{3}{|>{\raggedright\arraybackslash}p{\dimexpr\textwidth-2\tabcolsep-2\arrayrulewidth}|}{\centering\textbf{Antes (GammaLab 1.0)}} \\
\hline
\multicolumn{3}{|>{\raggedright\arraybackslash}p{\dimexpr\textwidth-2\tabcolsep-2\arrayrulewidth}|}{La versión 1.0 no tenía el concepto de proyecto: el avance se perdía al cerrar la aplicación.} \\
\hline
\multicolumn{3}{|>{\raggedright\arraybackslash}p{\dimexpr\textwidth-2\tabcolsep-2\arrayrulewidth}|}{\centering\textbf{Después (GammaLab 2.0, observado durante la prueba)}} \\
\hline
\multicolumn{3}{|>{\raggedright\arraybackslash}p{\dimexpr\textwidth-2\tabcolsep-2\arrayrulewidth}|}{\rule{0pt}{2.2em}} \\
\hline
\multicolumn{3}{|>{\raggedright\arraybackslash}p{\dimexpr\textwidth-2\tabcolsep-2\arrayrulewidth}|}{\centering\textbf{Retroalimentación del usuario}} \\
\hline
\multicolumn{3}{|>{\raggedright\arraybackslash}p{\dimexpr\textwidth-2\tabcolsep-2\arrayrulewidth}|}{\rule{0pt}{2.8em}} \\
\hline
\multicolumn{3}{|>{\raggedright\arraybackslash}p{\dimexpr\textwidth-2\tabcolsep-2\arrayrulewidth}|}{\centering\textbf{Resultado de la prueba}} \\
\hline
\multicolumn{3}{|>{\raggedright\arraybackslash}p{\dimexpr\textwidth-2\tabcolsep-2\arrayrulewidth}|}{\casilla{} Aceptado \qquad \casilla{} Aceptado parcialmente \qquad \casilla{} No aceptado} \\
\hline
\multicolumn{3}{|>{\raggedright\arraybackslash}p{\dimexpr\textwidth-2\tabcolsep-2\arrayrulewidth}|}{\centering\textbf{Observaciones}} \\
\hline
\multicolumn{3}{|>{\raggedright\arraybackslash}p{\dimexpr\textwidth-2\tabcolsep-2\arrayrulewidth}|}{\rule{0pt}{2.2em}} \\
\hline
\end{tabularx}
\caption{Prueba de aceptación del requisito FR-IO-03 (R07)}
\label{tab:uat_fr-io-03}
\end{table}
\begin{table}[H]
\centering
\small
\begin{tabularx}{\textwidth}{|>{\raggedright\arraybackslash}p{2.7cm}|X|>{\raggedright\arraybackslash}p{2.6cm}|}
\hline
\textbf{ID del requisito} & \textbf{Descripción breve} & \textbf{Rol del usuario} \\
\hline
FR-IO-04 (R78) & Abrir un proyecto guardado & Investigador \\
\hline
\multicolumn{3}{|>{\raggedright\arraybackslash}p{\dimexpr\textwidth-2\tabcolsep-2\arrayrulewidth}|}{\centering\textbf{Descripción}} \\
\hline
\multicolumn{3}{|>{\raggedright\arraybackslash}p{\dimexpr\textwidth-2\tabcolsep-2\arrayrulewidth}|}{El sistema debe abrir un proyecto guardado, restaurando la señal referenciada y el trabajo registrado en él, e informar cuando el archivo de señal referenciado ya no esté disponible.} \\
\hline
\multicolumn{3}{|>{\raggedright\arraybackslash}p{\dimexpr\textwidth-2\tabcolsep-2\arrayrulewidth}|}{\centering\textbf{Criterios de aceptación}} \\
\hline
\multicolumn{3}{|>{\raggedright\arraybackslash}p{\dimexpr\textwidth-2\tabcolsep-2\arrayrulewidth}|}{Al ejecutar la funcionalidad con señales reales del laboratorio, el usuario comprueba que el sistema abre un proyecto guardado, restaurando la señal referenciada y el trabajo registrado en él, e informar cuando el archivo de señal referenciado ya no esté disponible, sin errores y con un resultado coherente con lo que muestra la interfaz.} \\
\hline
\multicolumn{3}{|>{\raggedright\arraybackslash}p{\dimexpr\textwidth-2\tabcolsep-2\arrayrulewidth}|}{\centering\textbf{Antes (GammaLab 1.0)}} \\
\hline
\multicolumn{3}{|>{\raggedright\arraybackslash}p{\dimexpr\textwidth-2\tabcolsep-2\arrayrulewidth}|}{La versión 1.0 no permitía retomar un trabajo guardado, porque no existía la gestión de proyectos.} \\
\hline
\multicolumn{3}{|>{\raggedright\arraybackslash}p{\dimexpr\textwidth-2\tabcolsep-2\arrayrulewidth}|}{\centering\textbf{Después (GammaLab 2.0, observado durante la prueba)}} \\
\hline
\multicolumn{3}{|>{\raggedright\arraybackslash}p{\dimexpr\textwidth-2\tabcolsep-2\arrayrulewidth}|}{\rule{0pt}{2.2em}} \\
\hline
\multicolumn{3}{|>{\raggedright\arraybackslash}p{\dimexpr\textwidth-2\tabcolsep-2\arrayrulewidth}|}{\centering\textbf{Retroalimentación del usuario}} \\
\hline
\multicolumn{3}{|>{\raggedright\arraybackslash}p{\dimexpr\textwidth-2\tabcolsep-2\arrayrulewidth}|}{\rule{0pt}{2.8em}} \\
\hline
\multicolumn{3}{|>{\raggedright\arraybackslash}p{\dimexpr\textwidth-2\tabcolsep-2\arrayrulewidth}|}{\centering\textbf{Resultado de la prueba}} \\
\hline
\multicolumn{3}{|>{\raggedright\arraybackslash}p{\dimexpr\textwidth-2\tabcolsep-2\arrayrulewidth}|}{\casilla{} Aceptado \qquad \casilla{} Aceptado parcialmente \qquad \casilla{} No aceptado} \\
\hline
\multicolumn{3}{|>{\raggedright\arraybackslash}p{\dimexpr\textwidth-2\tabcolsep-2\arrayrulewidth}|}{\centering\textbf{Observaciones}} \\
\hline
\multicolumn{3}{|>{\raggedright\arraybackslash}p{\dimexpr\textwidth-2\tabcolsep-2\arrayrulewidth}|}{\rule{0pt}{2.2em}} \\
\hline
\end{tabularx}
\caption{Prueba de aceptación del requisito FR-IO-04 (R78)}
\label{tab:uat_fr-io-04}
\end{table}
\begin{table}[H]
\centering
\small
\begin{tabularx}{\textwidth}{|>{\raggedright\arraybackslash}p{2.7cm}|X|>{\raggedright\arraybackslash}p{2.6cm}|}
\hline
\textbf{ID del requisito} & \textbf{Descripción breve} & \textbf{Rol del usuario} \\
\hline
FR-IO-05 (R05) & Exportar los resultados de los análisis & Investigador \\
\hline
\multicolumn{3}{|>{\raggedright\arraybackslash}p{\dimexpr\textwidth-2\tabcolsep-2\arrayrulewidth}|}{\centering\textbf{Descripción}} \\
\hline
\multicolumn{3}{|>{\raggedright\arraybackslash}p{\dimexpr\textwidth-2\tabcolsep-2\arrayrulewidth}|}{El sistema debe exportar los resultados de la sección de análisis en los formatos Excel, CSV, JSON, JPEG, JPG, PNG y TIFF.} \\
\hline
\multicolumn{3}{|>{\raggedright\arraybackslash}p{\dimexpr\textwidth-2\tabcolsep-2\arrayrulewidth}|}{\centering\textbf{Criterios de aceptación}} \\
\hline
\multicolumn{3}{|>{\raggedright\arraybackslash}p{\dimexpr\textwidth-2\tabcolsep-2\arrayrulewidth}|}{Al ejecutar la funcionalidad con señales reales del laboratorio, el usuario comprueba que el sistema exporta los resultados de la sección de análisis en los formatos Excel, CSV, JSON, JPEG, JPG, PNG y TIFF, sin errores y con un resultado coherente con lo que muestra la interfaz.} \\
\hline
\multicolumn{3}{|>{\raggedright\arraybackslash}p{\dimexpr\textwidth-2\tabcolsep-2\arrayrulewidth}|}{\centering\textbf{Antes (GammaLab 1.0)}} \\
\hline
\multicolumn{3}{|>{\raggedright\arraybackslash}p{\dimexpr\textwidth-2\tabcolsep-2\arrayrulewidth}|}{\rule{0pt}{2.2em}} \\
\hline
\multicolumn{3}{|>{\raggedright\arraybackslash}p{\dimexpr\textwidth-2\tabcolsep-2\arrayrulewidth}|}{\centering\textbf{Después (GammaLab 2.0, observado durante la prueba)}} \\
\hline
\multicolumn{3}{|>{\raggedright\arraybackslash}p{\dimexpr\textwidth-2\tabcolsep-2\arrayrulewidth}|}{\rule{0pt}{2.2em}} \\
\hline
\multicolumn{3}{|>{\raggedright\arraybackslash}p{\dimexpr\textwidth-2\tabcolsep-2\arrayrulewidth}|}{\centering\textbf{Retroalimentación del usuario}} \\
\hline
\multicolumn{3}{|>{\raggedright\arraybackslash}p{\dimexpr\textwidth-2\tabcolsep-2\arrayrulewidth}|}{\rule{0pt}{2.8em}} \\
\hline
\multicolumn{3}{|>{\raggedright\arraybackslash}p{\dimexpr\textwidth-2\tabcolsep-2\arrayrulewidth}|}{\centering\textbf{Resultado de la prueba}} \\
\hline
\multicolumn{3}{|>{\raggedright\arraybackslash}p{\dimexpr\textwidth-2\tabcolsep-2\arrayrulewidth}|}{\casilla{} Aceptado \qquad \casilla{} Aceptado parcialmente \qquad \casilla{} No aceptado} \\
\hline
\multicolumn{3}{|>{\raggedright\arraybackslash}p{\dimexpr\textwidth-2\tabcolsep-2\arrayrulewidth}|}{\centering\textbf{Observaciones}} \\
\hline
\multicolumn{3}{|>{\raggedright\arraybackslash}p{\dimexpr\textwidth-2\tabcolsep-2\arrayrulewidth}|}{\rule{0pt}{2.2em}} \\
\hline
\end{tabularx}
\caption{Prueba de aceptación del requisito FR-IO-05 (R05)}
\label{tab:uat_fr-io-05}
\end{table}
\begin{table}[H]
\centering
\small
\begin{tabularx}{\textwidth}{|>{\raggedright\arraybackslash}p{2.7cm}|X|>{\raggedright\arraybackslash}p{2.6cm}|}
\hline
\textbf{ID del requisito} & \textbf{Descripción breve} & \textbf{Rol del usuario} \\
\hline
FR-IO-06 (R09) & Una sola columna X en el CSV & Investigador \\
\hline
\multicolumn{3}{|>{\raggedright\arraybackslash}p{\dimexpr\textwidth-2\tabcolsep-2\arrayrulewidth}|}{\centering\textbf{Descripción}} \\
\hline
\multicolumn{3}{|>{\raggedright\arraybackslash}p{\dimexpr\textwidth-2\tabcolsep-2\arrayrulewidth}|}{El sistema debe exportar a CSV una única columna de eje X cuando todas las series comparten el mismo eje.} \\
\hline
\multicolumn{3}{|>{\raggedright\arraybackslash}p{\dimexpr\textwidth-2\tabcolsep-2\arrayrulewidth}|}{\centering\textbf{Criterios de aceptación}} \\
\hline
\multicolumn{3}{|>{\raggedright\arraybackslash}p{\dimexpr\textwidth-2\tabcolsep-2\arrayrulewidth}|}{Al ejecutar la funcionalidad con señales reales del laboratorio, el usuario comprueba que el sistema exporta a CSV una única columna de eje X cuando todas las series comparten el mismo eje, sin errores y con un resultado coherente con lo que muestra la interfaz.} \\
\hline
\multicolumn{3}{|>{\raggedright\arraybackslash}p{\dimexpr\textwidth-2\tabcolsep-2\arrayrulewidth}|}{\centering\textbf{Antes (GammaLab 1.0)}} \\
\hline
\multicolumn{3}{|>{\raggedright\arraybackslash}p{\dimexpr\textwidth-2\tabcolsep-2\arrayrulewidth}|}{El archivo CSV repetía una columna de eje X por cada serie, aunque todas compartieran el mismo eje.} \\
\hline
\multicolumn{3}{|>{\raggedright\arraybackslash}p{\dimexpr\textwidth-2\tabcolsep-2\arrayrulewidth}|}{\centering\textbf{Después (GammaLab 2.0, observado durante la prueba)}} \\
\hline
\multicolumn{3}{|>{\raggedright\arraybackslash}p{\dimexpr\textwidth-2\tabcolsep-2\arrayrulewidth}|}{\rule{0pt}{2.2em}} \\
\hline
\multicolumn{3}{|>{\raggedright\arraybackslash}p{\dimexpr\textwidth-2\tabcolsep-2\arrayrulewidth}|}{\centering\textbf{Retroalimentación del usuario}} \\
\hline
\multicolumn{3}{|>{\raggedright\arraybackslash}p{\dimexpr\textwidth-2\tabcolsep-2\arrayrulewidth}|}{\rule{0pt}{2.8em}} \\
\hline
\multicolumn{3}{|>{\raggedright\arraybackslash}p{\dimexpr\textwidth-2\tabcolsep-2\arrayrulewidth}|}{\centering\textbf{Resultado de la prueba}} \\
\hline
\multicolumn{3}{|>{\raggedright\arraybackslash}p{\dimexpr\textwidth-2\tabcolsep-2\arrayrulewidth}|}{\casilla{} Aceptado \qquad \casilla{} Aceptado parcialmente \qquad \casilla{} No aceptado} \\
\hline
\multicolumn{3}{|>{\raggedright\arraybackslash}p{\dimexpr\textwidth-2\tabcolsep-2\arrayrulewidth}|}{\centering\textbf{Observaciones}} \\
\hline
\multicolumn{3}{|>{\raggedright\arraybackslash}p{\dimexpr\textwidth-2\tabcolsep-2\arrayrulewidth}|}{\rule{0pt}{2.2em}} \\
\hline
\end{tabularx}
\caption{Prueba de aceptación del requisito FR-IO-06 (R09)}
\label{tab:uat_fr-io-06}
\end{table}
\begin{table}[H]
\centering
\small
\begin{tabularx}{\textwidth}{|>{\raggedright\arraybackslash}p{2.7cm}|X|>{\raggedright\arraybackslash}p{2.6cm}|}
\hline
\textbf{ID del requisito} & \textbf{Descripción breve} & \textbf{Rol del usuario} \\
\hline
FR-IO-07 (R28) & Nombre de exportación descriptivo & Investigador \\
\hline
\multicolumn{3}{|>{\raggedright\arraybackslash}p{\dimexpr\textwidth-2\tabcolsep-2\arrayrulewidth}|}{\centering\textbf{Descripción}} \\
\hline
\multicolumn{3}{|>{\raggedright\arraybackslash}p{\dimexpr\textwidth-2\tabcolsep-2\arrayrulewidth}|}{El sistema debe asignar un nombre de exportación en el que se incluya el nombre de la señal} \\
\hline
\multicolumn{3}{|>{\raggedright\arraybackslash}p{\dimexpr\textwidth-2\tabcolsep-2\arrayrulewidth}|}{\centering\textbf{Criterios de aceptación}} \\
\hline
\multicolumn{3}{|>{\raggedright\arraybackslash}p{\dimexpr\textwidth-2\tabcolsep-2\arrayrulewidth}|}{Al ejecutar la funcionalidad con señales reales del laboratorio, el usuario comprueba que el sistema asigna un nombre de exportación en el que se incluya el nombre de la señal, sin errores y con un resultado coherente con lo que muestra la interfaz.} \\
\hline
\multicolumn{3}{|>{\raggedright\arraybackslash}p{\dimexpr\textwidth-2\tabcolsep-2\arrayrulewidth}|}{\centering\textbf{Antes (GammaLab 1.0)}} \\
\hline
\multicolumn{3}{|>{\raggedright\arraybackslash}p{\dimexpr\textwidth-2\tabcolsep-2\arrayrulewidth}|}{Los archivos exportados no incluían el nombre de la señal, lo que dificultaba identificarlos cuando se realizaban varias exportaciones.} \\
\hline
\multicolumn{3}{|>{\raggedright\arraybackslash}p{\dimexpr\textwidth-2\tabcolsep-2\arrayrulewidth}|}{\centering\textbf{Después (GammaLab 2.0, observado durante la prueba)}} \\
\hline
\multicolumn{3}{|>{\raggedright\arraybackslash}p{\dimexpr\textwidth-2\tabcolsep-2\arrayrulewidth}|}{\rule{0pt}{2.2em}} \\
\hline
\multicolumn{3}{|>{\raggedright\arraybackslash}p{\dimexpr\textwidth-2\tabcolsep-2\arrayrulewidth}|}{\centering\textbf{Retroalimentación del usuario}} \\
\hline
\multicolumn{3}{|>{\raggedright\arraybackslash}p{\dimexpr\textwidth-2\tabcolsep-2\arrayrulewidth}|}{\rule{0pt}{2.8em}} \\
\hline
\multicolumn{3}{|>{\raggedright\arraybackslash}p{\dimexpr\textwidth-2\tabcolsep-2\arrayrulewidth}|}{\centering\textbf{Resultado de la prueba}} \\
\hline
\multicolumn{3}{|>{\raggedright\arraybackslash}p{\dimexpr\textwidth-2\tabcolsep-2\arrayrulewidth}|}{\casilla{} Aceptado \qquad \casilla{} Aceptado parcialmente \qquad \casilla{} No aceptado} \\
\hline
\multicolumn{3}{|>{\raggedright\arraybackslash}p{\dimexpr\textwidth-2\tabcolsep-2\arrayrulewidth}|}{\centering\textbf{Observaciones}} \\
\hline
\multicolumn{3}{|>{\raggedright\arraybackslash}p{\dimexpr\textwidth-2\tabcolsep-2\arrayrulewidth}|}{\rule{0pt}{2.2em}} \\
\hline
\end{tabularx}
\caption{Prueba de aceptación del requisito FR-IO-07 (R28)}
\label{tab:uat_fr-io-07}
\end{table}
\begin{table}[H]
\centering
\small
\begin{tabularx}{\textwidth}{|>{\raggedright\arraybackslash}p{2.7cm}|X|>{\raggedright\arraybackslash}p{2.6cm}|}
\hline
\textbf{ID del requisito} & \textbf{Descripción breve} & \textbf{Rol del usuario} \\
\hline
FR-UI-01 (R01) & Cortar la señal en \textit{trials} & Investigador \\
\hline
\multicolumn{3}{|>{\raggedright\arraybackslash}p{\dimexpr\textwidth-2\tabcolsep-2\arrayrulewidth}|}{\centering\textbf{Descripción}} \\
\hline
\multicolumn{3}{|>{\raggedright\arraybackslash}p{\dimexpr\textwidth-2\tabcolsep-2\arrayrulewidth}|}{El sistema debe permitir realizar cortes de la señal} \\
\hline
\multicolumn{3}{|>{\raggedright\arraybackslash}p{\dimexpr\textwidth-2\tabcolsep-2\arrayrulewidth}|}{\centering\textbf{Criterios de aceptación}} \\
\hline
\multicolumn{3}{|>{\raggedright\arraybackslash}p{\dimexpr\textwidth-2\tabcolsep-2\arrayrulewidth}|}{Al ejecutar la funcionalidad con señales reales del laboratorio, el usuario comprueba que el sistema permite realizar cortes de la señal, sin errores y con un resultado coherente con lo que muestra la interfaz.} \\
\hline
\multicolumn{3}{|>{\raggedright\arraybackslash}p{\dimexpr\textwidth-2\tabcolsep-2\arrayrulewidth}|}{\centering\textbf{Antes (GammaLab 1.0)}} \\
\hline
\multicolumn{3}{|>{\raggedright\arraybackslash}p{\dimexpr\textwidth-2\tabcolsep-2\arrayrulewidth}|}{\rule{0pt}{2.2em}} \\
\hline
\multicolumn{3}{|>{\raggedright\arraybackslash}p{\dimexpr\textwidth-2\tabcolsep-2\arrayrulewidth}|}{\centering\textbf{Después (GammaLab 2.0, observado durante la prueba)}} \\
\hline
\multicolumn{3}{|>{\raggedright\arraybackslash}p{\dimexpr\textwidth-2\tabcolsep-2\arrayrulewidth}|}{\rule{0pt}{2.2em}} \\
\hline
\multicolumn{3}{|>{\raggedright\arraybackslash}p{\dimexpr\textwidth-2\tabcolsep-2\arrayrulewidth}|}{\centering\textbf{Retroalimentación del usuario}} \\
\hline
\multicolumn{3}{|>{\raggedright\arraybackslash}p{\dimexpr\textwidth-2\tabcolsep-2\arrayrulewidth}|}{\rule{0pt}{2.8em}} \\
\hline
\multicolumn{3}{|>{\raggedright\arraybackslash}p{\dimexpr\textwidth-2\tabcolsep-2\arrayrulewidth}|}{\centering\textbf{Resultado de la prueba}} \\
\hline
\multicolumn{3}{|>{\raggedright\arraybackslash}p{\dimexpr\textwidth-2\tabcolsep-2\arrayrulewidth}|}{\casilla{} Aceptado \qquad \casilla{} Aceptado parcialmente \qquad \casilla{} No aceptado} \\
\hline
\multicolumn{3}{|>{\raggedright\arraybackslash}p{\dimexpr\textwidth-2\tabcolsep-2\arrayrulewidth}|}{\centering\textbf{Observaciones}} \\
\hline
\multicolumn{3}{|>{\raggedright\arraybackslash}p{\dimexpr\textwidth-2\tabcolsep-2\arrayrulewidth}|}{\rule{0pt}{2.2em}} \\
\hline
\end{tabularx}
\caption{Prueba de aceptación del requisito FR-UI-01 (R01)}
\label{tab:uat_fr-ui-01}
\end{table}
\begin{table}[H]
\centering
\small
\begin{tabularx}{\textwidth}{|>{\raggedright\arraybackslash}p{2.7cm}|X|>{\raggedright\arraybackslash}p{2.6cm}|}
\hline
\textbf{ID del requisito} & \textbf{Descripción breve} & \textbf{Rol del usuario} \\
\hline
FR-UI-02 (R24) & Modificar un \textit{trial} ya creado & Investigador \\
\hline
\multicolumn{3}{|>{\raggedright\arraybackslash}p{\dimexpr\textwidth-2\tabcolsep-2\arrayrulewidth}|}{\centering\textbf{Descripción}} \\
\hline
\multicolumn{3}{|>{\raggedright\arraybackslash}p{\dimexpr\textwidth-2\tabcolsep-2\arrayrulewidth}|}{El sistema debe permitir modificar un conjunto de \textit{trials} ya creado, regenerándolo con otros parámetros, descartando \textit{trials} o aplicándole preprocesamiento, sin alterar el conjunto original.} \\
\hline
\multicolumn{3}{|>{\raggedright\arraybackslash}p{\dimexpr\textwidth-2\tabcolsep-2\arrayrulewidth}|}{\centering\textbf{Criterios de aceptación}} \\
\hline
\multicolumn{3}{|>{\raggedright\arraybackslash}p{\dimexpr\textwidth-2\tabcolsep-2\arrayrulewidth}|}{Al ejecutar la funcionalidad con señales reales del laboratorio, el usuario comprueba que el sistema permite modificar un conjunto de \textit{trials} ya creado, regenerándolo con otros parámetros, descartando \textit{trials} o aplicándole preprocesamiento, sin alterar el conjunto original, sin errores y con un resultado coherente con lo que muestra la interfaz.} \\
\hline
\multicolumn{3}{|>{\raggedright\arraybackslash}p{\dimexpr\textwidth-2\tabcolsep-2\arrayrulewidth}|}{\centering\textbf{Antes (GammaLab 1.0)}} \\
\hline
\multicolumn{3}{|>{\raggedright\arraybackslash}p{\dimexpr\textwidth-2\tabcolsep-2\arrayrulewidth}|}{\rule{0pt}{2.2em}} \\
\hline
\multicolumn{3}{|>{\raggedright\arraybackslash}p{\dimexpr\textwidth-2\tabcolsep-2\arrayrulewidth}|}{\centering\textbf{Después (GammaLab 2.0, observado durante la prueba)}} \\
\hline
\multicolumn{3}{|>{\raggedright\arraybackslash}p{\dimexpr\textwidth-2\tabcolsep-2\arrayrulewidth}|}{\rule{0pt}{2.2em}} \\
\hline
\multicolumn{3}{|>{\raggedright\arraybackslash}p{\dimexpr\textwidth-2\tabcolsep-2\arrayrulewidth}|}{\centering\textbf{Retroalimentación del usuario}} \\
\hline
\multicolumn{3}{|>{\raggedright\arraybackslash}p{\dimexpr\textwidth-2\tabcolsep-2\arrayrulewidth}|}{\rule{0pt}{2.8em}} \\
\hline
\multicolumn{3}{|>{\raggedright\arraybackslash}p{\dimexpr\textwidth-2\tabcolsep-2\arrayrulewidth}|}{\centering\textbf{Resultado de la prueba}} \\
\hline
\multicolumn{3}{|>{\raggedright\arraybackslash}p{\dimexpr\textwidth-2\tabcolsep-2\arrayrulewidth}|}{\casilla{} Aceptado \qquad \casilla{} Aceptado parcialmente \qquad \casilla{} No aceptado} \\
\hline
\multicolumn{3}{|>{\raggedright\arraybackslash}p{\dimexpr\textwidth-2\tabcolsep-2\arrayrulewidth}|}{\centering\textbf{Observaciones}} \\
\hline
\multicolumn{3}{|>{\raggedright\arraybackslash}p{\dimexpr\textwidth-2\tabcolsep-2\arrayrulewidth}|}{\rule{0pt}{2.2em}} \\
\hline
\end{tabularx}
\caption{Prueba de aceptación del requisito FR-UI-02 (R24)}
\label{tab:uat_fr-ui-02}
\end{table}
\begin{table}[H]
\centering
\small
\begin{tabularx}{\textwidth}{|>{\raggedright\arraybackslash}p{2.7cm}|X|>{\raggedright\arraybackslash}p{2.6cm}|}
\hline
\textbf{ID del requisito} & \textbf{Descripción breve} & \textbf{Rol del usuario} \\
\hline
FR-UI-03 (R25) & Restablecer parámetros por defecto & Investigador \\
\hline
\multicolumn{3}{|>{\raggedright\arraybackslash}p{\dimexpr\textwidth-2\tabcolsep-2\arrayrulewidth}|}{\centering\textbf{Descripción}} \\
\hline
\multicolumn{3}{|>{\raggedright\arraybackslash}p{\dimexpr\textwidth-2\tabcolsep-2\arrayrulewidth}|}{El sistema debe restablecer a sus valores por defecto todos los parámetros de una funcionalidad mediante un botón Clear, presente en cada funcionalidad que se ejecute con un botón.} \\
\hline
\multicolumn{3}{|>{\raggedright\arraybackslash}p{\dimexpr\textwidth-2\tabcolsep-2\arrayrulewidth}|}{\centering\textbf{Criterios de aceptación}} \\
\hline
\multicolumn{3}{|>{\raggedright\arraybackslash}p{\dimexpr\textwidth-2\tabcolsep-2\arrayrulewidth}|}{Al ejecutar la funcionalidad con señales reales del laboratorio, el usuario comprueba que el sistema restablece a sus valores por defecto todos los parámetros de una funcionalidad mediante un botón Clear, presente en cada funcionalidad que se ejecute con un botón, sin errores y con un resultado coherente con lo que muestra la interfaz.} \\
\hline
\multicolumn{3}{|>{\raggedright\arraybackslash}p{\dimexpr\textwidth-2\tabcolsep-2\arrayrulewidth}|}{\centering\textbf{Antes (GammaLab 1.0)}} \\
\hline
\multicolumn{3}{|>{\raggedright\arraybackslash}p{\dimexpr\textwidth-2\tabcolsep-2\arrayrulewidth}|}{Para ensayar otra configuración había que corregir los parámetros campo por campo o salir del módulo y volver a entrar.} \\
\hline
\multicolumn{3}{|>{\raggedright\arraybackslash}p{\dimexpr\textwidth-2\tabcolsep-2\arrayrulewidth}|}{\centering\textbf{Después (GammaLab 2.0, observado durante la prueba)}} \\
\hline
\multicolumn{3}{|>{\raggedright\arraybackslash}p{\dimexpr\textwidth-2\tabcolsep-2\arrayrulewidth}|}{\rule{0pt}{2.2em}} \\
\hline
\multicolumn{3}{|>{\raggedright\arraybackslash}p{\dimexpr\textwidth-2\tabcolsep-2\arrayrulewidth}|}{\centering\textbf{Retroalimentación del usuario}} \\
\hline
\multicolumn{3}{|>{\raggedright\arraybackslash}p{\dimexpr\textwidth-2\tabcolsep-2\arrayrulewidth}|}{\rule{0pt}{2.8em}} \\
\hline
\multicolumn{3}{|>{\raggedright\arraybackslash}p{\dimexpr\textwidth-2\tabcolsep-2\arrayrulewidth}|}{\centering\textbf{Resultado de la prueba}} \\
\hline
\multicolumn{3}{|>{\raggedright\arraybackslash}p{\dimexpr\textwidth-2\tabcolsep-2\arrayrulewidth}|}{\casilla{} Aceptado \qquad \casilla{} Aceptado parcialmente \qquad \casilla{} No aceptado} \\
\hline
\multicolumn{3}{|>{\raggedright\arraybackslash}p{\dimexpr\textwidth-2\tabcolsep-2\arrayrulewidth}|}{\centering\textbf{Observaciones}} \\
\hline
\multicolumn{3}{|>{\raggedright\arraybackslash}p{\dimexpr\textwidth-2\tabcolsep-2\arrayrulewidth}|}{\rule{0pt}{2.2em}} \\
\hline
\end{tabularx}
\caption{Prueba de aceptación del requisito FR-UI-03 (R25)}
\label{tab:uat_fr-ui-03}
\end{table}
\begin{table}[H]
\centering
\small
\begin{tabularx}{\textwidth}{|>{\raggedright\arraybackslash}p{2.7cm}|X|>{\raggedright\arraybackslash}p{2.6cm}|}
\hline
\textbf{ID del requisito} & \textbf{Descripción breve} & \textbf{Rol del usuario} \\
\hline
FR-UI-04 (R26) & Valores X/Y bajo el cursor & Investigador \\
\hline
\multicolumn{3}{|>{\raggedright\arraybackslash}p{\dimexpr\textwidth-2\tabcolsep-2\arrayrulewidth}|}{\centering\textbf{Descripción}} \\
\hline
\multicolumn{3}{|>{\raggedright\arraybackslash}p{\dimexpr\textwidth-2\tabcolsep-2\arrayrulewidth}|}{El sistema debe mostrar los valores de los ejes X y Y correspondientes al punto de la señal sobre el que se posiciona el cursor.} \\
\hline
\multicolumn{3}{|>{\raggedright\arraybackslash}p{\dimexpr\textwidth-2\tabcolsep-2\arrayrulewidth}|}{\centering\textbf{Criterios de aceptación}} \\
\hline
\multicolumn{3}{|>{\raggedright\arraybackslash}p{\dimexpr\textwidth-2\tabcolsep-2\arrayrulewidth}|}{Al ejecutar la funcionalidad con señales reales del laboratorio, el usuario comprueba que el sistema muestra los valores de los ejes X y Y correspondientes al punto de la señal sobre el que se posiciona el cursor, sin errores y con un resultado coherente con lo que muestra la interfaz.} \\
\hline
\multicolumn{3}{|>{\raggedright\arraybackslash}p{\dimexpr\textwidth-2\tabcolsep-2\arrayrulewidth}|}{\centering\textbf{Antes (GammaLab 1.0)}} \\
\hline
\multicolumn{3}{|>{\raggedright\arraybackslash}p{\dimexpr\textwidth-2\tabcolsep-2\arrayrulewidth}|}{La gráfica mostraba la forma de la señal, pero no permitía leer valores concretos sin exportar los datos.} \\
\hline
\multicolumn{3}{|>{\raggedright\arraybackslash}p{\dimexpr\textwidth-2\tabcolsep-2\arrayrulewidth}|}{\centering\textbf{Después (GammaLab 2.0, observado durante la prueba)}} \\
\hline
\multicolumn{3}{|>{\raggedright\arraybackslash}p{\dimexpr\textwidth-2\tabcolsep-2\arrayrulewidth}|}{\rule{0pt}{2.2em}} \\
\hline
\multicolumn{3}{|>{\raggedright\arraybackslash}p{\dimexpr\textwidth-2\tabcolsep-2\arrayrulewidth}|}{\centering\textbf{Retroalimentación del usuario}} \\
\hline
\multicolumn{3}{|>{\raggedright\arraybackslash}p{\dimexpr\textwidth-2\tabcolsep-2\arrayrulewidth}|}{\rule{0pt}{2.8em}} \\
\hline
\multicolumn{3}{|>{\raggedright\arraybackslash}p{\dimexpr\textwidth-2\tabcolsep-2\arrayrulewidth}|}{\centering\textbf{Resultado de la prueba}} \\
\hline
\multicolumn{3}{|>{\raggedright\arraybackslash}p{\dimexpr\textwidth-2\tabcolsep-2\arrayrulewidth}|}{\casilla{} Aceptado \qquad \casilla{} Aceptado parcialmente \qquad \casilla{} No aceptado} \\
\hline
\multicolumn{3}{|>{\raggedright\arraybackslash}p{\dimexpr\textwidth-2\tabcolsep-2\arrayrulewidth}|}{\centering\textbf{Observaciones}} \\
\hline
\multicolumn{3}{|>{\raggedright\arraybackslash}p{\dimexpr\textwidth-2\tabcolsep-2\arrayrulewidth}|}{\rule{0pt}{2.2em}} \\
\hline
\end{tabularx}
\caption{Prueba de aceptación del requisito FR-UI-04 (R26)}
\label{tab:uat_fr-ui-04}
\end{table}
\begin{table}[H]
\centering
\small
\begin{tabularx}{\textwidth}{|>{\raggedright\arraybackslash}p{2.7cm}|X|>{\raggedright\arraybackslash}p{2.6cm}|}
\hline
\textbf{ID del requisito} & \textbf{Descripción breve} & \textbf{Rol del usuario} \\
\hline
FR-UI-05 (R27) & Elegir los \textit{trials} de cada análisis & Investigador \\
\hline
\multicolumn{3}{|>{\raggedright\arraybackslash}p{\dimexpr\textwidth-2\tabcolsep-2\arrayrulewidth}|}{\centering\textbf{Descripción}} \\
\hline
\multicolumn{3}{|>{\raggedright\arraybackslash}p{\dimexpr\textwidth-2\tabcolsep-2\arrayrulewidth}|}{El sistema debe permitir escoger qué \textit{trials} participan en cada análisis y descartar los que no deban usarse, conservándolos en el conjunto.} \\
\hline
\multicolumn{3}{|>{\raggedright\arraybackslash}p{\dimexpr\textwidth-2\tabcolsep-2\arrayrulewidth}|}{\centering\textbf{Criterios de aceptación}} \\
\hline
\multicolumn{3}{|>{\raggedright\arraybackslash}p{\dimexpr\textwidth-2\tabcolsep-2\arrayrulewidth}|}{Al ejecutar la funcionalidad con señales reales del laboratorio, el usuario comprueba que el sistema permite escoger qué \textit{trials} participan en cada análisis y descartar los que no deban usarse, conservándolos en el conjunto, sin errores y con un resultado coherente con lo que muestra la interfaz.} \\
\hline
\multicolumn{3}{|>{\raggedright\arraybackslash}p{\dimexpr\textwidth-2\tabcolsep-2\arrayrulewidth}|}{\centering\textbf{Antes (GammaLab 1.0)}} \\
\hline
\multicolumn{3}{|>{\raggedright\arraybackslash}p{\dimexpr\textwidth-2\tabcolsep-2\arrayrulewidth}|}{\rule{0pt}{2.2em}} \\
\hline
\multicolumn{3}{|>{\raggedright\arraybackslash}p{\dimexpr\textwidth-2\tabcolsep-2\arrayrulewidth}|}{\centering\textbf{Después (GammaLab 2.0, observado durante la prueba)}} \\
\hline
\multicolumn{3}{|>{\raggedright\arraybackslash}p{\dimexpr\textwidth-2\tabcolsep-2\arrayrulewidth}|}{\rule{0pt}{2.2em}} \\
\hline
\multicolumn{3}{|>{\raggedright\arraybackslash}p{\dimexpr\textwidth-2\tabcolsep-2\arrayrulewidth}|}{\centering\textbf{Retroalimentación del usuario}} \\
\hline
\multicolumn{3}{|>{\raggedright\arraybackslash}p{\dimexpr\textwidth-2\tabcolsep-2\arrayrulewidth}|}{\rule{0pt}{2.8em}} \\
\hline
\multicolumn{3}{|>{\raggedright\arraybackslash}p{\dimexpr\textwidth-2\tabcolsep-2\arrayrulewidth}|}{\centering\textbf{Resultado de la prueba}} \\
\hline
\multicolumn{3}{|>{\raggedright\arraybackslash}p{\dimexpr\textwidth-2\tabcolsep-2\arrayrulewidth}|}{\casilla{} Aceptado \qquad \casilla{} Aceptado parcialmente \qquad \casilla{} No aceptado} \\
\hline
\multicolumn{3}{|>{\raggedright\arraybackslash}p{\dimexpr\textwidth-2\tabcolsep-2\arrayrulewidth}|}{\centering\textbf{Observaciones}} \\
\hline
\multicolumn{3}{|>{\raggedright\arraybackslash}p{\dimexpr\textwidth-2\tabcolsep-2\arrayrulewidth}|}{\rule{0pt}{2.2em}} \\
\hline
\end{tabularx}
\caption{Prueba de aceptación del requisito FR-UI-05 (R27)}
\label{tab:uat_fr-ui-05}
\end{table}
\begin{table}[H]
\centering
\small
\begin{tabularx}{\textwidth}{|>{\raggedright\arraybackslash}p{2.7cm}|X|>{\raggedright\arraybackslash}p{2.6cm}|}
\hline
\textbf{ID del requisito} & \textbf{Descripción breve} & \textbf{Rol del usuario} \\
\hline
FR-UI-06 (R29) & Ajustar la frecuencia de muestreo & Investigador \\
\hline
\multicolumn{3}{|>{\raggedright\arraybackslash}p{\dimexpr\textwidth-2\tabcolsep-2\arrayrulewidth}|}{\centering\textbf{Descripción}} \\
\hline
\multicolumn{3}{|>{\raggedright\arraybackslash}p{\dimexpr\textwidth-2\tabcolsep-2\arrayrulewidth}|}{El sistema debe ajustar el tiempo de muestreo} \\
\hline
\multicolumn{3}{|>{\raggedright\arraybackslash}p{\dimexpr\textwidth-2\tabcolsep-2\arrayrulewidth}|}{\centering\textbf{Criterios de aceptación}} \\
\hline
\multicolumn{3}{|>{\raggedright\arraybackslash}p{\dimexpr\textwidth-2\tabcolsep-2\arrayrulewidth}|}{Al ejecutar la funcionalidad con señales reales del laboratorio, el usuario comprueba que el sistema ajusta el tiempo de muestreo, sin errores y con un resultado coherente con lo que muestra la interfaz.} \\
\hline
\multicolumn{3}{|>{\raggedright\arraybackslash}p{\dimexpr\textwidth-2\tabcolsep-2\arrayrulewidth}|}{\centering\textbf{Antes (GammaLab 1.0)}} \\
\hline
\multicolumn{3}{|>{\raggedright\arraybackslash}p{\dimexpr\textwidth-2\tabcolsep-2\arrayrulewidth}|}{La frecuencia de muestreo era un valor predefinido que el usuario no podía corregir.} \\
\hline
\multicolumn{3}{|>{\raggedright\arraybackslash}p{\dimexpr\textwidth-2\tabcolsep-2\arrayrulewidth}|}{\centering\textbf{Después (GammaLab 2.0, observado durante la prueba)}} \\
\hline
\multicolumn{3}{|>{\raggedright\arraybackslash}p{\dimexpr\textwidth-2\tabcolsep-2\arrayrulewidth}|}{\rule{0pt}{2.2em}} \\
\hline
\multicolumn{3}{|>{\raggedright\arraybackslash}p{\dimexpr\textwidth-2\tabcolsep-2\arrayrulewidth}|}{\centering\textbf{Retroalimentación del usuario}} \\
\hline
\multicolumn{3}{|>{\raggedright\arraybackslash}p{\dimexpr\textwidth-2\tabcolsep-2\arrayrulewidth}|}{\rule{0pt}{2.8em}} \\
\hline
\multicolumn{3}{|>{\raggedright\arraybackslash}p{\dimexpr\textwidth-2\tabcolsep-2\arrayrulewidth}|}{\centering\textbf{Resultado de la prueba}} \\
\hline
\multicolumn{3}{|>{\raggedright\arraybackslash}p{\dimexpr\textwidth-2\tabcolsep-2\arrayrulewidth}|}{\casilla{} Aceptado \qquad \casilla{} Aceptado parcialmente \qquad \casilla{} No aceptado} \\
\hline
\multicolumn{3}{|>{\raggedright\arraybackslash}p{\dimexpr\textwidth-2\tabcolsep-2\arrayrulewidth}|}{\centering\textbf{Observaciones}} \\
\hline
\multicolumn{3}{|>{\raggedright\arraybackslash}p{\dimexpr\textwidth-2\tabcolsep-2\arrayrulewidth}|}{\rule{0pt}{2.2em}} \\
\hline
\end{tabularx}
\caption{Prueba de aceptación del requisito FR-UI-06 (R29)}
\label{tab:uat_fr-ui-06}
\end{table}
\begin{table}[H]
\centering
\small
\begin{tabularx}{\textwidth}{|>{\raggedright\arraybackslash}p{2.7cm}|X|>{\raggedright\arraybackslash}p{2.6cm}|}
\hline
\textbf{ID del requisito} & \textbf{Descripción breve} & \textbf{Rol del usuario} \\
\hline
FR-UI-07 (R30) & \textit{Zoom} en todas las direcciones & Investigador \\
\hline
\multicolumn{3}{|>{\raggedright\arraybackslash}p{\dimexpr\textwidth-2\tabcolsep-2\arrayrulewidth}|}{\centering\textbf{Descripción}} \\
\hline
\multicolumn{3}{|>{\raggedright\arraybackslash}p{\dimexpr\textwidth-2\tabcolsep-2\arrayrulewidth}|}{El sistema debe tener \textit{zoom} en todas las direcciones} \\
\hline
\multicolumn{3}{|>{\raggedright\arraybackslash}p{\dimexpr\textwidth-2\tabcolsep-2\arrayrulewidth}|}{\centering\textbf{Criterios de aceptación}} \\
\hline
\multicolumn{3}{|>{\raggedright\arraybackslash}p{\dimexpr\textwidth-2\tabcolsep-2\arrayrulewidth}|}{Al ejecutar la funcionalidad con señales reales del laboratorio, el usuario comprueba que el sistema tiene \textit{zoom} en todas las direcciones, sin errores y con un resultado coherente con lo que muestra la interfaz.} \\
\hline
\multicolumn{3}{|>{\raggedright\arraybackslash}p{\dimexpr\textwidth-2\tabcolsep-2\arrayrulewidth}|}{\centering\textbf{Antes (GammaLab 1.0)}} \\
\hline
\multicolumn{3}{|>{\raggedright\arraybackslash}p{\dimexpr\textwidth-2\tabcolsep-2\arrayrulewidth}|}{El \textit{zoom} estaba limitado a un solo eje.} \\
\hline
\multicolumn{3}{|>{\raggedright\arraybackslash}p{\dimexpr\textwidth-2\tabcolsep-2\arrayrulewidth}|}{\centering\textbf{Después (GammaLab 2.0, observado durante la prueba)}} \\
\hline
\multicolumn{3}{|>{\raggedright\arraybackslash}p{\dimexpr\textwidth-2\tabcolsep-2\arrayrulewidth}|}{\rule{0pt}{2.2em}} \\
\hline
\multicolumn{3}{|>{\raggedright\arraybackslash}p{\dimexpr\textwidth-2\tabcolsep-2\arrayrulewidth}|}{\centering\textbf{Retroalimentación del usuario}} \\
\hline
\multicolumn{3}{|>{\raggedright\arraybackslash}p{\dimexpr\textwidth-2\tabcolsep-2\arrayrulewidth}|}{\rule{0pt}{2.8em}} \\
\hline
\multicolumn{3}{|>{\raggedright\arraybackslash}p{\dimexpr\textwidth-2\tabcolsep-2\arrayrulewidth}|}{\centering\textbf{Resultado de la prueba}} \\
\hline
\multicolumn{3}{|>{\raggedright\arraybackslash}p{\dimexpr\textwidth-2\tabcolsep-2\arrayrulewidth}|}{\casilla{} Aceptado \qquad \casilla{} Aceptado parcialmente \qquad \casilla{} No aceptado} \\
\hline
\multicolumn{3}{|>{\raggedright\arraybackslash}p{\dimexpr\textwidth-2\tabcolsep-2\arrayrulewidth}|}{\centering\textbf{Observaciones}} \\
\hline
\multicolumn{3}{|>{\raggedright\arraybackslash}p{\dimexpr\textwidth-2\tabcolsep-2\arrayrulewidth}|}{\rule{0pt}{2.2em}} \\
\hline
\end{tabularx}
\caption{Prueba de aceptación del requisito FR-UI-07 (R30)}
\label{tab:uat_fr-ui-07}
\end{table}
\begin{table}[H]
\centering
\small
\begin{tabularx}{\textwidth}{|>{\raggedright\arraybackslash}p{2.7cm}|X|>{\raggedright\arraybackslash}p{2.6cm}|}
\hline
\textbf{ID del requisito} & \textbf{Descripción breve} & \textbf{Rol del usuario} \\
\hline
FR-UI-08 (R31) & Mostrar el umbral en \textit{mV} & Investigador \\
\hline
\multicolumn{3}{|>{\raggedright\arraybackslash}p{\dimexpr\textwidth-2\tabcolsep-2\arrayrulewidth}|}{\centering\textbf{Descripción}} \\
\hline
\multicolumn{3}{|>{\raggedright\arraybackslash}p{\dimexpr\textwidth-2\tabcolsep-2\arrayrulewidth}|}{El sistema debe mostrar el umbral en mV} \\
\hline
\multicolumn{3}{|>{\raggedright\arraybackslash}p{\dimexpr\textwidth-2\tabcolsep-2\arrayrulewidth}|}{\centering\textbf{Criterios de aceptación}} \\
\hline
\multicolumn{3}{|>{\raggedright\arraybackslash}p{\dimexpr\textwidth-2\tabcolsep-2\arrayrulewidth}|}{Al ejecutar la funcionalidad con señales reales del laboratorio, el usuario comprueba que el sistema muestra el umbral en mV, sin errores y con un resultado coherente con lo que muestra la interfaz.} \\
\hline
\multicolumn{3}{|>{\raggedright\arraybackslash}p{\dimexpr\textwidth-2\tabcolsep-2\arrayrulewidth}|}{\centering\textbf{Antes (GammaLab 1.0)}} \\
\hline
\multicolumn{3}{|>{\raggedright\arraybackslash}p{\dimexpr\textwidth-2\tabcolsep-2\arrayrulewidth}|}{\rule{0pt}{2.2em}} \\
\hline
\multicolumn{3}{|>{\raggedright\arraybackslash}p{\dimexpr\textwidth-2\tabcolsep-2\arrayrulewidth}|}{\centering\textbf{Después (GammaLab 2.0, observado durante la prueba)}} \\
\hline
\multicolumn{3}{|>{\raggedright\arraybackslash}p{\dimexpr\textwidth-2\tabcolsep-2\arrayrulewidth}|}{\rule{0pt}{2.2em}} \\
\hline
\multicolumn{3}{|>{\raggedright\arraybackslash}p{\dimexpr\textwidth-2\tabcolsep-2\arrayrulewidth}|}{\centering\textbf{Retroalimentación del usuario}} \\
\hline
\multicolumn{3}{|>{\raggedright\arraybackslash}p{\dimexpr\textwidth-2\tabcolsep-2\arrayrulewidth}|}{\rule{0pt}{2.8em}} \\
\hline
\multicolumn{3}{|>{\raggedright\arraybackslash}p{\dimexpr\textwidth-2\tabcolsep-2\arrayrulewidth}|}{\centering\textbf{Resultado de la prueba}} \\
\hline
\multicolumn{3}{|>{\raggedright\arraybackslash}p{\dimexpr\textwidth-2\tabcolsep-2\arrayrulewidth}|}{\casilla{} Aceptado \qquad \casilla{} Aceptado parcialmente \qquad \casilla{} No aceptado} \\
\hline
\multicolumn{3}{|>{\raggedright\arraybackslash}p{\dimexpr\textwidth-2\tabcolsep-2\arrayrulewidth}|}{\centering\textbf{Observaciones}} \\
\hline
\multicolumn{3}{|>{\raggedright\arraybackslash}p{\dimexpr\textwidth-2\tabcolsep-2\arrayrulewidth}|}{\rule{0pt}{2.2em}} \\
\hline
\end{tabularx}
\caption{Prueba de aceptación del requisito FR-UI-08 (R31)}
\label{tab:uat_fr-ui-08}
\end{table}
\begin{table}[H]
\centering
\small
\begin{tabularx}{\textwidth}{|>{\raggedright\arraybackslash}p{2.7cm}|X|>{\raggedright\arraybackslash}p{2.6cm}|}
\hline
\textbf{ID del requisito} & \textbf{Descripción breve} & \textbf{Rol del usuario} \\
\hline
FR-UI-09 (R32) & Mostrar \textit{Butterfly} explícitamente & Investigador \\
\hline
\multicolumn{3}{|>{\raggedright\arraybackslash}p{\dimexpr\textwidth-2\tabcolsep-2\arrayrulewidth}|}{\centering\textbf{Descripción}} \\
\hline
\multicolumn{3}{|>{\raggedright\arraybackslash}p{\dimexpr\textwidth-2\tabcolsep-2\arrayrulewidth}|}{El sistema debe mostrar \textit{Butterfly} explicitamente} \\
\hline
\multicolumn{3}{|>{\raggedright\arraybackslash}p{\dimexpr\textwidth-2\tabcolsep-2\arrayrulewidth}|}{\centering\textbf{Criterios de aceptación}} \\
\hline
\multicolumn{3}{|>{\raggedright\arraybackslash}p{\dimexpr\textwidth-2\tabcolsep-2\arrayrulewidth}|}{Al ejecutar la funcionalidad con señales reales del laboratorio, el usuario comprueba que el sistema muestra \textit{Butterfly} explicitamente, sin errores y con un resultado coherente con lo que muestra la interfaz.} \\
\hline
\multicolumn{3}{|>{\raggedright\arraybackslash}p{\dimexpr\textwidth-2\tabcolsep-2\arrayrulewidth}|}{\centering\textbf{Antes (GammaLab 1.0)}} \\
\hline
\multicolumn{3}{|>{\raggedright\arraybackslash}p{\dimexpr\textwidth-2\tabcolsep-2\arrayrulewidth}|}{El diagrama \textit{Butterfly} existía dentro del módulo ERP, pero no estaba señalado como tal en la interfaz.} \\
\hline
\multicolumn{3}{|>{\raggedright\arraybackslash}p{\dimexpr\textwidth-2\tabcolsep-2\arrayrulewidth}|}{\centering\textbf{Después (GammaLab 2.0, observado durante la prueba)}} \\
\hline
\multicolumn{3}{|>{\raggedright\arraybackslash}p{\dimexpr\textwidth-2\tabcolsep-2\arrayrulewidth}|}{\rule{0pt}{2.2em}} \\
\hline
\multicolumn{3}{|>{\raggedright\arraybackslash}p{\dimexpr\textwidth-2\tabcolsep-2\arrayrulewidth}|}{\centering\textbf{Retroalimentación del usuario}} \\
\hline
\multicolumn{3}{|>{\raggedright\arraybackslash}p{\dimexpr\textwidth-2\tabcolsep-2\arrayrulewidth}|}{\rule{0pt}{2.8em}} \\
\hline
\multicolumn{3}{|>{\raggedright\arraybackslash}p{\dimexpr\textwidth-2\tabcolsep-2\arrayrulewidth}|}{\centering\textbf{Resultado de la prueba}} \\
\hline
\multicolumn{3}{|>{\raggedright\arraybackslash}p{\dimexpr\textwidth-2\tabcolsep-2\arrayrulewidth}|}{\casilla{} Aceptado \qquad \casilla{} Aceptado parcialmente \qquad \casilla{} No aceptado} \\
\hline
\multicolumn{3}{|>{\raggedright\arraybackslash}p{\dimexpr\textwidth-2\tabcolsep-2\arrayrulewidth}|}{\centering\textbf{Observaciones}} \\
\hline
\multicolumn{3}{|>{\raggedright\arraybackslash}p{\dimexpr\textwidth-2\tabcolsep-2\arrayrulewidth}|}{\rule{0pt}{2.2em}} \\
\hline
\end{tabularx}
\caption{Prueba de aceptación del requisito FR-UI-09 (R32)}
\label{tab:uat_fr-ui-09}
\end{table}
\begin{table}[H]
\centering
\small
\begin{tabularx}{\textwidth}{|>{\raggedright\arraybackslash}p{2.7cm}|X|>{\raggedright\arraybackslash}p{2.6cm}|}
\hline
\textbf{ID del requisito} & \textbf{Descripción breve} & \textbf{Rol del usuario} \\
\hline
FR-UI-10 (R36) & Promediar todos los \textit{trials} & Investigador \\
\hline
\multicolumn{3}{|>{\raggedright\arraybackslash}p{\dimexpr\textwidth-2\tabcolsep-2\arrayrulewidth}|}{\centering\textbf{Descripción}} \\
\hline
\multicolumn{3}{|>{\raggedright\arraybackslash}p{\dimexpr\textwidth-2\tabcolsep-2\arrayrulewidth}|}{El sistema debe realizar el promedio en todos los \textit{trials}} \\
\hline
\multicolumn{3}{|>{\raggedright\arraybackslash}p{\dimexpr\textwidth-2\tabcolsep-2\arrayrulewidth}|}{\centering\textbf{Criterios de aceptación}} \\
\hline
\multicolumn{3}{|>{\raggedright\arraybackslash}p{\dimexpr\textwidth-2\tabcolsep-2\arrayrulewidth}|}{Al ejecutar la funcionalidad con señales reales del laboratorio, el usuario comprueba que el sistema realiza el promedio en todos los \textit{trials}, sin errores y con un resultado coherente con lo que muestra la interfaz.} \\
\hline
\multicolumn{3}{|>{\raggedright\arraybackslash}p{\dimexpr\textwidth-2\tabcolsep-2\arrayrulewidth}|}{\centering\textbf{Antes (GammaLab 1.0)}} \\
\hline
\multicolumn{3}{|>{\raggedright\arraybackslash}p{\dimexpr\textwidth-2\tabcolsep-2\arrayrulewidth}|}{\rule{0pt}{2.2em}} \\
\hline
\multicolumn{3}{|>{\raggedright\arraybackslash}p{\dimexpr\textwidth-2\tabcolsep-2\arrayrulewidth}|}{\centering\textbf{Después (GammaLab 2.0, observado durante la prueba)}} \\
\hline
\multicolumn{3}{|>{\raggedright\arraybackslash}p{\dimexpr\textwidth-2\tabcolsep-2\arrayrulewidth}|}{\rule{0pt}{2.2em}} \\
\hline
\multicolumn{3}{|>{\raggedright\arraybackslash}p{\dimexpr\textwidth-2\tabcolsep-2\arrayrulewidth}|}{\centering\textbf{Retroalimentación del usuario}} \\
\hline
\multicolumn{3}{|>{\raggedright\arraybackslash}p{\dimexpr\textwidth-2\tabcolsep-2\arrayrulewidth}|}{\rule{0pt}{2.8em}} \\
\hline
\multicolumn{3}{|>{\raggedright\arraybackslash}p{\dimexpr\textwidth-2\tabcolsep-2\arrayrulewidth}|}{\centering\textbf{Resultado de la prueba}} \\
\hline
\multicolumn{3}{|>{\raggedright\arraybackslash}p{\dimexpr\textwidth-2\tabcolsep-2\arrayrulewidth}|}{\casilla{} Aceptado \qquad \casilla{} Aceptado parcialmente \qquad \casilla{} No aceptado} \\
\hline
\multicolumn{3}{|>{\raggedright\arraybackslash}p{\dimexpr\textwidth-2\tabcolsep-2\arrayrulewidth}|}{\centering\textbf{Observaciones}} \\
\hline
\multicolumn{3}{|>{\raggedright\arraybackslash}p{\dimexpr\textwidth-2\tabcolsep-2\arrayrulewidth}|}{\rule{0pt}{2.2em}} \\
\hline
\end{tabularx}
\caption{Prueba de aceptación del requisito FR-UI-10 (R36)}
\label{tab:uat_fr-ui-10}
\end{table}
\begin{table}[H]
\centering
\small
\begin{tabularx}{\textwidth}{|>{\raggedright\arraybackslash}p{2.7cm}|X|>{\raggedright\arraybackslash}p{2.6cm}|}
\hline
\textbf{ID del requisito} & \textbf{Descripción breve} & \textbf{Rol del usuario} \\
\hline
FR-UI-11 (R37) & Indicar el \textit{trial} actual & Investigador \\
\hline
\multicolumn{3}{|>{\raggedright\arraybackslash}p{\dimexpr\textwidth-2\tabcolsep-2\arrayrulewidth}|}{\centering\textbf{Descripción}} \\
\hline
\multicolumn{3}{|>{\raggedright\arraybackslash}p{\dimexpr\textwidth-2\tabcolsep-2\arrayrulewidth}|}{El sistema debe indicar en todo momento en qué \textit{trial} se está trabajando.} \\
\hline
\multicolumn{3}{|>{\raggedright\arraybackslash}p{\dimexpr\textwidth-2\tabcolsep-2\arrayrulewidth}|}{\centering\textbf{Criterios de aceptación}} \\
\hline
\multicolumn{3}{|>{\raggedright\arraybackslash}p{\dimexpr\textwidth-2\tabcolsep-2\arrayrulewidth}|}{Al ejecutar la funcionalidad con señales reales del laboratorio, el usuario comprueba que el sistema indica en todo momento en qué \textit{trial} se está trabajando, sin errores y con un resultado coherente con lo que muestra la interfaz.} \\
\hline
\multicolumn{3}{|>{\raggedright\arraybackslash}p{\dimexpr\textwidth-2\tabcolsep-2\arrayrulewidth}|}{\centering\textbf{Antes (GammaLab 1.0)}} \\
\hline
\multicolumn{3}{|>{\raggedright\arraybackslash}p{\dimexpr\textwidth-2\tabcolsep-2\arrayrulewidth}|}{La interfaz no indicaba sobre qué \textit{trial} se estaba trabajando.} \\
\hline
\multicolumn{3}{|>{\raggedright\arraybackslash}p{\dimexpr\textwidth-2\tabcolsep-2\arrayrulewidth}|}{\centering\textbf{Después (GammaLab 2.0, observado durante la prueba)}} \\
\hline
\multicolumn{3}{|>{\raggedright\arraybackslash}p{\dimexpr\textwidth-2\tabcolsep-2\arrayrulewidth}|}{\rule{0pt}{2.2em}} \\
\hline
\multicolumn{3}{|>{\raggedright\arraybackslash}p{\dimexpr\textwidth-2\tabcolsep-2\arrayrulewidth}|}{\centering\textbf{Retroalimentación del usuario}} \\
\hline
\multicolumn{3}{|>{\raggedright\arraybackslash}p{\dimexpr\textwidth-2\tabcolsep-2\arrayrulewidth}|}{\rule{0pt}{2.8em}} \\
\hline
\multicolumn{3}{|>{\raggedright\arraybackslash}p{\dimexpr\textwidth-2\tabcolsep-2\arrayrulewidth}|}{\centering\textbf{Resultado de la prueba}} \\
\hline
\multicolumn{3}{|>{\raggedright\arraybackslash}p{\dimexpr\textwidth-2\tabcolsep-2\arrayrulewidth}|}{\casilla{} Aceptado \qquad \casilla{} Aceptado parcialmente \qquad \casilla{} No aceptado} \\
\hline
\multicolumn{3}{|>{\raggedright\arraybackslash}p{\dimexpr\textwidth-2\tabcolsep-2\arrayrulewidth}|}{\centering\textbf{Observaciones}} \\
\hline
\multicolumn{3}{|>{\raggedright\arraybackslash}p{\dimexpr\textwidth-2\tabcolsep-2\arrayrulewidth}|}{\rule{0pt}{2.2em}} \\
\hline
\end{tabularx}
\caption{Prueba de aceptación del requisito FR-UI-11 (R37)}
\label{tab:uat_fr-ui-11}
\end{table}
\begin{table}[H]
\centering
\small
\begin{tabularx}{\textwidth}{|>{\raggedright\arraybackslash}p{2.7cm}|X|>{\raggedright\arraybackslash}p{2.6cm}|}
\hline
\textbf{ID del requisito} & \textbf{Descripción breve} & \textbf{Rol del usuario} \\
\hline
FR-UI-12 (R38) & Indicar qué \textit{trials} participan & Investigador \\
\hline
\multicolumn{3}{|>{\raggedright\arraybackslash}p{\dimexpr\textwidth-2\tabcolsep-2\arrayrulewidth}|}{\centering\textbf{Descripción}} \\
\hline
\multicolumn{3}{|>{\raggedright\arraybackslash}p{\dimexpr\textwidth-2\tabcolsep-2\arrayrulewidth}|}{El sistema debe indicar qué \textit{trials} participan en cada operación.} \\
\hline
\multicolumn{3}{|>{\raggedright\arraybackslash}p{\dimexpr\textwidth-2\tabcolsep-2\arrayrulewidth}|}{\centering\textbf{Criterios de aceptación}} \\
\hline
\multicolumn{3}{|>{\raggedright\arraybackslash}p{\dimexpr\textwidth-2\tabcolsep-2\arrayrulewidth}|}{Al ejecutar la funcionalidad con señales reales del laboratorio, el usuario comprueba que el sistema indica qué \textit{trials} participan en cada operación, sin errores y con un resultado coherente con lo que muestra la interfaz.} \\
\hline
\multicolumn{3}{|>{\raggedright\arraybackslash}p{\dimexpr\textwidth-2\tabcolsep-2\arrayrulewidth}|}{\centering\textbf{Antes (GammaLab 1.0)}} \\
\hline
\multicolumn{3}{|>{\raggedright\arraybackslash}p{\dimexpr\textwidth-2\tabcolsep-2\arrayrulewidth}|}{\rule{0pt}{2.2em}} \\
\hline
\multicolumn{3}{|>{\raggedright\arraybackslash}p{\dimexpr\textwidth-2\tabcolsep-2\arrayrulewidth}|}{\centering\textbf{Después (GammaLab 2.0, observado durante la prueba)}} \\
\hline
\multicolumn{3}{|>{\raggedright\arraybackslash}p{\dimexpr\textwidth-2\tabcolsep-2\arrayrulewidth}|}{\rule{0pt}{2.2em}} \\
\hline
\multicolumn{3}{|>{\raggedright\arraybackslash}p{\dimexpr\textwidth-2\tabcolsep-2\arrayrulewidth}|}{\centering\textbf{Retroalimentación del usuario}} \\
\hline
\multicolumn{3}{|>{\raggedright\arraybackslash}p{\dimexpr\textwidth-2\tabcolsep-2\arrayrulewidth}|}{\rule{0pt}{2.8em}} \\
\hline
\multicolumn{3}{|>{\raggedright\arraybackslash}p{\dimexpr\textwidth-2\tabcolsep-2\arrayrulewidth}|}{\centering\textbf{Resultado de la prueba}} \\
\hline
\multicolumn{3}{|>{\raggedright\arraybackslash}p{\dimexpr\textwidth-2\tabcolsep-2\arrayrulewidth}|}{\casilla{} Aceptado \qquad \casilla{} Aceptado parcialmente \qquad \casilla{} No aceptado} \\
\hline
\multicolumn{3}{|>{\raggedright\arraybackslash}p{\dimexpr\textwidth-2\tabcolsep-2\arrayrulewidth}|}{\centering\textbf{Observaciones}} \\
\hline
\multicolumn{3}{|>{\raggedright\arraybackslash}p{\dimexpr\textwidth-2\tabcolsep-2\arrayrulewidth}|}{\rule{0pt}{2.2em}} \\
\hline
\end{tabularx}
\caption{Prueba de aceptación del requisito FR-UI-12 (R38)}
\label{tab:uat_fr-ui-12}
\end{table}
\begin{table}[H]
\centering
\small
\begin{tabularx}{\textwidth}{|>{\raggedright\arraybackslash}p{2.7cm}|X|>{\raggedright\arraybackslash}p{2.6cm}|}
\hline
\textbf{ID del requisito} & \textbf{Descripción breve} & \textbf{Rol del usuario} \\
\hline
FR-UI-13 (R40) & Remover un artefacto con medidas & Investigador \\
\hline
\multicolumn{3}{|>{\raggedright\arraybackslash}p{\dimexpr\textwidth-2\tabcolsep-2\arrayrulewidth}|}{\centering\textbf{Descripción}} \\
\hline
\multicolumn{3}{|>{\raggedright\arraybackslash}p{\dimexpr\textwidth-2\tabcolsep-2\arrayrulewidth}|}{El sistema debe remover un artefacto con medidas} \\
\hline
\multicolumn{3}{|>{\raggedright\arraybackslash}p{\dimexpr\textwidth-2\tabcolsep-2\arrayrulewidth}|}{\centering\textbf{Criterios de aceptación}} \\
\hline
\multicolumn{3}{|>{\raggedright\arraybackslash}p{\dimexpr\textwidth-2\tabcolsep-2\arrayrulewidth}|}{Al ejecutar la funcionalidad con señales reales del laboratorio, el usuario comprueba que el sistema remueve un artefacto con medidas, sin errores y con un resultado coherente con lo que muestra la interfaz.} \\
\hline
\multicolumn{3}{|>{\raggedright\arraybackslash}p{\dimexpr\textwidth-2\tabcolsep-2\arrayrulewidth}|}{\centering\textbf{Antes (GammaLab 1.0)}} \\
\hline
\multicolumn{3}{|>{\raggedright\arraybackslash}p{\dimexpr\textwidth-2\tabcolsep-2\arrayrulewidth}|}{\rule{0pt}{2.2em}} \\
\hline
\multicolumn{3}{|>{\raggedright\arraybackslash}p{\dimexpr\textwidth-2\tabcolsep-2\arrayrulewidth}|}{\centering\textbf{Después (GammaLab 2.0, observado durante la prueba)}} \\
\hline
\multicolumn{3}{|>{\raggedright\arraybackslash}p{\dimexpr\textwidth-2\tabcolsep-2\arrayrulewidth}|}{\rule{0pt}{2.2em}} \\
\hline
\multicolumn{3}{|>{\raggedright\arraybackslash}p{\dimexpr\textwidth-2\tabcolsep-2\arrayrulewidth}|}{\centering\textbf{Retroalimentación del usuario}} \\
\hline
\multicolumn{3}{|>{\raggedright\arraybackslash}p{\dimexpr\textwidth-2\tabcolsep-2\arrayrulewidth}|}{\rule{0pt}{2.8em}} \\
\hline
\multicolumn{3}{|>{\raggedright\arraybackslash}p{\dimexpr\textwidth-2\tabcolsep-2\arrayrulewidth}|}{\centering\textbf{Resultado de la prueba}} \\
\hline
\multicolumn{3}{|>{\raggedright\arraybackslash}p{\dimexpr\textwidth-2\tabcolsep-2\arrayrulewidth}|}{\casilla{} Aceptado \qquad \casilla{} Aceptado parcialmente \qquad \casilla{} No aceptado} \\
\hline
\multicolumn{3}{|>{\raggedright\arraybackslash}p{\dimexpr\textwidth-2\tabcolsep-2\arrayrulewidth}|}{\centering\textbf{Observaciones}} \\
\hline
\multicolumn{3}{|>{\raggedright\arraybackslash}p{\dimexpr\textwidth-2\tabcolsep-2\arrayrulewidth}|}{\rule{0pt}{2.2em}} \\
\hline
\end{tabularx}
\caption{Prueba de aceptación del requisito FR-UI-13 (R40)}
\label{tab:uat_fr-ui-13}
\end{table}
\begin{table}[H]
\centering
\small
\begin{tabularx}{\textwidth}{|>{\raggedright\arraybackslash}p{2.7cm}|X|>{\raggedright\arraybackslash}p{2.6cm}|}
\hline
\textbf{ID del requisito} & \textbf{Descripción breve} & \textbf{Rol del usuario} \\
\hline
FR-UI-14 (R41) & Promediar los \textit{trials} seleccionados & Investigador \\
\hline
\multicolumn{3}{|>{\raggedright\arraybackslash}p{\dimexpr\textwidth-2\tabcolsep-2\arrayrulewidth}|}{\centering\textbf{Descripción}} \\
\hline
\multicolumn{3}{|>{\raggedright\arraybackslash}p{\dimexpr\textwidth-2\tabcolsep-2\arrayrulewidth}|}{El sistema debe promediar los \textit{trials} específicos que el usuario seleccione.} \\
\hline
\multicolumn{3}{|>{\raggedright\arraybackslash}p{\dimexpr\textwidth-2\tabcolsep-2\arrayrulewidth}|}{\centering\textbf{Criterios de aceptación}} \\
\hline
\multicolumn{3}{|>{\raggedright\arraybackslash}p{\dimexpr\textwidth-2\tabcolsep-2\arrayrulewidth}|}{Al ejecutar la funcionalidad con señales reales del laboratorio, el usuario comprueba que el sistema promedia los \textit{trials} específicos que el usuario seleccione, sin errores y con un resultado coherente con lo que muestra la interfaz.} \\
\hline
\multicolumn{3}{|>{\raggedright\arraybackslash}p{\dimexpr\textwidth-2\tabcolsep-2\arrayrulewidth}|}{\centering\textbf{Antes (GammaLab 1.0)}} \\
\hline
\multicolumn{3}{|>{\raggedright\arraybackslash}p{\dimexpr\textwidth-2\tabcolsep-2\arrayrulewidth}|}{\rule{0pt}{2.2em}} \\
\hline
\multicolumn{3}{|>{\raggedright\arraybackslash}p{\dimexpr\textwidth-2\tabcolsep-2\arrayrulewidth}|}{\centering\textbf{Después (GammaLab 2.0, observado durante la prueba)}} \\
\hline
\multicolumn{3}{|>{\raggedright\arraybackslash}p{\dimexpr\textwidth-2\tabcolsep-2\arrayrulewidth}|}{\rule{0pt}{2.2em}} \\
\hline
\multicolumn{3}{|>{\raggedright\arraybackslash}p{\dimexpr\textwidth-2\tabcolsep-2\arrayrulewidth}|}{\centering\textbf{Retroalimentación del usuario}} \\
\hline
\multicolumn{3}{|>{\raggedright\arraybackslash}p{\dimexpr\textwidth-2\tabcolsep-2\arrayrulewidth}|}{\rule{0pt}{2.8em}} \\
\hline
\multicolumn{3}{|>{\raggedright\arraybackslash}p{\dimexpr\textwidth-2\tabcolsep-2\arrayrulewidth}|}{\centering\textbf{Resultado de la prueba}} \\
\hline
\multicolumn{3}{|>{\raggedright\arraybackslash}p{\dimexpr\textwidth-2\tabcolsep-2\arrayrulewidth}|}{\casilla{} Aceptado \qquad \casilla{} Aceptado parcialmente \qquad \casilla{} No aceptado} \\
\hline
\multicolumn{3}{|>{\raggedright\arraybackslash}p{\dimexpr\textwidth-2\tabcolsep-2\arrayrulewidth}|}{\centering\textbf{Observaciones}} \\
\hline
\multicolumn{3}{|>{\raggedright\arraybackslash}p{\dimexpr\textwidth-2\tabcolsep-2\arrayrulewidth}|}{\rule{0pt}{2.2em}} \\
\hline
\end{tabularx}
\caption{Prueba de aceptación del requisito FR-UI-14 (R41)}
\label{tab:uat_fr-ui-14}
\end{table}
\begin{table}[H]
\centering
\small
\begin{tabularx}{\textwidth}{|>{\raggedright\arraybackslash}p{2.7cm}|X|>{\raggedright\arraybackslash}p{2.6cm}|}
\hline
\textbf{ID del requisito} & \textbf{Descripción breve} & \textbf{Rol del usuario} \\
\hline
FR-UI-15 (R68) & \textit{Trials} identificados por archivo y canal & Investigador \\
\hline
\multicolumn{3}{|>{\raggedright\arraybackslash}p{\dimexpr\textwidth-2\tabcolsep-2\arrayrulewidth}|}{\centering\textbf{Descripción}} \\
\hline
\multicolumn{3}{|>{\raggedright\arraybackslash}p{\dimexpr\textwidth-2\tabcolsep-2\arrayrulewidth}|}{El sistema debe identificar cada conjunto de \textit{trials} por la pareja (archivo, canal) y devolver los \textit{trials} del canal solicitado, sin reemplazarlo por el último generado.} \\
\hline
\multicolumn{3}{|>{\raggedright\arraybackslash}p{\dimexpr\textwidth-2\tabcolsep-2\arrayrulewidth}|}{\centering\textbf{Criterios de aceptación}} \\
\hline
\multicolumn{3}{|>{\raggedright\arraybackslash}p{\dimexpr\textwidth-2\tabcolsep-2\arrayrulewidth}|}{Al ejecutar la funcionalidad con señales reales del laboratorio, el usuario comprueba que el sistema identifica cada conjunto de \textit{trials} por la pareja (archivo, canal) y devolver los \textit{trials} del canal solicitado, sin reemplazarlo por el último generado, sin errores y con un resultado coherente con lo que muestra la interfaz.} \\
\hline
\multicolumn{3}{|>{\raggedright\arraybackslash}p{\dimexpr\textwidth-2\tabcolsep-2\arrayrulewidth}|}{\centering\textbf{Antes (GammaLab 1.0)}} \\
\hline
\multicolumn{3}{|>{\raggedright\arraybackslash}p{\dimexpr\textwidth-2\tabcolsep-2\arrayrulewidth}|}{Al pedir los \textit{trials} de un canal, el sistema ignoraba el canal y devolvía siempre el último conjunto generado, lo que impedía analizar más de un canal.} \\
\hline
\multicolumn{3}{|>{\raggedright\arraybackslash}p{\dimexpr\textwidth-2\tabcolsep-2\arrayrulewidth}|}{\centering\textbf{Después (GammaLab 2.0, observado durante la prueba)}} \\
\hline
\multicolumn{3}{|>{\raggedright\arraybackslash}p{\dimexpr\textwidth-2\tabcolsep-2\arrayrulewidth}|}{\rule{0pt}{2.2em}} \\
\hline
\multicolumn{3}{|>{\raggedright\arraybackslash}p{\dimexpr\textwidth-2\tabcolsep-2\arrayrulewidth}|}{\centering\textbf{Retroalimentación del usuario}} \\
\hline
\multicolumn{3}{|>{\raggedright\arraybackslash}p{\dimexpr\textwidth-2\tabcolsep-2\arrayrulewidth}|}{\rule{0pt}{2.8em}} \\
\hline
\multicolumn{3}{|>{\raggedright\arraybackslash}p{\dimexpr\textwidth-2\tabcolsep-2\arrayrulewidth}|}{\centering\textbf{Resultado de la prueba}} \\
\hline
\multicolumn{3}{|>{\raggedright\arraybackslash}p{\dimexpr\textwidth-2\tabcolsep-2\arrayrulewidth}|}{\casilla{} Aceptado \qquad \casilla{} Aceptado parcialmente \qquad \casilla{} No aceptado} \\
\hline
\multicolumn{3}{|>{\raggedright\arraybackslash}p{\dimexpr\textwidth-2\tabcolsep-2\arrayrulewidth}|}{\centering\textbf{Observaciones}} \\
\hline
\multicolumn{3}{|>{\raggedright\arraybackslash}p{\dimexpr\textwidth-2\tabcolsep-2\arrayrulewidth}|}{\rule{0pt}{2.2em}} \\
\hline
\end{tabularx}
\caption{Prueba de aceptación del requisito FR-UI-15 (R68)}
\label{tab:uat_fr-ui-15}
\end{table}
\begin{table}[H]
\centering
\small
\begin{tabularx}{\textwidth}{|>{\raggedright\arraybackslash}p{2.7cm}|X|>{\raggedright\arraybackslash}p{2.6cm}|}
\hline
\textbf{ID del requisito} & \textbf{Descripción breve} & \textbf{Rol del usuario} \\
\hline
FR-WV-01 (R33) & Escala y rango de color del escalograma & Investigador \\
\hline
\multicolumn{3}{|>{\raggedright\arraybackslash}p{\dimexpr\textwidth-2\tabcolsep-2\arrayrulewidth}|}{\centering\textbf{Descripción}} \\
\hline
\multicolumn{3}{|>{\raggedright\arraybackslash}p{\dimexpr\textwidth-2\tabcolsep-2\arrayrulewidth}|}{El sistema debe permitir elegir el tipo de escala del escalograma de \textit{wavelet} y ajustar el rango de valores representado.} \\
\hline
\multicolumn{3}{|>{\raggedright\arraybackslash}p{\dimexpr\textwidth-2\tabcolsep-2\arrayrulewidth}|}{\centering\textbf{Criterios de aceptación}} \\
\hline
\multicolumn{3}{|>{\raggedright\arraybackslash}p{\dimexpr\textwidth-2\tabcolsep-2\arrayrulewidth}|}{Al ejecutar la funcionalidad con señales reales del laboratorio, el usuario comprueba que el sistema permite elegir el tipo de escala del escalograma de \textit{wavelet} y ajustar el rango de valores representado, sin errores y con un resultado coherente con lo que muestra la interfaz.} \\
\hline
\multicolumn{3}{|>{\raggedright\arraybackslash}p{\dimexpr\textwidth-2\tabcolsep-2\arrayrulewidth}|}{\centering\textbf{Antes (GammaLab 1.0)}} \\
\hline
\multicolumn{3}{|>{\raggedright\arraybackslash}p{\dimexpr\textwidth-2\tabcolsep-2\arrayrulewidth}|}{El escalograma tenía una escala fija, sin posibilidad de ajustar el rango de valores representado.} \\
\hline
\multicolumn{3}{|>{\raggedright\arraybackslash}p{\dimexpr\textwidth-2\tabcolsep-2\arrayrulewidth}|}{\centering\textbf{Después (GammaLab 2.0, observado durante la prueba)}} \\
\hline
\multicolumn{3}{|>{\raggedright\arraybackslash}p{\dimexpr\textwidth-2\tabcolsep-2\arrayrulewidth}|}{\rule{0pt}{2.2em}} \\
\hline
\multicolumn{3}{|>{\raggedright\arraybackslash}p{\dimexpr\textwidth-2\tabcolsep-2\arrayrulewidth}|}{\centering\textbf{Retroalimentación del usuario}} \\
\hline
\multicolumn{3}{|>{\raggedright\arraybackslash}p{\dimexpr\textwidth-2\tabcolsep-2\arrayrulewidth}|}{\rule{0pt}{2.8em}} \\
\hline
\multicolumn{3}{|>{\raggedright\arraybackslash}p{\dimexpr\textwidth-2\tabcolsep-2\arrayrulewidth}|}{\centering\textbf{Resultado de la prueba}} \\
\hline
\multicolumn{3}{|>{\raggedright\arraybackslash}p{\dimexpr\textwidth-2\tabcolsep-2\arrayrulewidth}|}{\casilla{} Aceptado \qquad \casilla{} Aceptado parcialmente \qquad \casilla{} No aceptado} \\
\hline
\multicolumn{3}{|>{\raggedright\arraybackslash}p{\dimexpr\textwidth-2\tabcolsep-2\arrayrulewidth}|}{\centering\textbf{Observaciones}} \\
\hline
\multicolumn{3}{|>{\raggedright\arraybackslash}p{\dimexpr\textwidth-2\tabcolsep-2\arrayrulewidth}|}{\rule{0pt}{2.2em}} \\
\hline
\end{tabularx}
\caption{Prueba de aceptación del requisito FR-WV-01 (R33)}
\label{tab:uat_fr-wv-01}
\end{table}
\begin{table}[H]
\centering
\small
\begin{tabularx}{\textwidth}{|>{\raggedright\arraybackslash}p{2.7cm}|X|>{\raggedright\arraybackslash}p{2.6cm}|}
\hline
\textbf{ID del requisito} & \textbf{Descripción breve} & \textbf{Rol del usuario} \\
\hline
FR-WV-02 (R34) & Normalizar el escalograma & Investigador \\
\hline
\multicolumn{3}{|>{\raggedright\arraybackslash}p{\dimexpr\textwidth-2\tabcolsep-2\arrayrulewidth}|}{\centering\textbf{Descripción}} \\
\hline
\multicolumn{3}{|>{\raggedright\arraybackslash}p{\dimexpr\textwidth-2\tabcolsep-2\arrayrulewidth}|}{El sistema debe normalizar el escalograma de \textit{wavelet} según el método que elija el usuario.} \\
\hline
\multicolumn{3}{|>{\raggedright\arraybackslash}p{\dimexpr\textwidth-2\tabcolsep-2\arrayrulewidth}|}{\centering\textbf{Criterios de aceptación}} \\
\hline
\multicolumn{3}{|>{\raggedright\arraybackslash}p{\dimexpr\textwidth-2\tabcolsep-2\arrayrulewidth}|}{Al ejecutar la funcionalidad con señales reales del laboratorio, el usuario comprueba que el sistema normaliza el escalograma de \textit{wavelet} según el método que elija el usuario, sin errores y con un resultado coherente con lo que muestra la interfaz.} \\
\hline
\multicolumn{3}{|>{\raggedright\arraybackslash}p{\dimexpr\textwidth-2\tabcolsep-2\arrayrulewidth}|}{\centering\textbf{Antes (GammaLab 1.0)}} \\
\hline
\multicolumn{3}{|>{\raggedright\arraybackslash}p{\dimexpr\textwidth-2\tabcolsep-2\arrayrulewidth}|}{La normalización del escalograma tenía una implementación incorrecta.} \\
\hline
\multicolumn{3}{|>{\raggedright\arraybackslash}p{\dimexpr\textwidth-2\tabcolsep-2\arrayrulewidth}|}{\centering\textbf{Después (GammaLab 2.0, observado durante la prueba)}} \\
\hline
\multicolumn{3}{|>{\raggedright\arraybackslash}p{\dimexpr\textwidth-2\tabcolsep-2\arrayrulewidth}|}{\rule{0pt}{2.2em}} \\
\hline
\multicolumn{3}{|>{\raggedright\arraybackslash}p{\dimexpr\textwidth-2\tabcolsep-2\arrayrulewidth}|}{\centering\textbf{Retroalimentación del usuario}} \\
\hline
\multicolumn{3}{|>{\raggedright\arraybackslash}p{\dimexpr\textwidth-2\tabcolsep-2\arrayrulewidth}|}{\rule{0pt}{2.8em}} \\
\hline
\multicolumn{3}{|>{\raggedright\arraybackslash}p{\dimexpr\textwidth-2\tabcolsep-2\arrayrulewidth}|}{\centering\textbf{Resultado de la prueba}} \\
\hline
\multicolumn{3}{|>{\raggedright\arraybackslash}p{\dimexpr\textwidth-2\tabcolsep-2\arrayrulewidth}|}{\casilla{} Aceptado \qquad \casilla{} Aceptado parcialmente \qquad \casilla{} No aceptado} \\
\hline
\multicolumn{3}{|>{\raggedright\arraybackslash}p{\dimexpr\textwidth-2\tabcolsep-2\arrayrulewidth}|}{\centering\textbf{Observaciones}} \\
\hline
\multicolumn{3}{|>{\raggedright\arraybackslash}p{\dimexpr\textwidth-2\tabcolsep-2\arrayrulewidth}|}{\rule{0pt}{2.2em}} \\
\hline
\end{tabularx}
\caption{Prueba de aceptación del requisito FR-WV-02 (R34)}
\label{tab:uat_fr-wv-02}
\end{table}
\begin{table}[H]
\centering
\small
\begin{tabularx}{\textwidth}{|>{\raggedright\arraybackslash}p{2.7cm}|X|>{\raggedright\arraybackslash}p{2.6cm}|}
\hline
\textbf{ID del requisito} & \textbf{Descripción breve} & \textbf{Rol del usuario} \\
\hline
FR-WV-03 (R35) & Exportar el escalograma en EMF & Investigador \\
\hline
\multicolumn{3}{|>{\raggedright\arraybackslash}p{\dimexpr\textwidth-2\tabcolsep-2\arrayrulewidth}|}{\centering\textbf{Descripción}} \\
\hline
\multicolumn{3}{|>{\raggedright\arraybackslash}p{\dimexpr\textwidth-2\tabcolsep-2\arrayrulewidth}|}{El sistema debe exportar \textit{wavelet} en EMF} \\
\hline
\multicolumn{3}{|>{\raggedright\arraybackslash}p{\dimexpr\textwidth-2\tabcolsep-2\arrayrulewidth}|}{\centering\textbf{Criterios de aceptación}} \\
\hline
\multicolumn{3}{|>{\raggedright\arraybackslash}p{\dimexpr\textwidth-2\tabcolsep-2\arrayrulewidth}|}{Al ejecutar la funcionalidad con señales reales del laboratorio, el usuario comprueba que el sistema exporta \textit{wavelet} en EMF, sin errores y con un resultado coherente con lo que muestra la interfaz.} \\
\hline
\multicolumn{3}{|>{\raggedright\arraybackslash}p{\dimexpr\textwidth-2\tabcolsep-2\arrayrulewidth}|}{\centering\textbf{Antes (GammaLab 1.0)}} \\
\hline
\multicolumn{3}{|>{\raggedright\arraybackslash}p{\dimexpr\textwidth-2\tabcolsep-2\arrayrulewidth}|}{La versión 1.0 no permitía exportar el escalograma en formato EMF.} \\
\hline
\multicolumn{3}{|>{\raggedright\arraybackslash}p{\dimexpr\textwidth-2\tabcolsep-2\arrayrulewidth}|}{\centering\textbf{Después (GammaLab 2.0, observado durante la prueba)}} \\
\hline
\multicolumn{3}{|>{\raggedright\arraybackslash}p{\dimexpr\textwidth-2\tabcolsep-2\arrayrulewidth}|}{\rule{0pt}{2.2em}} \\
\hline
\multicolumn{3}{|>{\raggedright\arraybackslash}p{\dimexpr\textwidth-2\tabcolsep-2\arrayrulewidth}|}{\centering\textbf{Retroalimentación del usuario}} \\
\hline
\multicolumn{3}{|>{\raggedright\arraybackslash}p{\dimexpr\textwidth-2\tabcolsep-2\arrayrulewidth}|}{\rule{0pt}{2.8em}} \\
\hline
\multicolumn{3}{|>{\raggedright\arraybackslash}p{\dimexpr\textwidth-2\tabcolsep-2\arrayrulewidth}|}{\centering\textbf{Resultado de la prueba}} \\
\hline
\multicolumn{3}{|>{\raggedright\arraybackslash}p{\dimexpr\textwidth-2\tabcolsep-2\arrayrulewidth}|}{\casilla{} Aceptado \qquad \casilla{} Aceptado parcialmente \qquad \casilla{} No aceptado} \\
\hline
\multicolumn{3}{|>{\raggedright\arraybackslash}p{\dimexpr\textwidth-2\tabcolsep-2\arrayrulewidth}|}{\centering\textbf{Observaciones}} \\
\hline
\multicolumn{3}{|>{\raggedright\arraybackslash}p{\dimexpr\textwidth-2\tabcolsep-2\arrayrulewidth}|}{\rule{0pt}{2.2em}} \\
\hline
\end{tabularx}
\caption{Prueba de aceptación del requisito FR-WV-03 (R35)}
\label{tab:uat_fr-wv-03}
\end{table}
\begin{table}[H]
\centering
\small
\begin{tabularx}{\textwidth}{|>{\raggedright\arraybackslash}p{2.7cm}|X|>{\raggedright\arraybackslash}p{2.6cm}|}
\hline
\textbf{ID del requisito} & \textbf{Descripción breve} & \textbf{Rol del usuario} \\
\hline
FR-PAC-01 (R13) & Ejecutar el análisis PAC & Investigador \\
\hline
\multicolumn{3}{|>{\raggedright\arraybackslash}p{\dimexpr\textwidth-2\tabcolsep-2\arrayrulewidth}|}{\centering\textbf{Descripción}} \\
\hline
\multicolumn{3}{|>{\raggedright\arraybackslash}p{\dimexpr\textwidth-2\tabcolsep-2\arrayrulewidth}|}{El sistema debe ejecutar el análisis de acoplamiento fase-amplitud (PAC) sobre una señal electrofisiológica cargada en el sistema.} \\
\hline
\multicolumn{3}{|>{\raggedright\arraybackslash}p{\dimexpr\textwidth-2\tabcolsep-2\arrayrulewidth}|}{\centering\textbf{Criterios de aceptación}} \\
\hline
\multicolumn{3}{|>{\raggedright\arraybackslash}p{\dimexpr\textwidth-2\tabcolsep-2\arrayrulewidth}|}{Al ejecutar la funcionalidad con señales reales del laboratorio, el usuario comprueba que el sistema ejecuta el análisis de acoplamiento fase-amplitud (PAC) sobre una señal electrofisiológica cargada en el sistema, sin errores y con un resultado coherente con lo que muestra la interfaz.} \\
\hline
\multicolumn{3}{|>{\raggedright\arraybackslash}p{\dimexpr\textwidth-2\tabcolsep-2\arrayrulewidth}|}{\centering\textbf{Antes (GammaLab 1.0)}} \\
\hline
\multicolumn{3}{|>{\raggedright\arraybackslash}p{\dimexpr\textwidth-2\tabcolsep-2\arrayrulewidth}|}{Funcionalidad nueva: la versión 1.0 no incluía el análisis de acoplamiento fase-amplitud (PAC).} \\
\hline
\multicolumn{3}{|>{\raggedright\arraybackslash}p{\dimexpr\textwidth-2\tabcolsep-2\arrayrulewidth}|}{\centering\textbf{Después (GammaLab 2.0, observado durante la prueba)}} \\
\hline
\multicolumn{3}{|>{\raggedright\arraybackslash}p{\dimexpr\textwidth-2\tabcolsep-2\arrayrulewidth}|}{\rule{0pt}{2.2em}} \\
\hline
\multicolumn{3}{|>{\raggedright\arraybackslash}p{\dimexpr\textwidth-2\tabcolsep-2\arrayrulewidth}|}{\centering\textbf{Retroalimentación del usuario}} \\
\hline
\multicolumn{3}{|>{\raggedright\arraybackslash}p{\dimexpr\textwidth-2\tabcolsep-2\arrayrulewidth}|}{\rule{0pt}{2.8em}} \\
\hline
\multicolumn{3}{|>{\raggedright\arraybackslash}p{\dimexpr\textwidth-2\tabcolsep-2\arrayrulewidth}|}{\centering\textbf{Resultado de la prueba}} \\
\hline
\multicolumn{3}{|>{\raggedright\arraybackslash}p{\dimexpr\textwidth-2\tabcolsep-2\arrayrulewidth}|}{\casilla{} Aceptado \qquad \casilla{} Aceptado parcialmente \qquad \casilla{} No aceptado} \\
\hline
\multicolumn{3}{|>{\raggedright\arraybackslash}p{\dimexpr\textwidth-2\tabcolsep-2\arrayrulewidth}|}{\centering\textbf{Observaciones}} \\
\hline
\multicolumn{3}{|>{\raggedright\arraybackslash}p{\dimexpr\textwidth-2\tabcolsep-2\arrayrulewidth}|}{\rule{0pt}{2.2em}} \\
\hline
\end{tabularx}
\caption{Prueba de aceptación del requisito FR-PAC-01 (R13)}
\label{tab:uat_fr-pac-01}
\end{table}
\begin{table}[H]
\centering
\small
\begin{tabularx}{\textwidth}{|>{\raggedright\arraybackslash}p{2.7cm}|X|>{\raggedright\arraybackslash}p{2.6cm}|}
\hline
\textbf{ID del requisito} & \textbf{Descripción breve} & \textbf{Rol del usuario} \\
\hline
FR-PAC-02 (R14) & Banda de la fase & Investigador \\
\hline
\multicolumn{3}{|>{\raggedright\arraybackslash}p{\dimexpr\textwidth-2\tabcolsep-2\arrayrulewidth}|}{\centering\textbf{Descripción}} \\
\hline
\multicolumn{3}{|>{\raggedright\arraybackslash}p{\dimexpr\textwidth-2\tabcolsep-2\arrayrulewidth}|}{El sistema debe seleccionar la banda de frecuencia correspondiente a la fase para el análisis PAC.} \\
\hline
\multicolumn{3}{|>{\raggedright\arraybackslash}p{\dimexpr\textwidth-2\tabcolsep-2\arrayrulewidth}|}{\centering\textbf{Criterios de aceptación}} \\
\hline
\multicolumn{3}{|>{\raggedright\arraybackslash}p{\dimexpr\textwidth-2\tabcolsep-2\arrayrulewidth}|}{Al ejecutar la funcionalidad con señales reales del laboratorio, el usuario comprueba que el sistema selecciona la banda de frecuencia correspondiente a la fase para el análisis PAC, sin errores y con un resultado coherente con lo que muestra la interfaz.} \\
\hline
\multicolumn{3}{|>{\raggedright\arraybackslash}p{\dimexpr\textwidth-2\tabcolsep-2\arrayrulewidth}|}{\centering\textbf{Antes (GammaLab 1.0)}} \\
\hline
\multicolumn{3}{|>{\raggedright\arraybackslash}p{\dimexpr\textwidth-2\tabcolsep-2\arrayrulewidth}|}{Funcionalidad nueva: la versión 1.0 no incluía el análisis de acoplamiento fase-amplitud (PAC).} \\
\hline
\multicolumn{3}{|>{\raggedright\arraybackslash}p{\dimexpr\textwidth-2\tabcolsep-2\arrayrulewidth}|}{\centering\textbf{Después (GammaLab 2.0, observado durante la prueba)}} \\
\hline
\multicolumn{3}{|>{\raggedright\arraybackslash}p{\dimexpr\textwidth-2\tabcolsep-2\arrayrulewidth}|}{\rule{0pt}{2.2em}} \\
\hline
\multicolumn{3}{|>{\raggedright\arraybackslash}p{\dimexpr\textwidth-2\tabcolsep-2\arrayrulewidth}|}{\centering\textbf{Retroalimentación del usuario}} \\
\hline
\multicolumn{3}{|>{\raggedright\arraybackslash}p{\dimexpr\textwidth-2\tabcolsep-2\arrayrulewidth}|}{\rule{0pt}{2.8em}} \\
\hline
\multicolumn{3}{|>{\raggedright\arraybackslash}p{\dimexpr\textwidth-2\tabcolsep-2\arrayrulewidth}|}{\centering\textbf{Resultado de la prueba}} \\
\hline
\multicolumn{3}{|>{\raggedright\arraybackslash}p{\dimexpr\textwidth-2\tabcolsep-2\arrayrulewidth}|}{\casilla{} Aceptado \qquad \casilla{} Aceptado parcialmente \qquad \casilla{} No aceptado} \\
\hline
\multicolumn{3}{|>{\raggedright\arraybackslash}p{\dimexpr\textwidth-2\tabcolsep-2\arrayrulewidth}|}{\centering\textbf{Observaciones}} \\
\hline
\multicolumn{3}{|>{\raggedright\arraybackslash}p{\dimexpr\textwidth-2\tabcolsep-2\arrayrulewidth}|}{\rule{0pt}{2.2em}} \\
\hline
\end{tabularx}
\caption{Prueba de aceptación del requisito FR-PAC-02 (R14)}
\label{tab:uat_fr-pac-02}
\end{table}
\begin{table}[H]
\centering
\small
\begin{tabularx}{\textwidth}{|>{\raggedright\arraybackslash}p{2.7cm}|X|>{\raggedright\arraybackslash}p{2.6cm}|}
\hline
\textbf{ID del requisito} & \textbf{Descripción breve} & \textbf{Rol del usuario} \\
\hline
FR-PAC-03 (R15) & Banda de la amplitud & Investigador \\
\hline
\multicolumn{3}{|>{\raggedright\arraybackslash}p{\dimexpr\textwidth-2\tabcolsep-2\arrayrulewidth}|}{\centering\textbf{Descripción}} \\
\hline
\multicolumn{3}{|>{\raggedright\arraybackslash}p{\dimexpr\textwidth-2\tabcolsep-2\arrayrulewidth}|}{El sistema debe seleccionar la banda de frecuencia correspondiente a la amplitud para el análisis PAC.} \\
\hline
\multicolumn{3}{|>{\raggedright\arraybackslash}p{\dimexpr\textwidth-2\tabcolsep-2\arrayrulewidth}|}{\centering\textbf{Criterios de aceptación}} \\
\hline
\multicolumn{3}{|>{\raggedright\arraybackslash}p{\dimexpr\textwidth-2\tabcolsep-2\arrayrulewidth}|}{Al ejecutar la funcionalidad con señales reales del laboratorio, el usuario comprueba que el sistema selecciona la banda de frecuencia correspondiente a la amplitud para el análisis PAC, sin errores y con un resultado coherente con lo que muestra la interfaz.} \\
\hline
\multicolumn{3}{|>{\raggedright\arraybackslash}p{\dimexpr\textwidth-2\tabcolsep-2\arrayrulewidth}|}{\centering\textbf{Antes (GammaLab 1.0)}} \\
\hline
\multicolumn{3}{|>{\raggedright\arraybackslash}p{\dimexpr\textwidth-2\tabcolsep-2\arrayrulewidth}|}{Funcionalidad nueva: la versión 1.0 no incluía el análisis de acoplamiento fase-amplitud (PAC).} \\
\hline
\multicolumn{3}{|>{\raggedright\arraybackslash}p{\dimexpr\textwidth-2\tabcolsep-2\arrayrulewidth}|}{\centering\textbf{Después (GammaLab 2.0, observado durante la prueba)}} \\
\hline
\multicolumn{3}{|>{\raggedright\arraybackslash}p{\dimexpr\textwidth-2\tabcolsep-2\arrayrulewidth}|}{\rule{0pt}{2.2em}} \\
\hline
\multicolumn{3}{|>{\raggedright\arraybackslash}p{\dimexpr\textwidth-2\tabcolsep-2\arrayrulewidth}|}{\centering\textbf{Retroalimentación del usuario}} \\
\hline
\multicolumn{3}{|>{\raggedright\arraybackslash}p{\dimexpr\textwidth-2\tabcolsep-2\arrayrulewidth}|}{\rule{0pt}{2.8em}} \\
\hline
\multicolumn{3}{|>{\raggedright\arraybackslash}p{\dimexpr\textwidth-2\tabcolsep-2\arrayrulewidth}|}{\centering\textbf{Resultado de la prueba}} \\
\hline
\multicolumn{3}{|>{\raggedright\arraybackslash}p{\dimexpr\textwidth-2\tabcolsep-2\arrayrulewidth}|}{\casilla{} Aceptado \qquad \casilla{} Aceptado parcialmente \qquad \casilla{} No aceptado} \\
\hline
\multicolumn{3}{|>{\raggedright\arraybackslash}p{\dimexpr\textwidth-2\tabcolsep-2\arrayrulewidth}|}{\centering\textbf{Observaciones}} \\
\hline
\multicolumn{3}{|>{\raggedright\arraybackslash}p{\dimexpr\textwidth-2\tabcolsep-2\arrayrulewidth}|}{\rule{0pt}{2.2em}} \\
\hline
\end{tabularx}
\caption{Prueba de aceptación del requisito FR-PAC-03 (R15)}
\label{tab:uat_fr-pac-03}
\end{table}
\begin{table}[H]
\centering
\small
\begin{tabularx}{\textwidth}{|>{\raggedright\arraybackslash}p{2.7cm}|X|>{\raggedright\arraybackslash}p{2.6cm}|}
\hline
\textbf{ID del requisito} & \textbf{Descripción breve} & \textbf{Rol del usuario} \\
\hline
FR-PAC-04 (R16) & Intervalo temporal o \textit{trials} del análisis & Investigador \\
\hline
\multicolumn{3}{|>{\raggedright\arraybackslash}p{\dimexpr\textwidth-2\tabcolsep-2\arrayrulewidth}|}{\centering\textbf{Descripción}} \\
\hline
\multicolumn{3}{|>{\raggedright\arraybackslash}p{\dimexpr\textwidth-2\tabcolsep-2\arrayrulewidth}|}{El sistema debe permitir definir el intervalo temporal (o los \textit{trials}) sobre el cual se ejecutará el análisis PAC.} \\
\hline
\multicolumn{3}{|>{\raggedright\arraybackslash}p{\dimexpr\textwidth-2\tabcolsep-2\arrayrulewidth}|}{\centering\textbf{Criterios de aceptación}} \\
\hline
\multicolumn{3}{|>{\raggedright\arraybackslash}p{\dimexpr\textwidth-2\tabcolsep-2\arrayrulewidth}|}{Al ejecutar la funcionalidad con señales reales del laboratorio, el usuario comprueba que el sistema permite definir el intervalo temporal (o los \textit{trials}) sobre el cual se ejecutará el análisis PAC, sin errores y con un resultado coherente con lo que muestra la interfaz.} \\
\hline
\multicolumn{3}{|>{\raggedright\arraybackslash}p{\dimexpr\textwidth-2\tabcolsep-2\arrayrulewidth}|}{\centering\textbf{Antes (GammaLab 1.0)}} \\
\hline
\multicolumn{3}{|>{\raggedright\arraybackslash}p{\dimexpr\textwidth-2\tabcolsep-2\arrayrulewidth}|}{Funcionalidad nueva: la versión 1.0 no incluía el análisis de acoplamiento fase-amplitud (PAC).} \\
\hline
\multicolumn{3}{|>{\raggedright\arraybackslash}p{\dimexpr\textwidth-2\tabcolsep-2\arrayrulewidth}|}{\centering\textbf{Después (GammaLab 2.0, observado durante la prueba)}} \\
\hline
\multicolumn{3}{|>{\raggedright\arraybackslash}p{\dimexpr\textwidth-2\tabcolsep-2\arrayrulewidth}|}{\rule{0pt}{2.2em}} \\
\hline
\multicolumn{3}{|>{\raggedright\arraybackslash}p{\dimexpr\textwidth-2\tabcolsep-2\arrayrulewidth}|}{\centering\textbf{Retroalimentación del usuario}} \\
\hline
\multicolumn{3}{|>{\raggedright\arraybackslash}p{\dimexpr\textwidth-2\tabcolsep-2\arrayrulewidth}|}{\rule{0pt}{2.8em}} \\
\hline
\multicolumn{3}{|>{\raggedright\arraybackslash}p{\dimexpr\textwidth-2\tabcolsep-2\arrayrulewidth}|}{\centering\textbf{Resultado de la prueba}} \\
\hline
\multicolumn{3}{|>{\raggedright\arraybackslash}p{\dimexpr\textwidth-2\tabcolsep-2\arrayrulewidth}|}{\casilla{} Aceptado \qquad \casilla{} Aceptado parcialmente \qquad \casilla{} No aceptado} \\
\hline
\multicolumn{3}{|>{\raggedright\arraybackslash}p{\dimexpr\textwidth-2\tabcolsep-2\arrayrulewidth}|}{\centering\textbf{Observaciones}} \\
\hline
\multicolumn{3}{|>{\raggedright\arraybackslash}p{\dimexpr\textwidth-2\tabcolsep-2\arrayrulewidth}|}{\rule{0pt}{2.2em}} \\
\hline
\end{tabularx}
\caption{Prueba de aceptación del requisito FR-PAC-04 (R16)}
\label{tab:uat_fr-pac-04}
\end{table}
\begin{table}[H]
\centering
\small
\begin{tabularx}{\textwidth}{|>{\raggedright\arraybackslash}p{2.7cm}|X|>{\raggedright\arraybackslash}p{2.6cm}|}
\hline
\textbf{ID del requisito} & \textbf{Descripción breve} & \textbf{Rol del usuario} \\
\hline
FR-PAC-05 (R17) & Calcular el acoplamiento por fase y frecuencia & Investigador \\
\hline
\multicolumn{3}{|>{\raggedright\arraybackslash}p{\dimexpr\textwidth-2\tabcolsep-2\arrayrulewidth}|}{\centering\textbf{Descripción}} \\
\hline
\multicolumn{3}{|>{\raggedright\arraybackslash}p{\dimexpr\textwidth-2\tabcolsep-2\arrayrulewidth}|}{El sistema debe calcular, para la banda de fase seleccionada y el rango de amplitud seleccionado, un mapa de acoplamiento fase-frecuencia mediante transformada} \\
\hline
\multicolumn{3}{|>{\raggedright\arraybackslash}p{\dimexpr\textwidth-2\tabcolsep-2\arrayrulewidth}|}{\centering\textbf{Criterios de aceptación}} \\
\hline
\multicolumn{3}{|>{\raggedright\arraybackslash}p{\dimexpr\textwidth-2\tabcolsep-2\arrayrulewidth}|}{Al ejecutar la funcionalidad con señales reales del laboratorio, el usuario comprueba que el sistema calcula, para la banda de fase seleccionada y el rango de amplitud seleccionado, un mapa de acoplamiento fase-frecuencia mediante transformada, sin errores y con un resultado coherente con lo que muestra la interfaz.} \\
\hline
\multicolumn{3}{|>{\raggedright\arraybackslash}p{\dimexpr\textwidth-2\tabcolsep-2\arrayrulewidth}|}{\centering\textbf{Antes (GammaLab 1.0)}} \\
\hline
\multicolumn{3}{|>{\raggedright\arraybackslash}p{\dimexpr\textwidth-2\tabcolsep-2\arrayrulewidth}|}{Funcionalidad nueva: la versión 1.0 no incluía el análisis de acoplamiento fase-amplitud (PAC).} \\
\hline
\multicolumn{3}{|>{\raggedright\arraybackslash}p{\dimexpr\textwidth-2\tabcolsep-2\arrayrulewidth}|}{\centering\textbf{Después (GammaLab 2.0, observado durante la prueba)}} \\
\hline
\multicolumn{3}{|>{\raggedright\arraybackslash}p{\dimexpr\textwidth-2\tabcolsep-2\arrayrulewidth}|}{\rule{0pt}{2.2em}} \\
\hline
\multicolumn{3}{|>{\raggedright\arraybackslash}p{\dimexpr\textwidth-2\tabcolsep-2\arrayrulewidth}|}{\centering\textbf{Retroalimentación del usuario}} \\
\hline
\multicolumn{3}{|>{\raggedright\arraybackslash}p{\dimexpr\textwidth-2\tabcolsep-2\arrayrulewidth}|}{\rule{0pt}{2.8em}} \\
\hline
\multicolumn{3}{|>{\raggedright\arraybackslash}p{\dimexpr\textwidth-2\tabcolsep-2\arrayrulewidth}|}{\centering\textbf{Resultado de la prueba}} \\
\hline
\multicolumn{3}{|>{\raggedright\arraybackslash}p{\dimexpr\textwidth-2\tabcolsep-2\arrayrulewidth}|}{\casilla{} Aceptado \qquad \casilla{} Aceptado parcialmente \qquad \casilla{} No aceptado} \\
\hline
\multicolumn{3}{|>{\raggedright\arraybackslash}p{\dimexpr\textwidth-2\tabcolsep-2\arrayrulewidth}|}{\centering\textbf{Observaciones}} \\
\hline
\multicolumn{3}{|>{\raggedright\arraybackslash}p{\dimexpr\textwidth-2\tabcolsep-2\arrayrulewidth}|}{\rule{0pt}{2.2em}} \\
\hline
\end{tabularx}
\caption{Prueba de aceptación del requisito FR-PAC-05 (R17)}
\label{tab:uat_fr-pac-05}
\end{table}
\begin{table}[H]
\centering
\small
\begin{tabularx}{\textwidth}{|>{\raggedright\arraybackslash}p{2.7cm}|X|>{\raggedright\arraybackslash}p{2.6cm}|}
\hline
\textbf{ID del requisito} & \textbf{Descripción breve} & \textbf{Rol del usuario} \\
\hline
FR-PAC-06 (R18) & Mapa de acoplamiento fase-frecuencia & Investigador \\
\hline
\multicolumn{3}{|>{\raggedright\arraybackslash}p{\dimexpr\textwidth-2\tabcolsep-2\arrayrulewidth}|}{\centering\textbf{Descripción}} \\
\hline
\multicolumn{3}{|>{\raggedright\arraybackslash}p{\dimexpr\textwidth-2\tabcolsep-2\arrayrulewidth}|}{El sistema debe representar los resultados del análisis PAC mediante un mapa fase-frecuencia (fase en el eje X, frecuencia de amplitud en el eje Y).} \\
\hline
\multicolumn{3}{|>{\raggedright\arraybackslash}p{\dimexpr\textwidth-2\tabcolsep-2\arrayrulewidth}|}{\centering\textbf{Criterios de aceptación}} \\
\hline
\multicolumn{3}{|>{\raggedright\arraybackslash}p{\dimexpr\textwidth-2\tabcolsep-2\arrayrulewidth}|}{Al ejecutar la funcionalidad con señales reales del laboratorio, el usuario comprueba que el sistema representa los resultados del análisis PAC mediante un mapa fase-frecuencia (fase en el eje X, frecuencia de amplitud en el eje Y), sin errores y con un resultado coherente con lo que muestra la interfaz.} \\
\hline
\multicolumn{3}{|>{\raggedright\arraybackslash}p{\dimexpr\textwidth-2\tabcolsep-2\arrayrulewidth}|}{\centering\textbf{Antes (GammaLab 1.0)}} \\
\hline
\multicolumn{3}{|>{\raggedright\arraybackslash}p{\dimexpr\textwidth-2\tabcolsep-2\arrayrulewidth}|}{Funcionalidad nueva: la versión 1.0 no incluía el análisis de acoplamiento fase-amplitud (PAC).} \\
\hline
\multicolumn{3}{|>{\raggedright\arraybackslash}p{\dimexpr\textwidth-2\tabcolsep-2\arrayrulewidth}|}{\centering\textbf{Después (GammaLab 2.0, observado durante la prueba)}} \\
\hline
\multicolumn{3}{|>{\raggedright\arraybackslash}p{\dimexpr\textwidth-2\tabcolsep-2\arrayrulewidth}|}{\rule{0pt}{2.2em}} \\
\hline
\multicolumn{3}{|>{\raggedright\arraybackslash}p{\dimexpr\textwidth-2\tabcolsep-2\arrayrulewidth}|}{\centering\textbf{Retroalimentación del usuario}} \\
\hline
\multicolumn{3}{|>{\raggedright\arraybackslash}p{\dimexpr\textwidth-2\tabcolsep-2\arrayrulewidth}|}{\rule{0pt}{2.8em}} \\
\hline
\multicolumn{3}{|>{\raggedright\arraybackslash}p{\dimexpr\textwidth-2\tabcolsep-2\arrayrulewidth}|}{\centering\textbf{Resultado de la prueba}} \\
\hline
\multicolumn{3}{|>{\raggedright\arraybackslash}p{\dimexpr\textwidth-2\tabcolsep-2\arrayrulewidth}|}{\casilla{} Aceptado \qquad \casilla{} Aceptado parcialmente \qquad \casilla{} No aceptado} \\
\hline
\multicolumn{3}{|>{\raggedright\arraybackslash}p{\dimexpr\textwidth-2\tabcolsep-2\arrayrulewidth}|}{\centering\textbf{Observaciones}} \\
\hline
\multicolumn{3}{|>{\raggedright\arraybackslash}p{\dimexpr\textwidth-2\tabcolsep-2\arrayrulewidth}|}{\rule{0pt}{2.2em}} \\
\hline
\end{tabularx}
\caption{Prueba de aceptación del requisito FR-PAC-06 (R18)}
\label{tab:uat_fr-pac-06}
\end{table}
\begin{table}[H]
\centering
\small
\begin{tabularx}{\textwidth}{|>{\raggedright\arraybackslash}p{2.7cm}|X|>{\raggedright\arraybackslash}p{2.6cm}|}
\hline
\textbf{ID del requisito} & \textbf{Descripción breve} & \textbf{Rol del usuario} \\
\hline
FR-PAC-07 (R19) & Distribución de amplitud según la fase & Investigador \\
\hline
\multicolumn{3}{|>{\raggedright\arraybackslash}p{\dimexpr\textwidth-2\tabcolsep-2\arrayrulewidth}|}{\centering\textbf{Descripción}} \\
\hline
\multicolumn{3}{|>{\raggedright\arraybackslash}p{\dimexpr\textwidth-2\tabcolsep-2\arrayrulewidth}|}{El sistema debe representar la distribución de amplitud en función de la fase (histograma de fase) obtenida durante el análisis PAC.} \\
\hline
\multicolumn{3}{|>{\raggedright\arraybackslash}p{\dimexpr\textwidth-2\tabcolsep-2\arrayrulewidth}|}{\centering\textbf{Criterios de aceptación}} \\
\hline
\multicolumn{3}{|>{\raggedright\arraybackslash}p{\dimexpr\textwidth-2\tabcolsep-2\arrayrulewidth}|}{Al ejecutar la funcionalidad con señales reales del laboratorio, el usuario comprueba que el sistema representa la distribución de amplitud en función de la fase (histograma de fase) obtenida durante el análisis PAC, sin errores y con un resultado coherente con lo que muestra la interfaz.} \\
\hline
\multicolumn{3}{|>{\raggedright\arraybackslash}p{\dimexpr\textwidth-2\tabcolsep-2\arrayrulewidth}|}{\centering\textbf{Antes (GammaLab 1.0)}} \\
\hline
\multicolumn{3}{|>{\raggedright\arraybackslash}p{\dimexpr\textwidth-2\tabcolsep-2\arrayrulewidth}|}{Funcionalidad nueva: la versión 1.0 no incluía el análisis de acoplamiento fase-amplitud (PAC).} \\
\hline
\multicolumn{3}{|>{\raggedright\arraybackslash}p{\dimexpr\textwidth-2\tabcolsep-2\arrayrulewidth}|}{\centering\textbf{Después (GammaLab 2.0, observado durante la prueba)}} \\
\hline
\multicolumn{3}{|>{\raggedright\arraybackslash}p{\dimexpr\textwidth-2\tabcolsep-2\arrayrulewidth}|}{\rule{0pt}{2.2em}} \\
\hline
\multicolumn{3}{|>{\raggedright\arraybackslash}p{\dimexpr\textwidth-2\tabcolsep-2\arrayrulewidth}|}{\centering\textbf{Retroalimentación del usuario}} \\
\hline
\multicolumn{3}{|>{\raggedright\arraybackslash}p{\dimexpr\textwidth-2\tabcolsep-2\arrayrulewidth}|}{\rule{0pt}{2.8em}} \\
\hline
\multicolumn{3}{|>{\raggedright\arraybackslash}p{\dimexpr\textwidth-2\tabcolsep-2\arrayrulewidth}|}{\centering\textbf{Resultado de la prueba}} \\
\hline
\multicolumn{3}{|>{\raggedright\arraybackslash}p{\dimexpr\textwidth-2\tabcolsep-2\arrayrulewidth}|}{\casilla{} Aceptado \qquad \casilla{} Aceptado parcialmente \qquad \casilla{} No aceptado} \\
\hline
\multicolumn{3}{|>{\raggedright\arraybackslash}p{\dimexpr\textwidth-2\tabcolsep-2\arrayrulewidth}|}{\centering\textbf{Observaciones}} \\
\hline
\multicolumn{3}{|>{\raggedright\arraybackslash}p{\dimexpr\textwidth-2\tabcolsep-2\arrayrulewidth}|}{\rule{0pt}{2.2em}} \\
\hline
\end{tabularx}
\caption{Prueba de aceptación del requisito FR-PAC-07 (R19)}
\label{tab:uat_fr-pac-07}
\end{table}
\begin{table}[H]
\centering
\small
\begin{tabularx}{\textwidth}{|>{\raggedright\arraybackslash}p{2.7cm}|X|>{\raggedright\arraybackslash}p{2.6cm}|}
\hline
\textbf{ID del requisito} & \textbf{Descripción breve} & \textbf{Rol del usuario} \\
\hline
FR-PAC-08 (R20) & Vector de acoplamiento medio & Investigador \\
\hline
\multicolumn{3}{|>{\raggedright\arraybackslash}p{\dimexpr\textwidth-2\tabcolsep-2\arrayrulewidth}|}{\centering\textbf{Descripción}} \\
\hline
\multicolumn{3}{|>{\raggedright\arraybackslash}p{\dimexpr\textwidth-2\tabcolsep-2\arrayrulewidth}|}{El sistema debe calcular y mostrar el valor del vector de acoplamiento medio (magnitud y ángulo) derivado del histograma de fase.} \\
\hline
\multicolumn{3}{|>{\raggedright\arraybackslash}p{\dimexpr\textwidth-2\tabcolsep-2\arrayrulewidth}|}{\centering\textbf{Criterios de aceptación}} \\
\hline
\multicolumn{3}{|>{\raggedright\arraybackslash}p{\dimexpr\textwidth-2\tabcolsep-2\arrayrulewidth}|}{Al ejecutar la funcionalidad con señales reales del laboratorio, el usuario comprueba que el sistema calcula y mostrar el valor del vector de acoplamiento medio (magnitud y ángulo) derivado del histograma de fase, sin errores y con un resultado coherente con lo que muestra la interfaz.} \\
\hline
\multicolumn{3}{|>{\raggedright\arraybackslash}p{\dimexpr\textwidth-2\tabcolsep-2\arrayrulewidth}|}{\centering\textbf{Antes (GammaLab 1.0)}} \\
\hline
\multicolumn{3}{|>{\raggedright\arraybackslash}p{\dimexpr\textwidth-2\tabcolsep-2\arrayrulewidth}|}{Funcionalidad nueva: la versión 1.0 no incluía el análisis de acoplamiento fase-amplitud (PAC).} \\
\hline
\multicolumn{3}{|>{\raggedright\arraybackslash}p{\dimexpr\textwidth-2\tabcolsep-2\arrayrulewidth}|}{\centering\textbf{Después (GammaLab 2.0, observado durante la prueba)}} \\
\hline
\multicolumn{3}{|>{\raggedright\arraybackslash}p{\dimexpr\textwidth-2\tabcolsep-2\arrayrulewidth}|}{\rule{0pt}{2.2em}} \\
\hline
\multicolumn{3}{|>{\raggedright\arraybackslash}p{\dimexpr\textwidth-2\tabcolsep-2\arrayrulewidth}|}{\centering\textbf{Retroalimentación del usuario}} \\
\hline
\multicolumn{3}{|>{\raggedright\arraybackslash}p{\dimexpr\textwidth-2\tabcolsep-2\arrayrulewidth}|}{\rule{0pt}{2.8em}} \\
\hline
\multicolumn{3}{|>{\raggedright\arraybackslash}p{\dimexpr\textwidth-2\tabcolsep-2\arrayrulewidth}|}{\centering\textbf{Resultado de la prueba}} \\
\hline
\multicolumn{3}{|>{\raggedright\arraybackslash}p{\dimexpr\textwidth-2\tabcolsep-2\arrayrulewidth}|}{\casilla{} Aceptado \qquad \casilla{} Aceptado parcialmente \qquad \casilla{} No aceptado} \\
\hline
\multicolumn{3}{|>{\raggedright\arraybackslash}p{\dimexpr\textwidth-2\tabcolsep-2\arrayrulewidth}|}{\centering\textbf{Observaciones}} \\
\hline
\multicolumn{3}{|>{\raggedright\arraybackslash}p{\dimexpr\textwidth-2\tabcolsep-2\arrayrulewidth}|}{\rule{0pt}{2.2em}} \\
\hline
\end{tabularx}
\caption{Prueba de aceptación del requisito FR-PAC-08 (R20)}
\label{tab:uat_fr-pac-08}
\end{table}
\begin{table}[H]
\centering
\small
\begin{tabularx}{\textwidth}{|>{\raggedright\arraybackslash}p{2.7cm}|X|>{\raggedright\arraybackslash}p{2.6cm}|}
\hline
\textbf{ID del requisito} & \textbf{Descripción breve} & \textbf{Rol del usuario} \\
\hline
FR-PAC-09 (R21) & Exportar los números de PAC & Investigador \\
\hline
\multicolumn{3}{|>{\raggedright\arraybackslash}p{\dimexpr\textwidth-2\tabcolsep-2\arrayrulewidth}|}{\centering\textbf{Descripción}} \\
\hline
\multicolumn{3}{|>{\raggedright\arraybackslash}p{\dimexpr\textwidth-2\tabcolsep-2\arrayrulewidth}|}{El sistema debe exportar los resultados numéricos del análisis PAC de un canal.} \\
\hline
\multicolumn{3}{|>{\raggedright\arraybackslash}p{\dimexpr\textwidth-2\tabcolsep-2\arrayrulewidth}|}{\centering\textbf{Criterios de aceptación}} \\
\hline
\multicolumn{3}{|>{\raggedright\arraybackslash}p{\dimexpr\textwidth-2\tabcolsep-2\arrayrulewidth}|}{Al ejecutar la funcionalidad con señales reales del laboratorio, el usuario comprueba que el sistema exporta los resultados numéricos del análisis PAC de un canal, sin errores y con un resultado coherente con lo que muestra la interfaz.} \\
\hline
\multicolumn{3}{|>{\raggedright\arraybackslash}p{\dimexpr\textwidth-2\tabcolsep-2\arrayrulewidth}|}{\centering\textbf{Antes (GammaLab 1.0)}} \\
\hline
\multicolumn{3}{|>{\raggedright\arraybackslash}p{\dimexpr\textwidth-2\tabcolsep-2\arrayrulewidth}|}{Funcionalidad nueva: la versión 1.0 no incluía el análisis de acoplamiento fase-amplitud (PAC).} \\
\hline
\multicolumn{3}{|>{\raggedright\arraybackslash}p{\dimexpr\textwidth-2\tabcolsep-2\arrayrulewidth}|}{\centering\textbf{Después (GammaLab 2.0, observado durante la prueba)}} \\
\hline
\multicolumn{3}{|>{\raggedright\arraybackslash}p{\dimexpr\textwidth-2\tabcolsep-2\arrayrulewidth}|}{\rule{0pt}{2.2em}} \\
\hline
\multicolumn{3}{|>{\raggedright\arraybackslash}p{\dimexpr\textwidth-2\tabcolsep-2\arrayrulewidth}|}{\centering\textbf{Retroalimentación del usuario}} \\
\hline
\multicolumn{3}{|>{\raggedright\arraybackslash}p{\dimexpr\textwidth-2\tabcolsep-2\arrayrulewidth}|}{\rule{0pt}{2.8em}} \\
\hline
\multicolumn{3}{|>{\raggedright\arraybackslash}p{\dimexpr\textwidth-2\tabcolsep-2\arrayrulewidth}|}{\centering\textbf{Resultado de la prueba}} \\
\hline
\multicolumn{3}{|>{\raggedright\arraybackslash}p{\dimexpr\textwidth-2\tabcolsep-2\arrayrulewidth}|}{\casilla{} Aceptado \qquad \casilla{} Aceptado parcialmente \qquad \casilla{} No aceptado} \\
\hline
\multicolumn{3}{|>{\raggedright\arraybackslash}p{\dimexpr\textwidth-2\tabcolsep-2\arrayrulewidth}|}{\centering\textbf{Observaciones}} \\
\hline
\multicolumn{3}{|>{\raggedright\arraybackslash}p{\dimexpr\textwidth-2\tabcolsep-2\arrayrulewidth}|}{\rule{0pt}{2.2em}} \\
\hline
\end{tabularx}
\caption{Prueba de aceptación del requisito FR-PAC-09 (R21)}
\label{tab:uat_fr-pac-09}
\end{table}
\begin{table}[H]
\centering
\small
\begin{tabularx}{\textwidth}{|>{\raggedright\arraybackslash}p{2.7cm}|X|>{\raggedright\arraybackslash}p{2.6cm}|}
\hline
\textbf{ID del requisito} & \textbf{Descripción breve} & \textbf{Rol del usuario} \\
\hline
FR-PAC-10 (R22) & Exportar las gráficas de PAC & Investigador \\
\hline
\multicolumn{3}{|>{\raggedright\arraybackslash}p{\dimexpr\textwidth-2\tabcolsep-2\arrayrulewidth}|}{\centering\textbf{Descripción}} \\
\hline
\multicolumn{3}{|>{\raggedright\arraybackslash}p{\dimexpr\textwidth-2\tabcolsep-2\arrayrulewidth}|}{El sistema debe exportar las representaciones gráficas generadas durante el análisis PAC de un canal.} \\
\hline
\multicolumn{3}{|>{\raggedright\arraybackslash}p{\dimexpr\textwidth-2\tabcolsep-2\arrayrulewidth}|}{\centering\textbf{Criterios de aceptación}} \\
\hline
\multicolumn{3}{|>{\raggedright\arraybackslash}p{\dimexpr\textwidth-2\tabcolsep-2\arrayrulewidth}|}{Al ejecutar la funcionalidad con señales reales del laboratorio, el usuario comprueba que el sistema exporta las representaciones gráficas generadas durante el análisis PAC de un canal, sin errores y con un resultado coherente con lo que muestra la interfaz.} \\
\hline
\multicolumn{3}{|>{\raggedright\arraybackslash}p{\dimexpr\textwidth-2\tabcolsep-2\arrayrulewidth}|}{\centering\textbf{Antes (GammaLab 1.0)}} \\
\hline
\multicolumn{3}{|>{\raggedright\arraybackslash}p{\dimexpr\textwidth-2\tabcolsep-2\arrayrulewidth}|}{Funcionalidad nueva: la versión 1.0 no incluía el análisis de acoplamiento fase-amplitud (PAC).} \\
\hline
\multicolumn{3}{|>{\raggedright\arraybackslash}p{\dimexpr\textwidth-2\tabcolsep-2\arrayrulewidth}|}{\centering\textbf{Después (GammaLab 2.0, observado durante la prueba)}} \\
\hline
\multicolumn{3}{|>{\raggedright\arraybackslash}p{\dimexpr\textwidth-2\tabcolsep-2\arrayrulewidth}|}{\rule{0pt}{2.2em}} \\
\hline
\multicolumn{3}{|>{\raggedright\arraybackslash}p{\dimexpr\textwidth-2\tabcolsep-2\arrayrulewidth}|}{\centering\textbf{Retroalimentación del usuario}} \\
\hline
\multicolumn{3}{|>{\raggedright\arraybackslash}p{\dimexpr\textwidth-2\tabcolsep-2\arrayrulewidth}|}{\rule{0pt}{2.8em}} \\
\hline
\multicolumn{3}{|>{\raggedright\arraybackslash}p{\dimexpr\textwidth-2\tabcolsep-2\arrayrulewidth}|}{\centering\textbf{Resultado de la prueba}} \\
\hline
\multicolumn{3}{|>{\raggedright\arraybackslash}p{\dimexpr\textwidth-2\tabcolsep-2\arrayrulewidth}|}{\casilla{} Aceptado \qquad \casilla{} Aceptado parcialmente \qquad \casilla{} No aceptado} \\
\hline
\multicolumn{3}{|>{\raggedright\arraybackslash}p{\dimexpr\textwidth-2\tabcolsep-2\arrayrulewidth}|}{\centering\textbf{Observaciones}} \\
\hline
\multicolumn{3}{|>{\raggedright\arraybackslash}p{\dimexpr\textwidth-2\tabcolsep-2\arrayrulewidth}|}{\rule{0pt}{2.2em}} \\
\hline
\end{tabularx}
\caption{Prueba de aceptación del requisito FR-PAC-10 (R22)}
\label{tab:uat_fr-pac-10}
\end{table}
\begin{table}[H]
\centering
\small
\begin{tabularx}{\textwidth}{|>{\raggedright\arraybackslash}p{2.7cm}|X|>{\raggedright\arraybackslash}p{2.6cm}|}
\hline
\textbf{ID del requisito} & \textbf{Descripción breve} & \textbf{Rol del usuario} \\
\hline
FR-PAC-11 (R23) & Guardar los parámetros de cada ejecución & Investigador \\
\hline
\multicolumn{3}{|>{\raggedright\arraybackslash}p{\dimexpr\textwidth-2\tabcolsep-2\arrayrulewidth}|}{\centering\textbf{Descripción}} \\
\hline
\multicolumn{3}{|>{\raggedright\arraybackslash}p{\dimexpr\textwidth-2\tabcolsep-2\arrayrulewidth}|}{El sistema debe almacenar los parámetros utilizados en cada ejecución del análisis PAC junto con sus resultados.} \\
\hline
\multicolumn{3}{|>{\raggedright\arraybackslash}p{\dimexpr\textwidth-2\tabcolsep-2\arrayrulewidth}|}{\centering\textbf{Criterios de aceptación}} \\
\hline
\multicolumn{3}{|>{\raggedright\arraybackslash}p{\dimexpr\textwidth-2\tabcolsep-2\arrayrulewidth}|}{Al ejecutar la funcionalidad con señales reales del laboratorio, el usuario comprueba que el sistema almacena los parámetros utilizados en cada ejecución del análisis PAC junto con sus resultados, sin errores y con un resultado coherente con lo que muestra la interfaz.} \\
\hline
\multicolumn{3}{|>{\raggedright\arraybackslash}p{\dimexpr\textwidth-2\tabcolsep-2\arrayrulewidth}|}{\centering\textbf{Antes (GammaLab 1.0)}} \\
\hline
\multicolumn{3}{|>{\raggedright\arraybackslash}p{\dimexpr\textwidth-2\tabcolsep-2\arrayrulewidth}|}{Funcionalidad nueva: la versión 1.0 no incluía el análisis de acoplamiento fase-amplitud (PAC).} \\
\hline
\multicolumn{3}{|>{\raggedright\arraybackslash}p{\dimexpr\textwidth-2\tabcolsep-2\arrayrulewidth}|}{\centering\textbf{Después (GammaLab 2.0, observado durante la prueba)}} \\
\hline
\multicolumn{3}{|>{\raggedright\arraybackslash}p{\dimexpr\textwidth-2\tabcolsep-2\arrayrulewidth}|}{\rule{0pt}{2.2em}} \\
\hline
\multicolumn{3}{|>{\raggedright\arraybackslash}p{\dimexpr\textwidth-2\tabcolsep-2\arrayrulewidth}|}{\centering\textbf{Retroalimentación del usuario}} \\
\hline
\multicolumn{3}{|>{\raggedright\arraybackslash}p{\dimexpr\textwidth-2\tabcolsep-2\arrayrulewidth}|}{\rule{0pt}{2.8em}} \\
\hline
\multicolumn{3}{|>{\raggedright\arraybackslash}p{\dimexpr\textwidth-2\tabcolsep-2\arrayrulewidth}|}{\centering\textbf{Resultado de la prueba}} \\
\hline
\multicolumn{3}{|>{\raggedright\arraybackslash}p{\dimexpr\textwidth-2\tabcolsep-2\arrayrulewidth}|}{\casilla{} Aceptado \qquad \casilla{} Aceptado parcialmente \qquad \casilla{} No aceptado} \\
\hline
\multicolumn{3}{|>{\raggedright\arraybackslash}p{\dimexpr\textwidth-2\tabcolsep-2\arrayrulewidth}|}{\centering\textbf{Observaciones}} \\
\hline
\multicolumn{3}{|>{\raggedright\arraybackslash}p{\dimexpr\textwidth-2\tabcolsep-2\arrayrulewidth}|}{\rule{0pt}{2.2em}} \\
\hline
\end{tabularx}
\caption{Prueba de aceptación del requisito FR-PAC-11 (R23)}
\label{tab:uat_fr-pac-11}
\end{table}
\begin{table}[H]
\centering
\small
\begin{tabularx}{\textwidth}{|>{\raggedright\arraybackslash}p{2.7cm}|X|>{\raggedright\arraybackslash}p{2.6cm}|}
\hline
\textbf{ID del requisito} & \textbf{Descripción breve} & \textbf{Rol del usuario} \\
\hline
FR-PAC-12 (R60) & Elegir el canal de fase y el de amplitud & Investigador \\
\hline
\multicolumn{3}{|>{\raggedright\arraybackslash}p{\dimexpr\textwidth-2\tabcolsep-2\arrayrulewidth}|}{\centering\textbf{Descripción}} \\
\hline
\multicolumn{3}{|>{\raggedright\arraybackslash}p{\dimexpr\textwidth-2\tabcolsep-2\arrayrulewidth}|}{El sistema debe permitir seleccionar dos canales de la misma señal: uno como fuente de fase y otro como fuente de amplitud.} \\
\hline
\multicolumn{3}{|>{\raggedright\arraybackslash}p{\dimexpr\textwidth-2\tabcolsep-2\arrayrulewidth}|}{\centering\textbf{Criterios de aceptación}} \\
\hline
\multicolumn{3}{|>{\raggedright\arraybackslash}p{\dimexpr\textwidth-2\tabcolsep-2\arrayrulewidth}|}{Al ejecutar la funcionalidad con señales reales del laboratorio, el usuario comprueba que el sistema permite seleccionar dos canales de la misma señal: uno como fuente de fase y otro como fuente de amplitud, sin errores y con un resultado coherente con lo que muestra la interfaz.} \\
\hline
\multicolumn{3}{|>{\raggedright\arraybackslash}p{\dimexpr\textwidth-2\tabcolsep-2\arrayrulewidth}|}{\centering\textbf{Antes (GammaLab 1.0)}} \\
\hline
\multicolumn{3}{|>{\raggedright\arraybackslash}p{\dimexpr\textwidth-2\tabcolsep-2\arrayrulewidth}|}{Funcionalidad nueva: la versión 1.0 no incluía el análisis de acoplamiento fase-amplitud (PAC).} \\
\hline
\multicolumn{3}{|>{\raggedright\arraybackslash}p{\dimexpr\textwidth-2\tabcolsep-2\arrayrulewidth}|}{\centering\textbf{Después (GammaLab 2.0, observado durante la prueba)}} \\
\hline
\multicolumn{3}{|>{\raggedright\arraybackslash}p{\dimexpr\textwidth-2\tabcolsep-2\arrayrulewidth}|}{\rule{0pt}{2.2em}} \\
\hline
\multicolumn{3}{|>{\raggedright\arraybackslash}p{\dimexpr\textwidth-2\tabcolsep-2\arrayrulewidth}|}{\centering\textbf{Retroalimentación del usuario}} \\
\hline
\multicolumn{3}{|>{\raggedright\arraybackslash}p{\dimexpr\textwidth-2\tabcolsep-2\arrayrulewidth}|}{\rule{0pt}{2.8em}} \\
\hline
\multicolumn{3}{|>{\raggedright\arraybackslash}p{\dimexpr\textwidth-2\tabcolsep-2\arrayrulewidth}|}{\centering\textbf{Resultado de la prueba}} \\
\hline
\multicolumn{3}{|>{\raggedright\arraybackslash}p{\dimexpr\textwidth-2\tabcolsep-2\arrayrulewidth}|}{\casilla{} Aceptado \qquad \casilla{} Aceptado parcialmente \qquad \casilla{} No aceptado} \\
\hline
\multicolumn{3}{|>{\raggedright\arraybackslash}p{\dimexpr\textwidth-2\tabcolsep-2\arrayrulewidth}|}{\centering\textbf{Observaciones}} \\
\hline
\multicolumn{3}{|>{\raggedright\arraybackslash}p{\dimexpr\textwidth-2\tabcolsep-2\arrayrulewidth}|}{\rule{0pt}{2.2em}} \\
\hline
\end{tabularx}
\caption{Prueba de aceptación del requisito FR-PAC-12 (R60)}
\label{tab:uat_fr-pac-12}
\end{table}
\begin{table}[H]
\centering
\small
\begin{tabularx}{\textwidth}{|>{\raggedright\arraybackslash}p{2.7cm}|X|>{\raggedright\arraybackslash}p{2.6cm}|}
\hline
\textbf{ID del requisito} & \textbf{Descripción breve} & \textbf{Rol del usuario} \\
\hline
FR-PAC-13 (R61) & Mapa de acoplamiento entre canales & Investigador \\
\hline
\multicolumn{3}{|>{\raggedright\arraybackslash}p{\dimexpr\textwidth-2\tabcolsep-2\arrayrulewidth}|}{\centering\textbf{Descripción}} \\
\hline
\multicolumn{3}{|>{\raggedright\arraybackslash}p{\dimexpr\textwidth-2\tabcolsep-2\arrayrulewidth}|}{El sistema debe calcular el mapa de acoplamiento fase-frecuencia entre el canal de fase y el canal de amplitud seleccionados.} \\
\hline
\multicolumn{3}{|>{\raggedright\arraybackslash}p{\dimexpr\textwidth-2\tabcolsep-2\arrayrulewidth}|}{\centering\textbf{Criterios de aceptación}} \\
\hline
\multicolumn{3}{|>{\raggedright\arraybackslash}p{\dimexpr\textwidth-2\tabcolsep-2\arrayrulewidth}|}{Al ejecutar la funcionalidad con señales reales del laboratorio, el usuario comprueba que el sistema calcula el mapa de acoplamiento fase-frecuencia entre el canal de fase y el canal de amplitud seleccionados, sin errores y con un resultado coherente con lo que muestra la interfaz.} \\
\hline
\multicolumn{3}{|>{\raggedright\arraybackslash}p{\dimexpr\textwidth-2\tabcolsep-2\arrayrulewidth}|}{\centering\textbf{Antes (GammaLab 1.0)}} \\
\hline
\multicolumn{3}{|>{\raggedright\arraybackslash}p{\dimexpr\textwidth-2\tabcolsep-2\arrayrulewidth}|}{Funcionalidad nueva: la versión 1.0 no incluía el análisis de acoplamiento fase-amplitud (PAC).} \\
\hline
\multicolumn{3}{|>{\raggedright\arraybackslash}p{\dimexpr\textwidth-2\tabcolsep-2\arrayrulewidth}|}{\centering\textbf{Después (GammaLab 2.0, observado durante la prueba)}} \\
\hline
\multicolumn{3}{|>{\raggedright\arraybackslash}p{\dimexpr\textwidth-2\tabcolsep-2\arrayrulewidth}|}{\rule{0pt}{2.2em}} \\
\hline
\multicolumn{3}{|>{\raggedright\arraybackslash}p{\dimexpr\textwidth-2\tabcolsep-2\arrayrulewidth}|}{\centering\textbf{Retroalimentación del usuario}} \\
\hline
\multicolumn{3}{|>{\raggedright\arraybackslash}p{\dimexpr\textwidth-2\tabcolsep-2\arrayrulewidth}|}{\rule{0pt}{2.8em}} \\
\hline
\multicolumn{3}{|>{\raggedright\arraybackslash}p{\dimexpr\textwidth-2\tabcolsep-2\arrayrulewidth}|}{\centering\textbf{Resultado de la prueba}} \\
\hline
\multicolumn{3}{|>{\raggedright\arraybackslash}p{\dimexpr\textwidth-2\tabcolsep-2\arrayrulewidth}|}{\casilla{} Aceptado \qquad \casilla{} Aceptado parcialmente \qquad \casilla{} No aceptado} \\
\hline
\multicolumn{3}{|>{\raggedright\arraybackslash}p{\dimexpr\textwidth-2\tabcolsep-2\arrayrulewidth}|}{\centering\textbf{Observaciones}} \\
\hline
\multicolumn{3}{|>{\raggedright\arraybackslash}p{\dimexpr\textwidth-2\tabcolsep-2\arrayrulewidth}|}{\rule{0pt}{2.2em}} \\
\hline
\end{tabularx}
\caption{Prueba de aceptación del requisito FR-PAC-13 (R61)}
\label{tab:uat_fr-pac-13}
\end{table}
\begin{table}[H]
\centering
\small
\begin{tabularx}{\textwidth}{|>{\raggedright\arraybackslash}p{2.7cm}|X|>{\raggedright\arraybackslash}p{2.6cm}|}
\hline
\textbf{ID del requisito} & \textbf{Descripción breve} & \textbf{Rol del usuario} \\
\hline
FR-PAC-14 (R62) & Mismas visualizaciones que PAC de un canal & Investigador \\
\hline
\multicolumn{3}{|>{\raggedright\arraybackslash}p{\dimexpr\textwidth-2\tabcolsep-2\arrayrulewidth}|}{\centering\textbf{Descripción}} \\
\hline
\multicolumn{3}{|>{\raggedright\arraybackslash}p{\dimexpr\textwidth-2\tabcolsep-2\arrayrulewidth}|}{El sistema debe representar los resultados del análisis PAC inter-canal con las mismas visualizaciones que el PAC de un canal.} \\
\hline
\multicolumn{3}{|>{\raggedright\arraybackslash}p{\dimexpr\textwidth-2\tabcolsep-2\arrayrulewidth}|}{\centering\textbf{Criterios de aceptación}} \\
\hline
\multicolumn{3}{|>{\raggedright\arraybackslash}p{\dimexpr\textwidth-2\tabcolsep-2\arrayrulewidth}|}{Al ejecutar la funcionalidad con señales reales del laboratorio, el usuario comprueba que el sistema representa los resultados del análisis PAC inter-canal con las mismas visualizaciones que el PAC de un canal, sin errores y con un resultado coherente con lo que muestra la interfaz.} \\
\hline
\multicolumn{3}{|>{\raggedright\arraybackslash}p{\dimexpr\textwidth-2\tabcolsep-2\arrayrulewidth}|}{\centering\textbf{Antes (GammaLab 1.0)}} \\
\hline
\multicolumn{3}{|>{\raggedright\arraybackslash}p{\dimexpr\textwidth-2\tabcolsep-2\arrayrulewidth}|}{Funcionalidad nueva: la versión 1.0 no incluía el análisis de acoplamiento fase-amplitud (PAC).} \\
\hline
\multicolumn{3}{|>{\raggedright\arraybackslash}p{\dimexpr\textwidth-2\tabcolsep-2\arrayrulewidth}|}{\centering\textbf{Después (GammaLab 2.0, observado durante la prueba)}} \\
\hline
\multicolumn{3}{|>{\raggedright\arraybackslash}p{\dimexpr\textwidth-2\tabcolsep-2\arrayrulewidth}|}{\rule{0pt}{2.2em}} \\
\hline
\multicolumn{3}{|>{\raggedright\arraybackslash}p{\dimexpr\textwidth-2\tabcolsep-2\arrayrulewidth}|}{\centering\textbf{Retroalimentación del usuario}} \\
\hline
\multicolumn{3}{|>{\raggedright\arraybackslash}p{\dimexpr\textwidth-2\tabcolsep-2\arrayrulewidth}|}{\rule{0pt}{2.8em}} \\
\hline
\multicolumn{3}{|>{\raggedright\arraybackslash}p{\dimexpr\textwidth-2\tabcolsep-2\arrayrulewidth}|}{\centering\textbf{Resultado de la prueba}} \\
\hline
\multicolumn{3}{|>{\raggedright\arraybackslash}p{\dimexpr\textwidth-2\tabcolsep-2\arrayrulewidth}|}{\casilla{} Aceptado \qquad \casilla{} Aceptado parcialmente \qquad \casilla{} No aceptado} \\
\hline
\multicolumn{3}{|>{\raggedright\arraybackslash}p{\dimexpr\textwidth-2\tabcolsep-2\arrayrulewidth}|}{\centering\textbf{Observaciones}} \\
\hline
\multicolumn{3}{|>{\raggedright\arraybackslash}p{\dimexpr\textwidth-2\tabcolsep-2\arrayrulewidth}|}{\rule{0pt}{2.2em}} \\
\hline
\end{tabularx}
\caption{Prueba de aceptación del requisito FR-PAC-14 (R62)}
\label{tab:uat_fr-pac-14}
\end{table}
\begin{table}[H]
\centering
\small
\begin{tabularx}{\textwidth}{|>{\raggedright\arraybackslash}p{2.7cm}|X|>{\raggedright\arraybackslash}p{2.6cm}|}
\hline
\textbf{ID del requisito} & \textbf{Descripción breve} & \textbf{Rol del usuario} \\
\hline
FR-PAC-15 (R63) & PPC: dos canales y rango de frecuencias & Investigador \\
\hline
\multicolumn{3}{|>{\raggedright\arraybackslash}p{\dimexpr\textwidth-2\tabcolsep-2\arrayrulewidth}|}{\centering\textbf{Descripción}} \\
\hline
\multicolumn{3}{|>{\raggedright\arraybackslash}p{\dimexpr\textwidth-2\tabcolsep-2\arrayrulewidth}|}{El sistema debe permitir seleccionar dos canales y un rango de frecuencias para calcular la coherencia de fase entre ellos.} \\
\hline
\multicolumn{3}{|>{\raggedright\arraybackslash}p{\dimexpr\textwidth-2\tabcolsep-2\arrayrulewidth}|}{\centering\textbf{Criterios de aceptación}} \\
\hline
\multicolumn{3}{|>{\raggedright\arraybackslash}p{\dimexpr\textwidth-2\tabcolsep-2\arrayrulewidth}|}{Al ejecutar la funcionalidad con señales reales del laboratorio, el usuario comprueba que el sistema permite seleccionar dos canales y un rango de frecuencias para calcular la coherencia de fase entre ellos, sin errores y con un resultado coherente con lo que muestra la interfaz.} \\
\hline
\multicolumn{3}{|>{\raggedright\arraybackslash}p{\dimexpr\textwidth-2\tabcolsep-2\arrayrulewidth}|}{\centering\textbf{Antes (GammaLab 1.0)}} \\
\hline
\multicolumn{3}{|>{\raggedright\arraybackslash}p{\dimexpr\textwidth-2\tabcolsep-2\arrayrulewidth}|}{Funcionalidad nueva: la versión 1.0 no incluía el análisis de consistencia de fase entre canales (PPC).} \\
\hline
\multicolumn{3}{|>{\raggedright\arraybackslash}p{\dimexpr\textwidth-2\tabcolsep-2\arrayrulewidth}|}{\centering\textbf{Después (GammaLab 2.0, observado durante la prueba)}} \\
\hline
\multicolumn{3}{|>{\raggedright\arraybackslash}p{\dimexpr\textwidth-2\tabcolsep-2\arrayrulewidth}|}{\rule{0pt}{2.2em}} \\
\hline
\multicolumn{3}{|>{\raggedright\arraybackslash}p{\dimexpr\textwidth-2\tabcolsep-2\arrayrulewidth}|}{\centering\textbf{Retroalimentación del usuario}} \\
\hline
\multicolumn{3}{|>{\raggedright\arraybackslash}p{\dimexpr\textwidth-2\tabcolsep-2\arrayrulewidth}|}{\rule{0pt}{2.8em}} \\
\hline
\multicolumn{3}{|>{\raggedright\arraybackslash}p{\dimexpr\textwidth-2\tabcolsep-2\arrayrulewidth}|}{\centering\textbf{Resultado de la prueba}} \\
\hline
\multicolumn{3}{|>{\raggedright\arraybackslash}p{\dimexpr\textwidth-2\tabcolsep-2\arrayrulewidth}|}{\casilla{} Aceptado \qquad \casilla{} Aceptado parcialmente \qquad \casilla{} No aceptado} \\
\hline
\multicolumn{3}{|>{\raggedright\arraybackslash}p{\dimexpr\textwidth-2\tabcolsep-2\arrayrulewidth}|}{\centering\textbf{Observaciones}} \\
\hline
\multicolumn{3}{|>{\raggedright\arraybackslash}p{\dimexpr\textwidth-2\tabcolsep-2\arrayrulewidth}|}{\rule{0pt}{2.2em}} \\
\hline
\end{tabularx}
\caption{Prueba de aceptación del requisito FR-PAC-15 (R63)}
\label{tab:uat_fr-pac-15}
\end{table}
\begin{table}[H]
\centering
\small
\begin{tabularx}{\textwidth}{|>{\raggedright\arraybackslash}p{2.7cm}|X|>{\raggedright\arraybackslash}p{2.6cm}|}
\hline
\textbf{ID del requisito} & \textbf{Descripción breve} & \textbf{Rol del usuario} \\
\hline
FR-PAC-16 (R64) & PPC: mapa de frecuencia y tiempo & Investigador \\
\hline
\multicolumn{3}{|>{\raggedright\arraybackslash}p{\dimexpr\textwidth-2\tabcolsep-2\arrayrulewidth}|}{\centering\textbf{Descripción}} \\
\hline
\multicolumn{3}{|>{\raggedright\arraybackslash}p{\dimexpr\textwidth-2\tabcolsep-2\arrayrulewidth}|}{El sistema debe representar la coherencia de fase mediante un mapa tiempo-frecuencia.} \\
\hline
\multicolumn{3}{|>{\raggedright\arraybackslash}p{\dimexpr\textwidth-2\tabcolsep-2\arrayrulewidth}|}{\centering\textbf{Criterios de aceptación}} \\
\hline
\multicolumn{3}{|>{\raggedright\arraybackslash}p{\dimexpr\textwidth-2\tabcolsep-2\arrayrulewidth}|}{Al ejecutar la funcionalidad con señales reales del laboratorio, el usuario comprueba que el sistema representa la coherencia de fase mediante un mapa tiempo-frecuencia, sin errores y con un resultado coherente con lo que muestra la interfaz.} \\
\hline
\multicolumn{3}{|>{\raggedright\arraybackslash}p{\dimexpr\textwidth-2\tabcolsep-2\arrayrulewidth}|}{\centering\textbf{Antes (GammaLab 1.0)}} \\
\hline
\multicolumn{3}{|>{\raggedright\arraybackslash}p{\dimexpr\textwidth-2\tabcolsep-2\arrayrulewidth}|}{Funcionalidad nueva: la versión 1.0 no incluía el análisis de consistencia de fase entre canales (PPC).} \\
\hline
\multicolumn{3}{|>{\raggedright\arraybackslash}p{\dimexpr\textwidth-2\tabcolsep-2\arrayrulewidth}|}{\centering\textbf{Después (GammaLab 2.0, observado durante la prueba)}} \\
\hline
\multicolumn{3}{|>{\raggedright\arraybackslash}p{\dimexpr\textwidth-2\tabcolsep-2\arrayrulewidth}|}{\rule{0pt}{2.2em}} \\
\hline
\multicolumn{3}{|>{\raggedright\arraybackslash}p{\dimexpr\textwidth-2\tabcolsep-2\arrayrulewidth}|}{\centering\textbf{Retroalimentación del usuario}} \\
\hline
\multicolumn{3}{|>{\raggedright\arraybackslash}p{\dimexpr\textwidth-2\tabcolsep-2\arrayrulewidth}|}{\rule{0pt}{2.8em}} \\
\hline
\multicolumn{3}{|>{\raggedright\arraybackslash}p{\dimexpr\textwidth-2\tabcolsep-2\arrayrulewidth}|}{\centering\textbf{Resultado de la prueba}} \\
\hline
\multicolumn{3}{|>{\raggedright\arraybackslash}p{\dimexpr\textwidth-2\tabcolsep-2\arrayrulewidth}|}{\casilla{} Aceptado \qquad \casilla{} Aceptado parcialmente \qquad \casilla{} No aceptado} \\
\hline
\multicolumn{3}{|>{\raggedright\arraybackslash}p{\dimexpr\textwidth-2\tabcolsep-2\arrayrulewidth}|}{\centering\textbf{Observaciones}} \\
\hline
\multicolumn{3}{|>{\raggedright\arraybackslash}p{\dimexpr\textwidth-2\tabcolsep-2\arrayrulewidth}|}{\rule{0pt}{2.2em}} \\
\hline
\end{tabularx}
\caption{Prueba de aceptación del requisito FR-PAC-16 (R64)}
\label{tab:uat_fr-pac-16}
\end{table}
\begin{table}[H]
\centering
\small
\begin{tabularx}{\textwidth}{|>{\raggedright\arraybackslash}p{2.7cm}|X|>{\raggedright\arraybackslash}p{2.6cm}|}
\hline
\textbf{ID del requisito} & \textbf{Descripción breve} & \textbf{Rol del usuario} \\
\hline
FR-PAC-17 (R65) & Exportar los números de PPC y del PAC entre canales & Investigador \\
\hline
\multicolumn{3}{|>{\raggedright\arraybackslash}p{\dimexpr\textwidth-2\tabcolsep-2\arrayrulewidth}|}{\centering\textbf{Descripción}} \\
\hline
\multicolumn{3}{|>{\raggedright\arraybackslash}p{\dimexpr\textwidth-2\tabcolsep-2\arrayrulewidth}|}{El sistema debe exportar los resultados numéricos del análisis PAC inter-canal y de coherencia de fase (PPC).} \\
\hline
\multicolumn{3}{|>{\raggedright\arraybackslash}p{\dimexpr\textwidth-2\tabcolsep-2\arrayrulewidth}|}{\centering\textbf{Criterios de aceptación}} \\
\hline
\multicolumn{3}{|>{\raggedright\arraybackslash}p{\dimexpr\textwidth-2\tabcolsep-2\arrayrulewidth}|}{Al ejecutar la funcionalidad con señales reales del laboratorio, el usuario comprueba que el sistema exporta los resultados numéricos del análisis PAC inter-canal y de coherencia de fase (PPC), sin errores y con un resultado coherente con lo que muestra la interfaz.} \\
\hline
\multicolumn{3}{|>{\raggedright\arraybackslash}p{\dimexpr\textwidth-2\tabcolsep-2\arrayrulewidth}|}{\centering\textbf{Antes (GammaLab 1.0)}} \\
\hline
\multicolumn{3}{|>{\raggedright\arraybackslash}p{\dimexpr\textwidth-2\tabcolsep-2\arrayrulewidth}|}{Funcionalidad nueva: la versión 1.0 no incluía el análisis de consistencia de fase entre canales (PPC).} \\
\hline
\multicolumn{3}{|>{\raggedright\arraybackslash}p{\dimexpr\textwidth-2\tabcolsep-2\arrayrulewidth}|}{\centering\textbf{Después (GammaLab 2.0, observado durante la prueba)}} \\
\hline
\multicolumn{3}{|>{\raggedright\arraybackslash}p{\dimexpr\textwidth-2\tabcolsep-2\arrayrulewidth}|}{\rule{0pt}{2.2em}} \\
\hline
\multicolumn{3}{|>{\raggedright\arraybackslash}p{\dimexpr\textwidth-2\tabcolsep-2\arrayrulewidth}|}{\centering\textbf{Retroalimentación del usuario}} \\
\hline
\multicolumn{3}{|>{\raggedright\arraybackslash}p{\dimexpr\textwidth-2\tabcolsep-2\arrayrulewidth}|}{\rule{0pt}{2.8em}} \\
\hline
\multicolumn{3}{|>{\raggedright\arraybackslash}p{\dimexpr\textwidth-2\tabcolsep-2\arrayrulewidth}|}{\centering\textbf{Resultado de la prueba}} \\
\hline
\multicolumn{3}{|>{\raggedright\arraybackslash}p{\dimexpr\textwidth-2\tabcolsep-2\arrayrulewidth}|}{\casilla{} Aceptado \qquad \casilla{} Aceptado parcialmente \qquad \casilla{} No aceptado} \\
\hline
\multicolumn{3}{|>{\raggedright\arraybackslash}p{\dimexpr\textwidth-2\tabcolsep-2\arrayrulewidth}|}{\centering\textbf{Observaciones}} \\
\hline
\multicolumn{3}{|>{\raggedright\arraybackslash}p{\dimexpr\textwidth-2\tabcolsep-2\arrayrulewidth}|}{\rule{0pt}{2.2em}} \\
\hline
\end{tabularx}
\caption{Prueba de aceptación del requisito FR-PAC-17 (R65)}
\label{tab:uat_fr-pac-17}
\end{table}
\begin{table}[H]
\centering
\small
\begin{tabularx}{\textwidth}{|>{\raggedright\arraybackslash}p{2.7cm}|X|>{\raggedright\arraybackslash}p{2.6cm}|}
\hline
\textbf{ID del requisito} & \textbf{Descripción breve} & \textbf{Rol del usuario} \\
\hline
FR-PAC-18 (R66) & Exportar las gráficas de PPC y del PAC entre canales & Investigador \\
\hline
\multicolumn{3}{|>{\raggedright\arraybackslash}p{\dimexpr\textwidth-2\tabcolsep-2\arrayrulewidth}|}{\centering\textbf{Descripción}} \\
\hline
\multicolumn{3}{|>{\raggedright\arraybackslash}p{\dimexpr\textwidth-2\tabcolsep-2\arrayrulewidth}|}{El sistema debe exportar las representaciones gráficas del análisis PAC inter-canal y de coherencia de fase (PPC) como imagen.} \\
\hline
\multicolumn{3}{|>{\raggedright\arraybackslash}p{\dimexpr\textwidth-2\tabcolsep-2\arrayrulewidth}|}{\centering\textbf{Criterios de aceptación}} \\
\hline
\multicolumn{3}{|>{\raggedright\arraybackslash}p{\dimexpr\textwidth-2\tabcolsep-2\arrayrulewidth}|}{Al ejecutar la funcionalidad con señales reales del laboratorio, el usuario comprueba que el sistema exporta las representaciones gráficas del análisis PAC inter-canal y de coherencia de fase (PPC) como imagen, sin errores y con un resultado coherente con lo que muestra la interfaz.} \\
\hline
\multicolumn{3}{|>{\raggedright\arraybackslash}p{\dimexpr\textwidth-2\tabcolsep-2\arrayrulewidth}|}{\centering\textbf{Antes (GammaLab 1.0)}} \\
\hline
\multicolumn{3}{|>{\raggedright\arraybackslash}p{\dimexpr\textwidth-2\tabcolsep-2\arrayrulewidth}|}{Funcionalidad nueva: la versión 1.0 no incluía el análisis de consistencia de fase entre canales (PPC).} \\
\hline
\multicolumn{3}{|>{\raggedright\arraybackslash}p{\dimexpr\textwidth-2\tabcolsep-2\arrayrulewidth}|}{\centering\textbf{Después (GammaLab 2.0, observado durante la prueba)}} \\
\hline
\multicolumn{3}{|>{\raggedright\arraybackslash}p{\dimexpr\textwidth-2\tabcolsep-2\arrayrulewidth}|}{\rule{0pt}{2.2em}} \\
\hline
\multicolumn{3}{|>{\raggedright\arraybackslash}p{\dimexpr\textwidth-2\tabcolsep-2\arrayrulewidth}|}{\centering\textbf{Retroalimentación del usuario}} \\
\hline
\multicolumn{3}{|>{\raggedright\arraybackslash}p{\dimexpr\textwidth-2\tabcolsep-2\arrayrulewidth}|}{\rule{0pt}{2.8em}} \\
\hline
\multicolumn{3}{|>{\raggedright\arraybackslash}p{\dimexpr\textwidth-2\tabcolsep-2\arrayrulewidth}|}{\centering\textbf{Resultado de la prueba}} \\
\hline
\multicolumn{3}{|>{\raggedright\arraybackslash}p{\dimexpr\textwidth-2\tabcolsep-2\arrayrulewidth}|}{\casilla{} Aceptado \qquad \casilla{} Aceptado parcialmente \qquad \casilla{} No aceptado} \\
\hline
\multicolumn{3}{|>{\raggedright\arraybackslash}p{\dimexpr\textwidth-2\tabcolsep-2\arrayrulewidth}|}{\centering\textbf{Observaciones}} \\
\hline
\multicolumn{3}{|>{\raggedright\arraybackslash}p{\dimexpr\textwidth-2\tabcolsep-2\arrayrulewidth}|}{\rule{0pt}{2.2em}} \\
\hline
\end{tabularx}
\caption{Prueba de aceptación del requisito FR-PAC-18 (R66)}
\label{tab:uat_fr-pac-18}
\end{table}
\begin{table}[H]
\centering
\small
\begin{tabularx}{\textwidth}{|>{\raggedright\arraybackslash}p{2.7cm}|X|>{\raggedright\arraybackslash}p{2.6cm}|}
\hline
\textbf{ID del requisito} & \textbf{Descripción breve} & \textbf{Rol del usuario} \\
\hline
FR-PAC-19 (R75) & PAC Psel segmenta por ciclo de la onda lenta & Investigador \\
\hline
\multicolumn{3}{|>{\raggedright\arraybackslash}p{\dimexpr\textwidth-2\tabcolsep-2\arrayrulewidth}|}{\centering\textbf{Descripción}} \\
\hline
\multicolumn{3}{|>{\raggedright\arraybackslash}p{\dimexpr\textwidth-2\tabcolsep-2\arrayrulewidth}|}{El sistema debe segmentar la señal filtrada en la banda de frecuencia rápida tomando cada ciclo de la onda lenta como una época, y superponer todas las épocas en una misma gráfica de tipo butterfly con la fase del ciclo lento en el eje X.} \\
\hline
\multicolumn{3}{|>{\raggedright\arraybackslash}p{\dimexpr\textwidth-2\tabcolsep-2\arrayrulewidth}|}{\centering\textbf{Criterios de aceptación}} \\
\hline
\multicolumn{3}{|>{\raggedright\arraybackslash}p{\dimexpr\textwidth-2\tabcolsep-2\arrayrulewidth}|}{Al ejecutar la funcionalidad con señales reales del laboratorio, el usuario comprueba que el sistema segmenta la señal filtrada en la banda de frecuencia rápida tomando cada ciclo de la onda lenta como una época, y superponer todas las épocas en una misma gráfica de tipo butterfly con la fase del ciclo lento en el eje X, sin errores y con un resultado coherente con lo que muestra la interfaz.} \\
\hline
\multicolumn{3}{|>{\raggedright\arraybackslash}p{\dimexpr\textwidth-2\tabcolsep-2\arrayrulewidth}|}{\centering\textbf{Antes (GammaLab 1.0)}} \\
\hline
\multicolumn{3}{|>{\raggedright\arraybackslash}p{\dimexpr\textwidth-2\tabcolsep-2\arrayrulewidth}|}{Funcionalidad nueva: la versión 1.0 no incluía el análisis de acoplamiento fase-amplitud (PAC).} \\
\hline
\multicolumn{3}{|>{\raggedright\arraybackslash}p{\dimexpr\textwidth-2\tabcolsep-2\arrayrulewidth}|}{\centering\textbf{Después (GammaLab 2.0, observado durante la prueba)}} \\
\hline
\multicolumn{3}{|>{\raggedright\arraybackslash}p{\dimexpr\textwidth-2\tabcolsep-2\arrayrulewidth}|}{\rule{0pt}{2.2em}} \\
\hline
\multicolumn{3}{|>{\raggedright\arraybackslash}p{\dimexpr\textwidth-2\tabcolsep-2\arrayrulewidth}|}{\centering\textbf{Retroalimentación del usuario}} \\
\hline
\multicolumn{3}{|>{\raggedright\arraybackslash}p{\dimexpr\textwidth-2\tabcolsep-2\arrayrulewidth}|}{\rule{0pt}{2.8em}} \\
\hline
\multicolumn{3}{|>{\raggedright\arraybackslash}p{\dimexpr\textwidth-2\tabcolsep-2\arrayrulewidth}|}{\centering\textbf{Resultado de la prueba}} \\
\hline
\multicolumn{3}{|>{\raggedright\arraybackslash}p{\dimexpr\textwidth-2\tabcolsep-2\arrayrulewidth}|}{\casilla{} Aceptado \qquad \casilla{} Aceptado parcialmente \qquad \casilla{} No aceptado} \\
\hline
\multicolumn{3}{|>{\raggedright\arraybackslash}p{\dimexpr\textwidth-2\tabcolsep-2\arrayrulewidth}|}{\centering\textbf{Observaciones}} \\
\hline
\multicolumn{3}{|>{\raggedright\arraybackslash}p{\dimexpr\textwidth-2\tabcolsep-2\arrayrulewidth}|}{\rule{0pt}{2.2em}} \\
\hline
\end{tabularx}
\caption{Prueba de aceptación del requisito FR-PAC-19 (R75)}
\label{tab:uat_fr-pac-19}
\end{table}
\begin{table}[H]
\centering
\small
\begin{tabularx}{\textwidth}{|>{\raggedright\arraybackslash}p{2.7cm}|X|>{\raggedright\arraybackslash}p{2.6cm}|}
\hline
\textbf{ID del requisito} & \textbf{Descripción breve} & \textbf{Rol del usuario} \\
\hline
FR-PAC-20 (R76) & Filtro Butterworth pasa-banda de fase cero & Investigador \\
\hline
\multicolumn{3}{|>{\raggedright\arraybackslash}p{\dimexpr\textwidth-2\tabcolsep-2\arrayrulewidth}|}{\centering\textbf{Descripción}} \\
\hline
\multicolumn{3}{|>{\raggedright\arraybackslash}p{\dimexpr\textwidth-2\tabcolsep-2\arrayrulewidth}|}{El sistema debe aplicar sobre la señal activa un filtro Butterworth pasa-banda de fase cero, con frecuencia de corte inferior, frecuencia de corte superior y orden configurables por el usuario.} \\
\hline
\multicolumn{3}{|>{\raggedright\arraybackslash}p{\dimexpr\textwidth-2\tabcolsep-2\arrayrulewidth}|}{\centering\textbf{Criterios de aceptación}} \\
\hline
\multicolumn{3}{|>{\raggedright\arraybackslash}p{\dimexpr\textwidth-2\tabcolsep-2\arrayrulewidth}|}{Al ejecutar la funcionalidad con señales reales del laboratorio, el usuario comprueba que el sistema aplica sobre la señal activa un filtro Butterworth pasa-banda de fase cero, con frecuencia de corte inferior, frecuencia de corte superior y orden configurables por el usuario, sin errores y con un resultado coherente con lo que muestra la interfaz.} \\
\hline
\multicolumn{3}{|>{\raggedright\arraybackslash}p{\dimexpr\textwidth-2\tabcolsep-2\arrayrulewidth}|}{\centering\textbf{Antes (GammaLab 1.0)}} \\
\hline
\multicolumn{3}{|>{\raggedright\arraybackslash}p{\dimexpr\textwidth-2\tabcolsep-2\arrayrulewidth}|}{La versión 1.0 aplicaba el filtro Butterworth; los filtros elíptico y Chebyshev provocaban bloqueos de la aplicación.} \\
\hline
\multicolumn{3}{|>{\raggedright\arraybackslash}p{\dimexpr\textwidth-2\tabcolsep-2\arrayrulewidth}|}{\centering\textbf{Después (GammaLab 2.0, observado durante la prueba)}} \\
\hline
\multicolumn{3}{|>{\raggedright\arraybackslash}p{\dimexpr\textwidth-2\tabcolsep-2\arrayrulewidth}|}{\rule{0pt}{2.2em}} \\
\hline
\multicolumn{3}{|>{\raggedright\arraybackslash}p{\dimexpr\textwidth-2\tabcolsep-2\arrayrulewidth}|}{\centering\textbf{Retroalimentación del usuario}} \\
\hline
\multicolumn{3}{|>{\raggedright\arraybackslash}p{\dimexpr\textwidth-2\tabcolsep-2\arrayrulewidth}|}{\rule{0pt}{2.8em}} \\
\hline
\multicolumn{3}{|>{\raggedright\arraybackslash}p{\dimexpr\textwidth-2\tabcolsep-2\arrayrulewidth}|}{\centering\textbf{Resultado de la prueba}} \\
\hline
\multicolumn{3}{|>{\raggedright\arraybackslash}p{\dimexpr\textwidth-2\tabcolsep-2\arrayrulewidth}|}{\casilla{} Aceptado \qquad \casilla{} Aceptado parcialmente \qquad \casilla{} No aceptado} \\
\hline
\multicolumn{3}{|>{\raggedright\arraybackslash}p{\dimexpr\textwidth-2\tabcolsep-2\arrayrulewidth}|}{\centering\textbf{Observaciones}} \\
\hline
\multicolumn{3}{|>{\raggedright\arraybackslash}p{\dimexpr\textwidth-2\tabcolsep-2\arrayrulewidth}|}{\rule{0pt}{2.2em}} \\
\hline
\end{tabularx}
\caption{Prueba de aceptación del requisito FR-PAC-20 (R76)}
\label{tab:uat_fr-pac-20}
\end{table}

\subsection{Consolidado de resultados}

\begin{center}
\small
\begin{tabularx}{\textwidth}{|X|c|c|}
\hline
\textbf{Resultado} & \textbf{Requisitos} & \textbf{Porcentaje} \\
\hline
Aceptados & & \\
\hline
Aceptados parcialmente & & \\
\hline
No aceptados & & \\
\hline
\textbf{Total} & \textbf{45} & \textbf{100\,\%} \\
\hline
\end{tabularx}
\captionof{table}{Consolidado de las pruebas de aceptación}
\label{tab:td_uat_consolidado}
\end{center}

\subsection{Conclusiones de las pruebas de aceptación}
\pendiente[Se redactarán una vez ejecutadas las pruebas con el usuario, a partir de los cuadros de aceptación y del consolidado de resultados.]
\capitulo{TRAZABILIDAD ENTRE REQUISITOS Y PRUEBAS}
La Tabla~\ref{tab:td_trazabilidad} relaciona cada uno de los 81 requisitos del SRS con las pruebas automáticas, las mediciones o las verificaciones manuales que lo comprueban. Los requisitos funcionales se validan, además, con el usuario en los cuadros de aceptación de la Sección~\ref{sec:td_uat}.

\begin{center}
\footnotesize
\begin{longtable}{|>{\raggedright\arraybackslash}p{0.12\textwidth}|>{\raggedright\arraybackslash}p{0.07\textwidth}|>{\raggedright\arraybackslash}p{0.24\textwidth}|>{\raggedright\arraybackslash}p{0.45\textwidth}|}
\hline
\textbf{Requisito} & \textbf{Matriz} & \textbf{Descripción breve} & \textbf{Prueba o verificación} \\
\hline
\endfirsthead
\hline
\textbf{Requisito} & \textbf{Matriz} & \textbf{Descripción breve} & \textbf{Prueba o verificación} \\
\hline
\endhead
\hline
\endfoot
\hline
\caption{Trazabilidad entre requisitos y pruebas} \label{tab:td_trazabilidad} \\
\endlastfoot
FR-IO-01 & R08 & Persistir la señal cargada & Pruebas automáticas del almacén de datos y de la lectura sin duplicación (\textit{test\_dataset}, \textit{test\_fileio\_sin\_duplicar}). \\
\hline
FR-IO-02 & R80 & Cargar señales ABF, EDF y MAT & Pruebas automáticas de la lectura de ABF y EDF; la de MAT se comprueba en la prueba del programa instalado (\textit{test\_fileio\_abf}, \textit{test\_fileio\_edf}, \textit{prueba\_congelado}). \\
\hline
FR-IO-03 & R07 & Guardar un proyecto & Pruebas automáticas del guardado en pasos, de su cancelación y del registro de la ruta (\textit{test\_guardado\_proyecto}, \textit{test\_correcciones\_interaccion}). \\
\hline
FR-IO-04 & R78 & Abrir un proyecto guardado & Pruebas automáticas de la apertura, de la señal que falta y de la ruta de origen, y de la persistencia de las modificaciones de artefactos (\textit{test\_apertura\_proyecto}, \textit{test\_persistencia\_artifact\_remove}, \textit{test\_correcciones\_interaccion}). \\
\hline
FR-IO-05 & R05 & Exportar los resultados de los análisis & Pruebas automáticas de la exportación y del escritor de Excel. \textit{Relative PSD} dibuja su gráfica con el menú contextual y exporta parámetros, potencias y fracción por banda (\textit{test\_export}, \textit{test\_xlsx\_writer}, \textit{test\_resultados\_xlsx}, \textit{test\_relative\_psd\_exportacion}). \\
\hline
FR-IO-06 & R09 & Una sola columna X en el CSV & Pruebas automáticas de la exportación con ejes iguales y distintos (\textit{test\_export}). \\
\hline
FR-IO-07 & R28 & Nombre de exportación descriptivo & Verificación manual del nombre propuesto al exportar. \\
\hline
FR-UI-01 & R01 & Cortar la señal en \textit{trials} & Pruebas automáticas de la detección de estímulos y del corte (\textit{test\_trials}). \\
\hline
FR-UI-02 & R24 & Modificar un \textit{trial} ya creado & Pruebas automáticas de los conjuntos derivados, que dejan el original intacto (\textit{test\_trials\_derivados}). \\
\hline
FR-UI-03 & R25 & Restablecer parámetros por defecto & Prueba automática del Clear de \textit{Trials} y verificación manual en las demás extensiones (\textit{test\_correcciones\_interaccion}). \\
\hline
FR-UI-04 & R26 & Valores X/Y bajo el cursor & Pruebas automáticas de la conversión de píxeles a datos, del formato de la lectura y del punto más cercano (\textit{test\_indicadores\_ui}, \textit{test\_correcciones\_interaccion}). \\
\hline
FR-UI-05 & R27 & Elegir los \textit{trials} de cada análisis & Pruebas automáticas de la selección y de los descartes (\textit{test\_average\_seleccion}, \textit{test\_indicadores\_ui}, \textit{test\_correcciones\_interaccion}). \\
\hline
FR-UI-06 & R29 & Ajustar la frecuencia de muestreo & Pruebas automáticas del cambio y de la restauración de la frecuencia (\textit{test\_cu018\_sampling\_rate}). \\
\hline
FR-UI-07 & R30 & \textit{Zoom} en todas las direcciones & Pruebas automáticas de los ejes del \textit{zoom} y verificación manual en cada gráfica (\textit{test\_correcciones\_interaccion}). \\
\hline
FR-UI-08 & R31 & Mostrar el umbral en \textit{mV} & Pruebas automáticas de la conversión de unidades y verificación manual del campo (\textit{test\_trials}). \\
\hline
FR-UI-09 & R32 & Mostrar \textit{Butterfly} explícitamente & Prueba automática del título dibujado (\textit{test\_correcciones\_interaccion}). \\
\hline
FR-UI-10 & R36 & Promediar todos los \textit{trials} & Pruebas automáticas del promedio (\textit{test\_average\_plugin}). \\
\hline
FR-UI-11 & R37 & Indicar el \textit{trial} actual & Pruebas automáticas del estado y del texto del indicador (\textit{test\_indicadores\_ui}). \\
\hline
FR-UI-12 & R38 & Indicar qué \textit{trials} participan & Pruebas automáticas de la descripción de los \textit{trials} y de la selección de \textit{Average} (\textit{test\_indicadores\_ui}, \textit{test\_average\_seleccion}). \\
\hline
FR-UI-13 & R40 & Remover un artefacto con medidas & Pruebas automáticas del cálculo de las medidas y del uso del servicio de tareas (\textit{test\_medidas\_artefacto}, \textit{test\_artifact\_remove\_orquestador}). \\
\hline
FR-UI-14 & R41 & Promediar los \textit{trials} seleccionados & Pruebas automáticas de la selección y del promedio (\textit{test\_average\_seleccion}). \\
\hline
FR-UI-15 & R68 & \textit{Trials} identificados por archivo y canal & Pruebas automáticas de la consulta de \textit{trials} por canal (\textit{test\_trials\_por\_canal}). \\
\hline
FR-WV-01 & R33 & Escala y rango de color del escalograma & Pruebas automáticas del rango de color (\textit{test\_escalograma\_rango\_color}). \\
\hline
FR-WV-02 & R34 & Normalizar el escalograma & Pruebas automáticas de las dos extensiones, con las fallas de comparación numérica descritas en DT-03 (\textit{test\_wavelet\_plugin}, \textit{test\_wavelet\_average\_plugin}, \textit{test\_submuestreo\_wavelet}). \\
\hline
FR-WV-03 & R35 & Exportar el escalograma en EMF & Pruebas automáticas del archivo EMF (\textit{test\_emf\_export}). \\
\hline
FR-PAC-01 & R13 & Ejecutar el análisis PAC & Pruebas automáticas del cálculo de PAC (\textit{test\_pac\_core}). \\
\hline
FR-PAC-02 & R14 & Banda de la fase & Pruebas automáticas del cálculo de PAC (\textit{test\_pac\_core}). \\
\hline
FR-PAC-03 & R15 & Banda de la amplitud & Pruebas automáticas del cálculo de PAC (\textit{test\_pac\_core}). \\
\hline
FR-PAC-04 & R16 & Intervalo temporal o \textit{trials} del análisis & Pruebas automáticas del selector de \textit{trial} y verificación manual de los modos (\textit{test\_pac\_intercanal\_plugin}). \\
\hline
FR-PAC-05 & R17 & Calcular el acoplamiento por fase y frecuencia & Pruebas automáticas contra la referencia de MATLAB (\textit{test\_pac\_core}). \\
\hline
FR-PAC-06 & R18 & Mapa de acoplamiento fase-frecuencia & Pruebas automáticas del eje de fase y verificación manual del mapa (\textit{test\_pac\_grados}). \\
\hline
FR-PAC-07 & R19 & Distribución de amplitud según la fase & Pruebas automáticas del cálculo del histograma y del eje de fase (\textit{test\_pac\_core}, \textit{test\_pac\_grados}). \\
\hline
FR-PAC-08 & R20 & Vector de acoplamiento medio & Pruebas automáticas del cálculo del vector (\textit{test\_pac\_core}). \\
\hline
FR-PAC-09 & R21 & Exportar los números de PAC & Pruebas automáticas de la exportación y de los ángulos en grados (\textit{test\_export}, \textit{test\_pac\_grados}). \\
\hline
FR-PAC-10 & R22 & Exportar las gráficas de PAC & Pruebas automáticas de la exportación de imágenes y de EMF (\textit{test\_export}, \textit{test\_emf\_export}). \\
\hline
FR-PAC-11 & R23 & Guardar los parámetros de cada ejecución & Prueba automática de la restauración de análisis al abrir un proyecto (\textit{test\_apertura\_proyecto}). \\
\hline
FR-PAC-12 & R60 & Elegir el canal de fase y el de amplitud & Pruebas automáticas del selector y de su comportamiento (\textit{test\_pac\_intercanal\_plugin}). \\
\hline
FR-PAC-13 & R61 & Mapa de acoplamiento entre canales & Pruebas automáticas del cálculo entre canales (\textit{test\_pac\_intercanal}). \\
\hline
FR-PAC-14 & R62 & Mismas visualizaciones que PAC de un canal & Pruebas automáticas de que ambos modos dibujan los mismos paneles (\textit{test\_pac\_intercanal\_plugin}). \\
\hline
FR-PAC-15 & R63 & PPC: dos canales y rango de frecuencias & Pruebas automáticas del cálculo de \textit{PPC} y de su ejecución en el servicio de tareas (\textit{test\_ppc\_core}, \textit{test\_ppc\_plugin\_orquestado}). \\
\hline
FR-PAC-16 & R64 & PPC: mapa de frecuencia y tiempo & Pruebas automáticas del cálculo y de la transformada de Gabor compartida (\textit{test\_ppc\_core}, \textit{test\_gabor\_en\_core}). \\
\hline
FR-PAC-17 & R65 & Exportar los números de PPC y del PAC entre canales & Pruebas automáticas que exportan archivos reales de PAC entre canales (\textit{test\_pac\_intercanal\_exportacion}). \\
\hline
FR-PAC-18 & R66 & Exportar las gráficas de PPC y del PAC entre canales & Pruebas automáticas de la exportación de imágenes y de EMF; no hay prueba específica del modo entre canales (\textit{test\_emf\_export}). \\
\hline
FR-PAC-19 & R75 & PAC Psel segmenta por ciclo de la onda lenta & Pruebas automáticas de la selección de ciclos y de la transformada de Gabor (\textit{test\_pac\_core}, \textit{test\_gabor\_en\_core}). \\
\hline
FR-PAC-20 & R76 & Filtro Butterworth pasa-banda de fase cero & Pruebas automáticas de los filtros y medición de su tiempo (\textit{test\_pac\_core}, \textit{test\_preprocessing\_bench}, \textit{test\_correcciones\_interaccion}). \\
\hline
NFR-U-01 & R02 & Diferenciar información, aviso y error & Pruebas automáticas de los niveles de mensaje (\textit{test\_niveles\_de\_mensaje}, \textit{test\_indicadores\_ui}). \\
\hline
NFR-U-02 & R03 & Resultado en la misma pantalla & Verificación manual en cada extensión. \\
\hline
NFR-U-03 & R04 & Retroalimentación en operaciones de más de 1 s & Pruebas automáticas del guardado y de la apertura. Medición: abrir un proyecto bloquea la interfaz unos 80 ms y guardar un archivo de 314 MB, 18 ms (\textit{test\_task\_service}, \textit{test\_guardado\_proyecto}, \textit{test\_apertura\_proyecto}). \\
\hline
NFR-U-04 & R06 & Iconos en las funcionalidades & Verificación manual de la cinta. \\
\hline
NFR-U-05 & R10 & Bienvenida a usuarios nuevos & Prueba automática de que la ventana sale solo la primera vez; verificación manual de la pantalla (\textit{test\_correcciones\_interaccion}). \\
\hline
NFR-U-06 & R12 & Comunicar el cierre del proyecto & Pruebas automáticas del cierre al cargar otra señal y verificación en la aplicación (\textit{test\_correcciones\_interaccion}). \\
\hline
NFR-U-07 & R42 & Descripción emergente de cada control & Prueba automática que recorre la interfaz de cada extensión y falla si algún control queda sin descripción (\textit{test\_ayudas\_controles}). \\
\hline
NFR-U-08 & R43 & Funciones de edición en \textit{Home} & Verificación manual de la pantalla. \\
\hline
NFR-U-09 & R44 & Pantalla de inicio durante la carga & Verificación manual del arranque. \\
\hline
NFR-P-01 & R11 & Liberar la señal al cerrar el proyecto & Pruebas automáticas y medición: tras cerrar el proyecto no queda ninguna copia viva de la señal (\textit{test\_liberar\_senal}). \\
\hline
NFR-P-02 & R45 & Liberar recursos gráficos al salir del módulo & Pruebas automáticas de la liberación. Medición: al recorrer los módulos de análisis, el proceso pasa de 3.118 MB sin liberar a 2.585 MB con la configuración actual (\textit{test\_liberar\_ventanas}). \\
\hline
NFR-P-03 & R47 & Cargar el archivo en 2 s o menos & Pruebas automáticas de la lectura sin duplicación. Medición con un archivo de 2 canales de 1,8 millones de muestras: 0,04 s de lectura y 0,2 s hasta verlo en pantalla (límite: 2 s) (\textit{test\_io\_bench}, \textit{test\_fileio\_sin\_duplicar}). \\
\hline
NFR-P-04 & R48 & \textit{Wavelet} en 10 s o menos & Medición: de 1,1 a 1,4 s para un \textit{trial} de 30.500 muestras con los parámetros por defecto (límite: 10 s) (\textit{test\_compute\_bench}). \\
\hline
NFR-P-05 & R49 & Graficar 3 canales sin degradar la interacción & Pruebas automáticas del submuestreo. Medición en el equipo de desarrollo, con tres canales de 1,8 millones de muestras: mediana de 29 a 44 FPS al hacer \textit{zoom} y de 24 a 36 FPS al desplazar (límite: 24 FPS); el peor cuadro de cada corrida varió entre 16 y 31 FPS según la carga del equipo. Reducir los puntos dibujados hasta 250 no modificó la tasa, por lo que el tiempo restante es fijo del renderizado (\textit{test\_submuestreo}, \textit{test\_render\_bench}). \\
\hline
NFR-P-06 & R50 & Cambiar de módulo en 500 ms o menos & Medición: el cambio de módulo tarda entre 10 y 310 ms con una señal de 1,7 millones de muestras (límite: 500 ms). \\
\hline
NFR-P-07 & R51 & Exportar 100\,000 filas en 2 s o menos & Pruebas automáticas del escritor. Medición con 100.000 filas: 0,39 s en Excel y 0,11 s en CSV (límite: 2 s) (\textit{test\_xlsx\_writer}). \\
\hline
NFR-P-08 & R52 & Liberar la señal anterior en menos de 1 s & Pruebas automáticas y medición: la señal anterior se libera en 243 a 353 ms; tras cuatro cargas seguidas queda una sola señal viva y la memoria sube 178 MB en vez de 631 MB (\textit{test\_liberar\_senal}). \\
\hline
NFR-P-09 & R53 & Cierre en menos de 3 s con recursos liberados & Medición: la ventana se cierra en 172 ms y el proceso termina en 0,63 s (límite: 3 s). \\
\hline
NFR-P-10 & R55 & Interfaz responsiva (0,3 s o menos) & Medición: la interfaz se congela entre 270 y 480 ms la primera vez que se abre PAC, según la corrida, y entre 110 y 200 ms en adelante (límite: 300 ms). \\
\hline
NFR-P-11 & R56 & Promedio de todos los \textit{trials} en 1 s o menos & Medición: de 0,32 a 0,41 s para 50 \textit{trials} de 30.500 muestras (límite: 1 s) (\textit{test\_compute\_bench}). \\
\hline
NFR-P-12 & R57 & Exportar una gráfica en 1 s (3 s si es vectorial) & Pruebas automáticas de la exportación. Medición: PNG en 0,26 s, TIFF en 0,06 s y EMF en 0,12 s (límites: 1 s y 3 s) (\textit{test\_export}, \textit{test\_emf\_export}). \\
\hline
NFR-P-13 & R59 & Conversión en bloque a las estructuras de VTK & Medición: 1,8 millones de filas convertidas en 0,03 s. Prueba automática del llenado vectorizado del mapa de calor (\textit{test\_vtk\_adapters\_bench}, \textit{test\_erp\_vtk\_llenado}). \\
\hline
NFR-P-14 & R71 & Cómputo pesado por el servicio de tareas & Pruebas automáticas del patrón y de las extensiones que lo usan (\textit{test\_run\_task}, \textit{test\_task\_service}, \textit{test\_artifact\_remove\_orquestador}, \textit{test\_ppc\_plugin\_orquestado}). \\
\hline
NFR-P-15 & R73 & Hasta 2000 puntos por serie y 200 curvas & Pruebas automáticas del submuestreo: como máximo 2000 puntos por serie, sin alterar los datos (\textit{test\_submuestreo}, \textit{test\_correcciones\_interaccion}). \\
\hline
NFR-P-16 & R74 & PAC por el servicio de tareas & Medición: un comodulograma de 20×20 pares tarda entre 0,2 y 0,5 s (límite: 30 s); la cancelación se comprueba en la aplicación. \\
\hline
NFR-P-17 & R81 & Listo en 8 s o menos & Medición: 3,1 s hasta la ventana lista (límite: 8 s) y prueba del programa instalado (\textit{test\_kernel\_bench}, \textit{prueba\_congelado}). \\
\hline
NFR-R-01 & R46 & Terminar los hilos al cancelar o cerrar & Pruebas automáticas de la cancelación cooperativa, del desligado a los dos segundos y del cambio de sección (\textit{test\_task\_service}, \textit{test\_cancelacion\_al\_cambiar\_seccion}). \\
\hline
NFR-R-02 & R54 & Sincronizar el estado compartido & Pruebas automáticas de la concurrencia sobre el almacén (\textit{test\_data\_store\_concurrencia}). \\
\hline
NFR-R-03 & R58 & Retirar los detectores de eventos de VTK & Pruebas automáticas del retiro, la reposición y el reemplazo del menú; verificación en la aplicación con cinco cálculos seguidos (\textit{test\_vtk\_observers}). \\
\hline
NFR-R-04 & R72 & Reutilizar la representación gráfica & Pruebas automáticas que verifican que la segunda petición devuelve la misma tabla, que corregir o restaurar la frecuencia la descarta y que la tabla nueva trae el eje corregido sin tocar la anterior (\textit{test\_cache\_tabla\_vtk}, \textit{test\_cu018\_sampling\_rate}). \\
\hline
NFR-R-05 & R79 & No modificar la señal ni los \textit{trials} originales & Pruebas automáticas de los conjuntos derivados (\textit{test\_trials\_derivados}). \\
\hline
NFR-S-01 & R39 & Ejecutable instalable que registra la extensión del proyecto & Pruebas automáticas de las rutas del programa instalado y prueba del ejecutable, además de la instalación real verificada en un equipo (\textit{test\_instalacion\_rutas}, \textit{prueba\_congelado}). \\
\hline
NFR-S-02 & R67 & Vía pública única sobre los \textit{trials} & Pruebas automáticas de los \textit{trials} originales, derivados y por canal (\textit{test\_trials\_derivados}, \textit{test\_trials\_por\_canal}). \\
\hline
NFR-S-03 & R69 & Servicios compartidos registrados en el núcleo & Prueba automática que recorre el código y falla si otro módulo crea uno de los servicios, y pruebas de cada servicio. Verificado en la aplicación con una exportación y una medición reales (\textit{test\_servicios\_compartidos}). \\
\hline
NFR-S-04 & R70 & Cómputo pesado por un servicio del núcleo & Pruebas automáticas del servicio de tareas y del patrón de uso (\textit{test\_task\_service}, \textit{test\_run\_task}). \\
\hline
NFR-S-05 & R77 & Nuevo \textit{plugin} sin modificar el núcleo & Pruebas automáticas con una extensión de prueba que se guarda, se restaura y se ordena sin tocar el núcleo (\textit{test\_plugin\_sin\_tocar\_nucleo}). \\
\hline
\end{longtable}
\end{center}
\capitulo{CONCLUSIONES}
La suite automática de GammaLab 2.0 consta de 761 pruebas, de las cuales 755 se aprobaron, 5 se mantienen como fallos esperados con su causa documentada y 1 falla por estar desactualizada frente a la incorporación de PPC. Frente a la versión 1.0, cuyo TD reportó 60 pruebas unitarias \cite{td2025}, la cobertura se amplió a todos los servicios del núcleo, a los módulos nuevos de análisis y a los requisitos no funcionales.

La validación frente a BOARD-FTD-PACC confirma la equivalencia de \textit{Average}, ERP, FFT, \textit{FFT Average} y PSD. En \textit{PSD Average} y en la transformada \textit{wavelet} las diferencias se explican por decisiones de implementación concretas, medidas y documentadas, y no por errores de cálculo. Los requisitos de rendimiento se cumplen con dos matices: el peor cuadro durante la interacción con tres canales y la primera apertura del módulo PAC.

Quedan pendientes las pruebas de usabilidad y de aceptación con el usuario del laboratorio, cuyos casos y cuadros se encuentran preparados en las Secciones~\ref{sec:td_usabilidad} y~\ref{sec:td_uat}, así como la actualización de la prueba \textit{test\_el\_orden\_de\_botones\_viene\_de\_los\_metadatos}.

\begin{thebibliography}{9}

\bibitem{td2025}
S. Asencio Rodríguez, M. F. González Valencia, J. E. Martínez Clavijo, y G. Villalobos, ``Test Document (TD): Gamma Lab,'' versión 1.1, Programa de Ingeniería de Sistemas, Facultad de Ingeniería, Pontificia Universidad Javeriana, Bogotá, D.C., Colombia, nov. 2025.

\bibitem{brooke1996}
J. Brooke, ``SUS: A quick and dirty usability scale,'' en \textit{Usability Evaluation in Industry}, P. W. Jordan, B. Thomas, B. A. Weerdmeester, e I. L. McClelland, Eds. Londres, Reino Unido: Taylor \& Francis, 1996, pp. 189--194.

\end{thebibliography}

\end{document}
